% Rozdzial 17 do skryptu "Nowoczesny C++" (C++20).
% Plik zawiera wylacznie tresc rozdzialu: bez \documentclass i \begin{document}.
% Wymaga pakietow: listings, amsmath, amssymb, booktabs, tabularx,
% array, enumitem, hyperref; makr \code, \term, \nota z dotychczasowego skryptu.
% Uzycie: % Rozdzial 17 do skryptu "Nowoczesny C++" (C++20).
% Plik zawiera wylacznie tresc rozdzialu: bez \documentclass i \begin{document}.
% Wymaga pakietow: listings, amsmath, amssymb, booktabs, tabularx,
% array, enumitem, hyperref; makr \code, \term, \nota z dotychczasowego skryptu.
% Uzycie: % Rozdzial 17 do skryptu "Nowoczesny C++" (C++20).
% Plik zawiera wylacznie tresc rozdzialu: bez \documentclass i \begin{document}.
% Wymaga pakietow: listings, amsmath, amssymb, booktabs, tabularx,
% array, enumitem, hyperref; makr \code, \term, \nota z dotychczasowego skryptu.
% Uzycie: % Rozdzial 17 do skryptu "Nowoczesny C++" (C++20).
% Plik zawiera wylacznie tresc rozdzialu: bez \documentclass i \begin{document}.
% Wymaga pakietow: listings, amsmath, amssymb, booktabs, tabularx,
% array, enumitem, hyperref; makr \code, \term, \nota z dotychczasowego skryptu.
% Uzycie: \input{rozdzial_17_kompilacja_linkowanie_abi_moduly_cpp.tex}
% Komendy powloki sa przeznaczone glownie dla GCC/Clang na Linux/ELF;
% roznice Windows/MSVC i macOS zaznaczono w tekscie.

\chapter{Kompilacja, linkowanie, ABI i moduły C++20}
\label{chap:cpp17-build-abi-modules}

\section{Dlaczego trzeba rozumieć proces budowania?}

W~poprzednich rozdziałach projektowaliśmy klasy, szablony, mechanizmy
zarządzania pamięcią i~współbieżność. Dopóki cały program mieści się
w~jednym pliku, wystarcza zwykle wywołanie \code{g++ main.cpp -o app}.
W~realnym projekcie kod dzieli się jednak na dziesiątki lub tysiące
\emph{jednostek translacji}, biblioteki zewnętrzne i~moduły. To wtedy
pojawiają się komunikaty \code{undefined reference}, \code{multiple definition},
konflikty wersji bibliotek, długi czas budowania oraz błędy ujawniające
się dopiero po wymianie biblioteki współdzielonej.

Musimy rozdzielić trzy różne zagadnienia:
\begin{description}[style=nextline]
\item[Poprawność programu w~języku C++]
Czy deklaracje, definicje, typy i~wywołania spełniają reguły C++,
w~szczególności \emph{One Definition Rule} (ODR)?
\item[Model budowania i~linkowania]
Jakie pliki są kompilowane osobno, gdzie kompilator zapisuje symbole,
w~jakiej kolejności linker przeszukuje biblioteki oraz jak ładowane są
biblioteki dynamiczne?
\item[Zgodność binarna (ABI)]
Czy dwa osobno zbudowane komponenty zgadzają się co do reprezentacji
typów, sposobu wywoływania funkcji, nazw symboli i~pozostałych reguł
interfejsu binarnego?
\end{description}

\nota{Norma ISO C++ opisuje semantykę programu, ale nie narzuca jednego
formatu pliku obiektowego, konwencji \emph{name mangling}, organizacji
bibliotek \code{.so}/\code{.dll} ani poleceń linkera. Poniżej
odróżniamy reguły języka od typowych rozwiązań narzędziowych.}

\section{Od pliku źródłowego do programu}

\begin{center}\small
\begin{tabularx}{\textwidth}{@{}p{0.31\textwidth}X@{}}
\toprule
Etap & Za co odpowiada \\
\midrule
Preprocesor & Rozwija dyrektywy \code{#include}, makra i warunkową kompilację. \\
Kompilator & Sprawdza składnię i typy oraz wytwarza kod dla jednostki translacji. \\
Asembler & Zamienia kod asemblera na plik obiektowy z symbolami i relokacjami. \\
Linker & Łączy pliki obiektowe i biblioteki, rozwiązując odwołania do symboli. \\
Loader systemowy & Ładuje program i wymagane biblioteki współdzielone przed uruchomieniem. \\
\bottomrule
\end{tabularx}
\end{center}
\label{sec:cpp17-pipeline}

\subsection{Preprocesor, kompilator, asembler, linker i~loader}

W~uproszczeniu klasyczny projekt z~nagłówkami przechodzi przez etapy:
\[
\begin{aligned}
\text{źródła i nagłówki}&\ \longrightarrow\ \text{preprocessing}
\ \longrightarrow\ \text{kompilacja}\\
&\phantom{\ \longrightarrow\ }\longrightarrow\ \text{plik obiektowy}
\ \longrightarrow\ \text{linkowanie}
\ \longrightarrow\ \text{program}.
\end{aligned}
\]

Na niektórych platformach kompilator wytwarza tekst asemblerowy
(\code{.s}) jako etap pośredni; w~innych przejście do kodu maszynowego
jest zintegrowane. Linkowanie może odbywać się przy budowaniu programu,
a~dodatkowe rozwiązywanie symboli może następować w~czasie ładowania
lub później.

Prześledźmy etapy na samodzielnym programie.

\begin{lstlisting}[caption={Plik hello.cpp: prosty program do obserwowania kompilacji.}]
#include <iostream>

int twice(int x) { return 2 * x; }

int main() {
    std::cout << twice(21) << '\n';
}
\end{lstlisting}

\begin{lstlisting}[language=bash,caption={Wybrane etapy budowania przy uzyciu GCC/Clang.}]
g++ -std=c++20 -E hello.cpp -o hello.ii  # Preprocesor
# Produkt poprzedniego etapu moze byc duzy.
g++ -std=c++20 -S hello.cpp -o hello.s   # Asembler

g++ -std=c++20 -c hello.cpp -o hello.o   # Kod obiektowy
g++ hello.o -o hello                      # Linkowanie
./hello                                   # 42
\end{lstlisting}

Opcja \code{-c} nie tworzy wykonywalnego programu; produkuje plik
obiektowy, który może zawierać nierozwiązane odwołania do symboli
pochodzących z~innych jednostek translacji. Polecenie \code{-S}
kończy budowanie na asemblerze, a~\code{-E} na wstępnym przetwarzaniu.
Warto odróżnić program sterujący kompilacją (np.~\code{g++}) od samego
frontendu C++, asemblera i~linkera.

\subsection{Pliki obiektowe, symbole i~relokacje}

Plik obiektowy zawiera zwykle kod maszynowy, dane, tablice symboli
oraz \emph{relokacje}, czyli informacje, które adresy muszą zostać
skorygowane po ustaleniu ostatecznego układu programu. W~środowisku
Linux powszechny jest format ELF; Windows używa PE/COFF, a~macOS
Mach-O. Szczegóły tych formatów nie są częścią języka C++.

\begin{lstlisting}[language=bash,caption={Ogladanie symboli i relokacji (Linux/ELF).}]
nm -C hello.o                 # Symbole; -C demangluje nazwy C++
readelf -Ws hello.o          # Tablica symboli ELF
readelf -r hello.o           # Relokacje
objdump -drC hello.o         # Asembler i relokacje
\end{lstlisting}

Przykładowo litera \code{T} w~typowym wyjściu \code{nm} informuje
o~symbolu definiowanym w~sekcji kodu, a~\code{U} o~symbolu
niezdefiniowanym w~danym obiekcie, który trzeba rozwiązać. Nie
należy traktować oznaczeń \code{nm} jako przenośnej własności C++.

\subsection{Jednostka translacji i~kompilacja oddzielna}

W~tradycyjnym modelu \emph{jednostka translacji} (\emph{translation unit},
TU) powstaje z~pliku źródłowego po uwzględnieniu dyrektyw preprocesora,
w~szczególności zawartości dołączanych nagłówków. Każda TU jest
analizowana niezależnie, a~linker łączy otrzymane obiekty.

Rozważmy projekt:
\begin{lstlisting}[caption={Trzy pliki klasycznego projektu C++: API, implementacja i klient.}]
// Plik: math.hpp
#ifndef CPP17_MATH_HPP
#define CPP17_MATH_HPP
int square(int x);
#endif

// Plik: math.cpp
#include "math.hpp"
int square(int x) { return x * x; }

// Plik: main.cpp
#include "math.hpp"
#include <iostream>
int main() { std::cout << square(7) << '\n'; }
\end{lstlisting}

\begin{lstlisting}[language=bash,caption={Niezalezne budowanie TU i koncowe linkowanie.}]
g++ -std=c++20 -Wall -Wextra -c math.cpp -o math.o
g++ -std=c++20 -Wall -Wextra -c main.cpp -o main.o
g++ main.o math.o -o app
./app                        # 49
\end{lstlisting}

\code{math.hpp} dostarcza deklarację klientowi i~implementacji,
ale nie jest zwykle oddzielnym obiektem linkowanym. Zmiana samego
\code{math.cpp} nie musi wymuszać ponownej kompilacji \code{main.cpp},
o~ile nie zmienił się kontrakt w~nagłówku lub pozostałe zależności.

\subsection{Przykładowe komendy w~MSVC}

W~\emph{Developer Command Prompt for Visual Studio} można zbudować
analogiczny przykład za pomocą \code{cl.exe} i~\code{link.exe}:

\begin{lstlisting}[language=bash,caption={Kompilacja dwóch plików w środowisku MSVC/Windows.}]
cl /std:c++20 /EHsc /c math.cpp
cl /std:c++20 /EHsc /c main.cpp
link main.obj math.obj /OUT:app.exe
app.exe
\end{lstlisting}

Przy budowaniu programów C++ w~Visual Studio warto również
kontrolować zgodność wersji MSVC, ustawienia biblioteki uruchomieniowej
(np.~\code{/MD} wobec \code{/MT}) i~konfigurację Debug/Release.
Takie flagi nie mają bezpośrednich odpowiedników narzuconych przez
standard języka; należą do konkretnego środowiska.

\section{Nagłówki, makra i~zasada jednej definicji}
\label{sec:cpp17-headers-odr}

\subsection{Include guards i~\texttt{\#pragma once}}

Zawartość \code{\#include} jest w~uproszczeniu wstawiana do źródła
przez preprocesor. Wielokrotne dołączenie tego samego nagłówka
w~jednej TU mogłoby więc prowadzić do powtórzenia definicji klasy.
Chronimy się przez \emph{include guard}:

\begin{lstlisting}[caption={Klasyczna os?ona nag?owka.}]
#ifndef LIB_WIDGET_HPP
#define LIB_WIDGET_HPP

struct Widget {
    int value{};
};

#endif // LIB_WIDGET_HPP
\end{lstlisting}

W~praktyce popularna jest też dyrektywa \code{\#pragma once},
obsługiwana przez najważniejsze kompilatory, lecz niestandaryzowana
przez ISO C++. Oba rozwiązania chronią \emph{jedną TU} przed
powtórnym dołączeniem pliku. Nie rozwiązują automatycznie problemu
\emph{wielu definicji tego samego nie-inline symbolu w~różnych TU}.

\subsection{ODR: deklaracja kontra definicja}

Deklaracja wprowadza nazwę i~informację o~typie, natomiast definicja
zwykle dostarcza byt danej funkcji lub obiektu. Deklaracji pewnej
funkcji może być wiele, ale \emph{nie-inline} funkcja z~zewnętrzną
łącznością, używana w~programie w~sposób wymagający definicji,
powinna mieć jedną definicję w~programie. Reguły ODR zawierają
wyjątki pozwalające na identyczne definicje w~różnych TU, np. dla
odpowiednich funkcji \code{inline}, klas i~szablonów.

\begin{lstlisting}[caption={Typowa pu?apka: definicja zwyk?ej funkcji w nag?owku.}]
// BAD: wrong.hpp, dolaczony do dwoch plikow .cpp
int twice(int x) { return 2 * x; }
// Dwie zewnetrzne definicje -> naruszenie ODR.
\end{lstlisting}

\begin{lstlisting}[caption={Definicja inline przeznaczona do nag?owka.}]
// Plik: twice.hpp
#pragma once
inline int twice(int x) { return 2 * x; }
\end{lstlisting}

\code{inline} nie nakazuje kompilatorowi wstawienia instrukcji funkcji
bezpośrednio w~miejscu wywołania. Jego znaczenie dla organizacji
programu polega przede wszystkim na dopuszczeniu definicji w~wielu TU
zgodnie z~ODR (z~odpowiednimi warunkami). Wszystkie takie definicje
muszą być zgodne i~spełniać reguły ODR; samo dopisanie \code{inline}
nie naprawia różniących się implementacji w~dwóch plikach.

\subsection{Zmienne \texttt{extern}, \texttt{inline}, \texttt{const} i~\texttt{static}}

\begin{lstlisting}[caption={Jedna definicja zmiennej globalnej, wiele deklaracji.}]
// Plik: settings.hpp
#pragma once
extern int request_limit;       // Tylko deklaracja

// Plik: settings.cpp
#include "settings.hpp"
int request_limit = 100;        // Definicja

// Plik: main.cpp
#include "settings.hpp"
int main() { return request_limit == 100 ? 0 : 1; }
\end{lstlisting}

Od C++17 w~nagłówkach można również definiować \emph{zmienne inline}:
\begin{lstlisting}[caption={Zmienna inline w nag?owku (C++17).}]
#pragma once
inline constexpr int default_timeout_ms = 250;
inline int global_counter = 0;
\end{lstlisting}

W~tym przypadku odpowiednie definicje we wszystkich TU mogą
wskazywać na ten sam obiekt o~zewnętrznej łączności. W~kontraście
zmienna \code{const} zdefiniowana w~przestrzeni nazw bez \code{extern}
ma w~typowym przypadku łączność wewnętrzną; oddzielne TU mogą mieć
wówczas odrębne egzemplarze. \code{static} na poziomie przestrzeni
nazw także zapewnia łączność wewnętrzną, lecz na współczesnym C++
zwykle preferujemy anonimową przestrzeń nazw.

\begin{lstlisting}[caption={Nazwy pomocnicze ograniczone do jednej TU.}]
// Plik: helper.cpp
namespace {
    int internal_counter = 0;
    int twice_local(int x) { return x * 2; }
}

int perform(int x) {
    ++internal_counter;
    return twice_local(x);
}
\end{lstlisting}

\subsection{Szablony: dlaczego definicje zwykle znajdują się w~nagłówkach?}

Kompilator instancjuje szablon dla określonych argumentów typowych
w~miejscu, w~którym wymagana jest jego definicja. Samo udostępnienie
deklaracji z~pliku \code{.hpp} i~ukrycie ogólnej definicji szablonu
w~oddzielnym \code{.cpp} może zatem prowadzić do błędu linkera.

\begin{lstlisting}[caption={Najprostszy, tradycyjny sposob udostepniania szablonu.}]
// Plik: multiply.hpp
#pragma once

template<class T>
T multiply(T a, T b) {
    return a * b;
}
\end{lstlisting}

Można także jawnie instancjować \emph{konkretne} typy w~pliku
implementacji. Oddzielne jednostki translacji nie muszą wtedy
instancjować tego samego kodu niezależnie:

\begin{lstlisting}[caption={Jawna instancjacja okreslonego wariantu szablonu.}]
// Plik: operation.hpp
#pragma once
template<class T>
T add(T a, T b);
extern template int add<int>(int, int);

// Plik: operation.cpp
#include "operation.hpp"
template<class T>
T add(T a, T b) { return a + b; }
template int add<int>(int, int);

// Plik: main.cpp
#include "operation.hpp"
int main() { return add(10, 32) == 42 ? 0 : 1; }
\end{lstlisting}

W~tym przykładzie klient może korzystać z~\code{add<int>}, ale
instancjacja innych typów nie zostaje mu udostępniona automatycznie.
Definicje szablonów umieszczane w~interfejsach modułów podlegają
nieco innemu modelowi widoczności, który omówimy dalej.

\subsection{Naruszenie ODR bez łatwego komunikatu}

Nie wszystkie naruszenia ODR skutkują natychmiastowym błędem
linkera. Szczególnie groźne jest warunkowe definiowanie tej samej
klasy w~dwóch TU w~inny sposób, np. przy użyciu sprzecznych
makr kompilacji:

\begin{lstlisting}[caption={Przyk?ad niezgodnych definicji typu (celowo b?edny).}]
// Plik: config.hpp
struct Config {
#ifdef EXTRA_FIELD
    int extra;
#endif
    int value;
};

// Jedna TU z -DEXTRA_FIELD, inna bez tej flagi:
// rozne definicje klasy -> naruszenie ODR.
\end{lstlisting}

Taki program może się nawet zlinkować, lecz pozostaje niepoprawny.
Dlatego opcje zmieniające publiczne definicje klas i~struktur muszą
być spójne we wszystkich TU danego programu.

\section{Symbole, łączność i~typowe błędy linkera}
\label{sec:cpp17-symbols}

\subsection{Łączność a~widoczność symbolu}

\emph{Linkage} w~języku C++ mówi, czy deklaracje w~różnych zakresach
i~TU mogą nazywać tę samą encję. W~praktyce spotykamy między innymi
łączność zewnętrzną, wewnętrzną, a~w~jednostkach modułów także
łączność modułową. Nie należy jednak utożsamiać tej własności
języka z~\emph{visibility} na poziomie ELF ani z~eksportem symbolu
z~biblioteki dynamicznej.

\begin{center}
\begin{tabularx}{\textwidth}{@{}lX@{}}
\toprule
Mechanizm & Co określa?\\
\midrule
\code{static} w~przestrzeni nazw & Łączność wewnętrzna nazwy w~TU.\\
\code{extern} & M.in. deklarowanie obiektu definiowanego w~innym miejscu.\\
Anonimowa przestrzeń nazw & Wewnętrzna łączność elementów pomocniczych.\\
\code{inline} & Warunki wielokrotnych definicji zgodnych z~ODR.\\
\code{-fvisibility=hidden} & Domyślna widoczność symboli w~typowych bibliotekach ELF.\\
\code{__declspec(dllexport)} & Udostępnianie symboli biblioteki DLL w~MSVC.\\
\bottomrule
\end{tabularx}
\end{center}

\subsection{Name mangling i~\texttt{extern~\textquotedbl{}C\textquotedbl{}}}

C++ dopuszcza przeciążanie funkcji, przestrzenie nazw oraz rozbudowane
typy szablonowe. Aby móc odróżniać symbole po linkowaniu,
implementacje zwykle kodują część tych informacji w~nazwach
binarnych (\emph{name mangling}). Dokładny sposób kodowania nie jest
ustalony przez ISO C++.

\begin{lstlisting}[caption={Przeciazanie i punkt wejscia C ABI.}]
namespace calc {
    int f(int x) { return x + 1; }
    double f(double x) { return x + 1.0; }
}

extern "C" int api_add(int a, int b) {
    return a + b;
}
\end{lstlisting}

Na typowej platformie ELF/Itanium ABI funkcje przeciążone otrzymują
różne symbole zakodowane przez kompilator, podczas gdy \code{api_add}
ma nazwę zgodną z~łącznością języka C. Nie jest to gwarancja pełnej
zgodności ABI z~dowolną biblioteką C: znaczenie mają także deklaracje
typów i~konwencja wywołania ustalona dla danej platformy.

\begin{lstlisting}[language=bash,caption={Demanglowanie i inspekcja symboli (przyk?ady narzedziowe).}]
g++ -std=c++20 -c overload.cpp -o overload.o
nm overload.o
nm -C overload.o
c++filt _Z3fooi      # Ilustracja popularnego schematu manglingu
\end{lstlisting}

\subsection{\texttt{undefined reference}, \texttt{multiple definition}}

Błąd \code{undefined reference} oznacza zwykle, że linker znalazł
odwołanie do symbolu, lecz nie znalazł odpowiedniej definicji.
Przykładowe przyczyny to pominięcie \code{math.o}, pominięcie
biblioteki, niezgodny podpis funkcji albo brak jawnej instancjacji
szablonu. Z~kolei \code{multiple definition} sugeruje m.in. zwykłą
funkcję lub zmienną zdefiniowaną w~nagłówku dołączanym do wielu TU.

Przydatna lista kontrolna:
\begin{enumerate}
\item Sprawdź, czy wszystkie wymagane \code{.cpp} są kompilowane i~linkowane.
\item Porównaj podpis funkcji: przestrzeń nazw, typy, kwalifikatory i~ABI.
\item Upewnij się, że deklaracja i~definicja pochodzą z~tego samego interfejsu.
\item Sprawdź kolejność bibliotek statycznych i~ich zależności.
\item Dla klas polimorficznych sprawdź brakujące definicje funkcji
      wirtualnych, w~tym destruktora wirtualnego, który może wymagać
      definicji nawet jako funkcja \code{= 0}.
\item Dla szablonów sprawdź dostępność definicji w~punkcie instancjacji
      lub odpowiednią instancjację jawną.
\end{enumerate}

\section{Biblioteki statyczne}
\label{sec:cpp17-static-libraries}

\emph{Biblioteka statyczna} jest typowo archiwum plików obiektowych.
W~Linux ma zwykle rozszerzenie \code{.a}, w~MSVC \code{.lib}.
Biblioteka sama nie jest osobnym procesem ani automatycznie tworzonym
wątkiem. Linker wybiera z~niej obiekty potrzebne do rozwiązania
zależności i~tworzy wynikowy program lub inną bibliotekę.

\begin{lstlisting}[language=bash,caption={Budowanie i uzycie biblioteki statycznej (Linux/GCC).}]
g++ -std=c++20 -O2 -c math.cpp -o math.o
ar rcs libmathstuff.a math.o

g++ -std=c++20 -O2 -c main.cpp -o main.o
g++ main.o -L. -lmathstuff -o app
./app
\end{lstlisting}

W~typowym poleceniu GNU linkera biblioteki archiwalne warto umieszczać
\emph{po} obiektach, które odwołują się do ich symboli. Przykładowo
\code{g++ main.o -lmathstuff} jest zazwyczaj właściwszym porządkiem
niż \code{g++ -lmathstuff main.o}. Linkowanie zależności cyklicznych
bibliotek statycznych może wymagać ponownego wskazania biblioteki
lub użycia grupy linkera (np.~\code{--start-group}/\code{--end-group}
w~GNU ld); lepiej jednak unikać takich zależności architektonicznie.

\nota{Statyczna biblioteka nie oznacza automatycznie w~pełni
statycznego programu. Aplikacja może dalej korzystać z~bibliotek
współdzielonych systemu i~standardowej biblioteki C++.}

\section{Biblioteki współdzielone i~dynamiczne ładowanie}
\label{sec:cpp17-shared-libraries}

\subsection{Budowanie biblioteki współdzielonej}

Na systemach ELF biblioteki współdzielone są zazwyczaj plikami
\code{.so}. Obiekty wykorzystywane w~takich bibliotekach zwykle
buduje się jako kod niezależny od położenia w~pamięci (\emph{PIC}).

\begin{lstlisting}[language=bash,caption={Biblioteka .so oraz program klienta (Linux/GCC).}]
g++ -std=c++20 -O2 -fPIC -c math.cpp -o math_pic.o
g++ -shared math_pic.o -o libmathstuff.so

g++ -std=c++20 -O2 main.cpp -L. -lmathstuff \
    '-Wl,-rpath,$ORIGIN' -o app
./app
\end{lstlisting}

Wyrażenie \code{\$ORIGIN} odnosi się w~typowym loaderze ELF do
katalogu, w~którym znajduje się ładowany obiekt. Apostrofy w~poleceniu
powłoki zapobiegają wcześniejszej ekspansji \code{\$ORIGIN} przez
powłokę. W~projekcie produkcyjnym należy przemyśleć ścieżki
instalacyjne, pakowanie bibliotek i~bezpieczeństwo ładowania.

Biblioteki współdzielone pozwalają wielu procesom mapować ich kod
z~tego samego pliku oraz wymieniać implementację bez konieczności
pełnego linkowania programu w~czasie jego budowania. Nie gwarantuje
to jednak, że \emph{dowolna} nowa wersja zachowa zgodność ABI.

\subsection{Wyszukiwanie bibliotek: RPATH/RUNPATH, SONAME}

Na Linux/ELF linker i~loader korzystają z~mechanizmów takich jak
\emph{SONAME}, \emph{RPATH} i~\emph{RUNPATH}. Reguły wyszukiwania
zależą od systemu i~konfiguracji; środowisko wykonania może również
uwzględniać ścieżki systemowe lub zmienne środowiskowe.

\begin{lstlisting}[language=bash,caption={Diagnostyka zaleznosci dynamicznych (Linux).}]
readelf -d ./app | grep -E 'NEEDED|RPATH|RUNPATH'
ldd ./app                  # Tylko dla zaufanych plikow
readelf -Ws libmathstuff.so
\end{lstlisting}

Programów nieznanego pochodzenia nie należy lekkomyślnie przekazywać
do narzędzi, które mogą uruchamiać kod lub loader; do statycznej
inspekcji lepiej nadaje się \code{readelf}.

\subsection{Eksportowanie symboli i~kontrola API}

Domyślne eksportowanie zbyt wielu symboli może powodować kolizje
nazw i~powiększać powierzchnię ABI. W~ELF można budować bibliotekę
z~\code{-fvisibility=hidden}, a~publiczne symbole jawnie oznaczać:

\begin{lstlisting}[caption={Przenosne w granicach GCC/Clang/MSVC makro eksportu.}]
// Plik: api.hpp
#pragma once

#if defined(_WIN32)
  #if defined(MYLIB_BUILD_SHARED)
    #define MYLIB_API __declspec(dllexport)
  #elif defined(MYLIB_USE_SHARED)
    #define MYLIB_API __declspec(dllimport)
  #else
    #define MYLIB_API
  #endif
#elif defined(__GNUC__) || defined(__clang__)
  #define MYLIB_API __attribute__((visibility("default")))
#else
  #define MYLIB_API
#endif

extern "C" MYLIB_API int mylib_add(int, int);
\end{lstlisting}

W~MSVC zwykle stosuje się projekt DLL oraz odpowiednią bibliotekę
importową; na macOS biblioteki mają najczęściej rozszerzenie
\code{.dylib}. Flagi \code{-fPIC}, \code{-shared} i~\code{-Wl,...}
w~powyższej postaci dotyczą konkretnego toolchainu, nie ISO C++.

\subsection{\texttt{dlopen}, \texttt{dlsym} i~granica interfejsu wtyczki}

Bibliotekę można także załadować podczas działania procesu, bez
linkowania jej wprost z~aplikacją. Na POSIX często używa się
\code{dlopen}/\code{dlsym}, natomiast na Windows odpowiedników
\code{LoadLibrary}/\code{GetProcAddress}.

\begin{lstlisting}[caption={Prosta wtyczka C ABI: biblioteka i klient POSIX (dwa pliki).}]
// Plik: plugin.cpp
extern "C" int plugin_square(int x) {
    return x * x;
}

// Plik: loader.cpp
#include <dlfcn.h>
#include <iostream>
#include <stdexcept>

int main() {
    void* lib = dlopen("./libplugin.so", RTLD_NOW);
    if (!lib) throw std::runtime_error(dlerror());

    using Fn = int (*)(int);
    void* address = dlsym(lib, "plugin_square");
    if (!address) {
        dlclose(lib);
        throw std::runtime_error("Missing plugin_square");
    }

    // POSIX dopuszcza konwersje dlsym -> wskaznik do funkcji.
    Fn square = reinterpret_cast<Fn>(address);
    std::cout << square(7) << '\n'; // 49
    dlclose(lib);
}
\end{lstlisting}

\begin{lstlisting}[language=bash,caption={Budowanie przyk?adu wtyczki (Linux/ELF).}]
g++ -std=c++20 -fPIC -shared plugin.cpp -o libplugin.so
g++ -std=c++20 loader.cpp -ldl -o loader
./loader
\end{lstlisting}

\nota{To przykład dla POSIX, a~nie w~pełni przenośna konstrukcja
ISO C++. Również rzutowanie wyniku \code{dlsym} na wskaźnik funkcji
opiera się na wymaganiach POSIX, a~nie na ogólnej gwarancji C++.
Dodatkowo nie wolno wyładowywać biblioteki, dopóki jakakolwiek funkcja
lub obiekt z~tej biblioteki nadal może być użyty. Przy bardziej
złożonych wtyczkach potrzebny jest kontrakt wersji ABI i~zarządzania
czasem życia.}

\section{ABI --- Application Binary Interface}
\label{sec:cpp17-abi}

\subsection{API a~ABI: dwa różne kontrakty}

\emph{API} opisuje przede wszystkim dostępne nazwy, typy, funkcje
i~semantykę na poziomie źródła. \emph{ABI} określa, jak te elementy
reprezentowane są w~kodzie maszynowym. ABI obejmuje między innymi:
\begin{itemize}
\item konwencje wywołania: przekazywanie parametrów, wartość zwracana,
      rejestry i~ułożenie stosu;
\item reprezentację i~wyrównanie obiektów, w~tym pól klas;
\item nazwy symboli C++, RTTI, mechanizmy wyjątków i~wirtualnych wywołań;
\item reguły udostępniania i~ładowania symboli bibliotek;
\item ABI standardowej biblioteki oraz ustawienia kompilatora.
\end{itemize}

Dwie biblioteki mogą być zgodne źródłowo i~jednocześnie binarnie
niezgodne. W~szczególności identyczny nagłówek nie oznacza, że plik
\code{.so} zbudowany innym toolchainem lub z~innym zestawem opcji
na pewno współpracuje z~aplikacją.

\subsection{Zmiana układu obiektu łamie zgodność binarną}

Rozważmy klasę publiczną:
\begin{lstlisting}[caption={Zmiana prywatnych pol nadal moze zmienic ABI.}]
// Wersja biblioteki 1
class Record {
    int id_{};
public:
    int id() const;
};

// Wersja biblioteki 2
class Record {
    int id_{};
    double score_{}; // Nawet prywatne pole zmienia uklad.
public:
    int id() const;
};
\end{lstlisting}

Klient kompiluje np.~\code{sizeof(Record)} według nagłówka,
który widział podczas kompilacji. Gdy zamienimy jedynie bibliotekę
na wersję z~innym układem \code{Record}, może dojść do naruszenia
kontraktu pamięci, nawet jeśli \code{id()} nadal ma taki sam podpis.
Podobne ryzyko dotyczy dodawania wirtualnych funkcji i~zmian
hierarchii dziedziczenia.

\subsection{Wzorzec PImpl}

Aby zmniejszyć zależności kompilacyjne i~ukryć układ implementacji,
można przechowywać wskaźnik do niekompletnego typu. Dzięki temu
zmiana \emph{pól implementacyjnych} często nie zmienia rozmiaru
publicznego obiektu.

\begin{lstlisting}[caption={PImpl: plik widget.hpp.}]
#pragma once
#include <memory>

class Widget {
public:
    Widget();
    ~Widget();

    Widget(Widget&&) noexcept;
    Widget& operator=(Widget&&) noexcept;

    Widget(const Widget&) = delete;
    Widget& operator=(const Widget&) = delete;

    int value() const;
    void set_value(int value);

private:
    struct Impl;
    std::unique_ptr<Impl> impl_;
};
\end{lstlisting}

\begin{lstlisting}[caption={PImpl: plik widget.cpp.}]
#include "widget.hpp"
#include <utility>

struct Widget::Impl {
    int value = 0;
    // Tu moga pojawiac sie dalsze pola.
};

Widget::Widget() : impl_(std::make_unique<Impl>()) {}
Widget::~Widget() = default;
Widget::Widget(Widget&&) noexcept = default;
Widget& Widget::operator=(Widget&&) noexcept = default;

int Widget::value() const { return impl_->value; }
void Widget::set_value(int value) { impl_->value = value; }
\end{lstlisting}

Destruktor jest zdefiniowany w~\code{widget.cpp}, gdzie \code{Impl}
jest kompletny; to istotne przy \code{std::unique_ptr<Impl>}.
PImpl nie zapewnia automatycznie pełnej stabilności ABI: nadal
liczy się zgodność kompilatora, standardowej biblioteki, eksportu
symboli, konwencji wywołania i~pozostałych publicznych deklaracji.
Dodatkowy wskaźnik i~alokacja na stercie mogą zwiększać koszt dostępu.

\subsection{C ABI jako ostrożna granica interoperacyjności}

Dla wtyczek i~bibliotek wykorzystywanych z~różnych języków często
projektuje się wąski interfejs oparty na funkcjach \code{extern "C"}
oraz nieprzezroczystych uchwytach (\emph{opaque handles}). Zamiast
przekazywać \code{std::string} lub klasy z~funkcjami wirtualnymi,
przekazuje się proste typy i~wskaźniki.

\begin{lstlisting}[caption={Przyk?adowy szkic stabilniejszego interfejsu binarnego.}]
#ifdef __cplusplus
extern "C" {
#endif

struct engine_handle;

engine_handle* engine_create(void);
void engine_destroy(engine_handle* p);
int engine_eval(engine_handle* p, double x, double* result);

#ifdef __cplusplus
}
#endif
\end{lstlisting}

Podpis funkcji nadal wymaga uzgodnienia ABI danej platformy,
konwencji wywołania i~kontraktu błędów. Przekraczanie takiej granicy
wyjątkami C++, pamięcią zwalnianą przez niezgodny alokator lub obiektami
standardowej biblioteki jest niebezpieczne. Bezpieczniej zwracać
kody błędu, a~alokację i~zwalnianie łączyć w~pary po tej samej
stronie interfejsu.

\subsection{Wersje kompilatora i~biblioteki standardowej}

Istnieją różne rodziny ABI C++ (np. rodzina Itanium ABI stosowana na
wielu platformach uniksowych oraz ABI MSVC w~Windows). Nawet w~ramach
jednej rodziny znaczenie mają flagi kompilatora, ustawienia standardowej
biblioteki, opcje wyjątków, RTTI i~ABI konkretnych typów.
Przykładowo GCC/libstdc++ używa mechanizmu tzw. \emph{dual ABI}
dla niektórych typów standardowych (m.in. \code{std::string}); mieszanie
niezgodnych ustawień może skutkować błędami linkera lub działania.
Nie należy zakładać, że biblioteki zbudowane dowolnymi wersjami
GCC, Clanga i~MSVC są wymienne binarnie.

\section{System budowania: CMake i~projekt wieloplikowy}
\label{sec:cpp17-cmake}

\subsection{Dlaczego targety są ważniejsze niż globalne flagi?}

Ręczne wywoływanie kompilatora dobrze tłumaczy mechanikę, lecz dla
większych projektów potrzebujemy grafu zależności, wariantów
Debug/Release, testów i~kontroli opcji kompilatora. Współczesny
CMake modeluje budowanie przez \emph{targety}: bibliotekę, program,
interfejsową bibliotekę nagłówkową i~ich właściwości.

Przykładowa struktura:
\begin{lstlisting}[caption={Uk?ad plikow klasycznej biblioteki CMake.}]
project/
  CMakeLists.txt
  include/
    mathstuff/math.hpp
  src/
    math.cpp
    main.cpp
\end{lstlisting}

W~tej konkretnej strukturze pliki źródłowe mogą zaczynać się od
\code{\#include "mathstuff/math.hpp"}, a~\code{main.cpp} również
korzysta z~tej samej ścieżki nagłówka. To dzięki
\code{target_include_directories} kompilator znajduje katalog
\code{include} bez wpisywania ręcznych flag \code{-I}.

\begin{lstlisting}[language={},caption={CMakeLists.txt dla klasycznej biblioteki.}]
cmake_minimum_required(VERSION 3.20)
project(MathStuff LANGUAGES CXX)

add_library(mathstuff STATIC src/math.cpp)
target_compile_features(mathstuff PUBLIC cxx_std_20)
target_include_directories(mathstuff PUBLIC
    $<BUILD_INTERFACE:${CMAKE_CURRENT_SOURCE_DIR}/include>
    $<INSTALL_INTERFACE:include>
)

add_executable(app src/main.cpp)
target_link_libraries(app PRIVATE mathstuff)
\end{lstlisting}

\begin{lstlisting}[language=bash,caption={Konfigurowanie i budowanie projektu z CMake.}]
cmake -S . -B build -DCMAKE_BUILD_TYPE=Debug
cmake --build build -j
./build/app
\end{lstlisting}

\code{PRIVATE}, \code{PUBLIC} i~\code{INTERFACE} określają
\emph{propagację wymagań użycia} (usage requirements) między
targetami:
\begin{center}
\begin{tabularx}{\textwidth}{@{}lX@{}}
\toprule
Słowo & Znaczenie\\
\midrule
\code{PRIVATE} & Wymaganie tylko na potrzeby budowania danego targetu.\\
\code{INTERFACE} & Wymaganie wyłącznie dla jego konsumentów.\\
\code{PUBLIC} & Wymaganie dla targetu i~jego konsumentów.\\
\bottomrule
\end{tabularx}
\end{center}

Gdy publiczny nagłówek biblioteki wymaga pewnego katalogu
nagłówków, ten katalog powinien zwykle propagować się do klientów.
Gdy jednak biblioteka korzysta z~zależności jedynie wewnątrz
\code{.cpp}, często wystarcza \code{PRIVATE}. Właściwy zakres
zależy też od tego, czy target jest biblioteką statyczną lub
współdzieloną i~jakie zależności są konieczne przy ostatecznym
linkowaniu.

\subsection{Biblioteka nagłówkowa i~\texttt{INTERFACE}}

\begin{lstlisting}[language={},caption={Header-only target i wymagania dla klientow.}]
add_library(my_headers INTERFACE)
target_include_directories(my_headers INTERFACE
    ${CMAKE_CURRENT_SOURCE_DIR}/include
)
target_compile_features(my_headers INTERFACE cxx_std_20)

add_executable(tool src/tool.cpp)
target_link_libraries(tool PRIVATE my_headers)
\end{lstlisting}

Taki target nie tworzy archiwum \code{.a}. Reprezentuje zestaw
wymagań i~publiczne nagłówki używane przez konsumentów.

\subsection{Diagnostyka, sanitizatory, LTO i~build types}

\begin{lstlisting}[language=bash,caption={Przyk?adowe flagi diagnostyczne GCC/Clang.}]
g++ -std=c++20 -O0 -g -Wall -Wextra -Wpedantic \
    -fsanitize=address,undefined main.cpp -o app_debug

g++ -std=c++20 -O2 -flto main.cpp math.cpp -o app_lto
\end{lstlisting}

\emph{Link-Time Optimization} (LTO) może umożliwiać optymalizacje
pomiędzy TU, np. dodatkowe wstawianie funkcji do miejsc wywołania.
Jest to jednak kosztowna funkcja toolchainu; jej efekty zależą od
programu i~opcji. Profilowanie i~powtarzalne pomiary są ważniejsze
niż sam fakt włączenia \code{-O3} lub \code{-flto}.

\section{Moduły C++20: nowy mechanizm organizacji kodu}
\label{sec:cpp17-modules}

\subsection{Motywacja: problem tekstowego dołączania nagłówków}

W~klasycznym modelu ten sam popularny nagłówek może być parsowany
oddzielnie w~wielu TU. Makra i~kolejność \code{\#include} wpływają
na wynik preprocessingu. Moduły C++20 umożliwiają jawne eksportowanie
deklaracji i~importowanie ich bez tekstowego wklejania całych
interfejsów do każdego klienta.

Zaletami modułów mogą być izolacja makr, bardziej jednoznaczne
zależności i~krótsze czasy kompilacji. Nie ma jednak gwarancji,
że \emph{każdy} program po migracji przyspieszy, a~wdrożenie
wymaga wsparcia kompilatora i~systemu budowania.

\subsection{Pierwszy moduł: \texttt{export module} i~\texttt{import}}

Rozważmy dwa pliki:
\begin{lstlisting}[caption={Interfejs modu?u geometry: plik geometry.cppm.}]
export module geometry;

export int square(int x) {
    return x * x;
}

int hidden_helper(int x) {
    return x + 1;
}
\end{lstlisting}

\begin{lstlisting}[caption={Klient modu?u: plik main.cpp.}]
#include <iostream>
import geometry;

int main() {
    std::cout << square(7) << '\n'; // 49
    // hidden_helper(7); // Nie jest eksportowana.
}
\end{lstlisting}

\code{export module geometry;} deklaruje \emph{główną jednostkę
interfejsu modułu} (\emph{primary module interface unit}).
\code{export int square(...)} udostępnia deklarację klientowi.
Samo \code{import geometry;} nie oznacza linkowania kodu funkcji
z~powietrza: jeśli moduł generuje kod obiektowy, trzeba go również
uwzględnić podczas linkowania.

\subsection{BMI/CMI: moduł nie jest biblioteką binarną}

Kompilator wytwarza dla interfejsu modułu własną reprezentację,
często nazywaną \emph{BMI} (\emph{Binary Module Interface}) lub
\emph{CMI} (\emph{Compiled Module Interface}). W~Clangu spotyka się
pliki \code{.pcm}, w~GCC \code{.gcm}, a~w~MSVC pliki \code{.ifc}.
Nie są to uniwersalne, przenośne między kompilatorami artefakty
podobne do biblioteki \code{.a}. Importer musi znać interfejs
przetworzony przez zgodną konfigurację kompilatora.

Zależności muszą być budowane w~poprawnym porządku:
\[
\text{źródło interfejsu modułu}
\longrightarrow \text{BMI/CMI}
\longrightarrow \text{kompilacja importera}
\longrightarrow \text{linkowanie obiektów}.
\]

\subsection{Rzeczywiste komendy: Clang i~GCC}

Dla prostego modułu \code{geometry.cppm} w~Clangu obsługującym
standardowe moduły można zastosować:
\begin{lstlisting}[language=bash,caption={Budowanie przyk?adowego modu?u w Clangu (tryb BMI).}]
clang++ -std=c++20 geometry.cppm --precompile -o geometry.pcm
clang++ -std=c++20 -c geometry.pcm -o geometry.o
clang++ -std=c++20 -fmodule-file=geometry=geometry.pcm \
    -c main.cpp -o main.o
clang++ geometry.o main.o -o app
./app
\end{lstlisting}

Dla GCC~14 (wsparcie modułów C++20 było w~tej wersji nadal rozwijane)
trzeba dodatkowo włączyć eksperymentalny mechanizm
\code{-fmodules-ts}. W~przykładach z~GCC~14 przydatne bywa jawne
wskazanie języka przy rozszerzeniu \code{.cppm}:
\begin{lstlisting}[language=bash,caption={Budowanie tego samego modu?u z GCC 14.}]
g++ -std=c++20 -fmodules-ts -x c++ -c geometry.cppm \
    -o geometry.o
# GCC tworzy potrzebne dane interfejsu w gcm.cache/.
g++ -std=c++20 -fmodules-ts -c main.cpp -o main.o
g++ geometry.o main.o -o app
./app
\end{lstlisting}

Dokładne opcje i~formaty modułów mogą zmieniać się między wersjami
narzędzi. W~dużym projekcie lepiej powierzyć zależności systemowi
budowania obsługującemu skanowanie importów, zamiast ręcznie
zarządzać plikami BMI.

\subsection{Jednostka interfejsu i~jednostka implementacji}

Moduł może rozdzielać interfejs od implementacji bez tradycyjnego
nagłówka. Jednostka implementacji deklaruje \code{module name;},
a~nie \code{export module name;}:

\begin{lstlisting}[caption={Dwa pliki jednego modu?u: interface.cppm i implementation.cpp.}]
// Plik: engine.cppm
export module engine;
export int run(int value);

// Plik: engine.cpp
module engine;
int run(int value) { return 2 * value; }
\end{lstlisting}

Importer widzi tylko deklaracje eksportowane przez interfejs.
Jednostkę implementacji również trzeba skompilować i~zlinkować.
Konkretny system budowania musi uwzględnić fakt, że jednostka
implementacji importuje interfejs swojego modułu niejawnie.

\subsection{Global module fragment i~nagłówki klasyczne}

Tradycyjny nagłówek może wymagać makr konfigurujących jego treść.
W~takim przypadku można skorzystać z~\emph{global module fragment}:

\begin{lstlisting}[caption={Global module fragment: tradycyjny include przed modu?em.}]
module;
#include <string>
#include <utility>

export module greeting;

export std::string greet(std::string name) {
    return "Hello " + std::move(name);
}
\end{lstlisting}

Dyrektywa \code{module;} rozpoczyna fragment globalny, a~następna
deklaracja \code{export module greeting;} rozpoczyna właściwą
jednostkę modułu. Importowanie nagłówków jednostkowych przez
\code{import <string>;} jest innym mechanizmem, tzw. \emph{header units},
których dostępność zależy od toolchainu. Nie utożsamiajmy go z~nazwanym
modułem standardowej biblioteki \code{import std;}.

\subsection{Partycje i~ponowny eksport importu}

Duży moduł można podzielić na partycje. Nazwa po dwukropku oznacza
partycję wewnątrz jednego modułu, a~nie odrębny nazwany moduł.

\begin{lstlisting}[caption={Dwie jednostki interfejsu modu?u z partycja.}]
// Plik: math_detail.cppm
export module mathlib:detail;
export int triple(int x) { return 3 * x; }

// Plik: mathlib.cppm
export module mathlib;
export import :detail;
export int times_six(int x) { return 2 * triple(x); }

// Plik: main.cpp
import mathlib;
int main() { return times_six(7) == 42 ? 0 : 1; }
\end{lstlisting}

Klient importuje \code{mathlib}, a~nie \code{mathlib:detail}.
W~tym przykładzie \code{export import :detail;} re-eksportuje
eksporty partycji. Wymaga to prawidłowej kolejności budowania
\emph{najpierw partycji, potem głównego interfejsu, a~na końcu klienta}.

\subsection{Łączność modułowa i~ODR}

C++20 wprowadza reguły własności deklaracji (\emph{module ownership}).
Deklaracje przypisane do nazwanego modułu należą do niego, a~pewne
nieeksportowane nazwy uzyskują \emph{module linkage}. Różni się to od
mechanicznego zastąpienia \code{\#include} przez \code{import}.

Przykładowo dodanie \code{export} nie sprawia automatycznie, że
interfejs C++ będzie zgodny ABI między kompilatorami. Ponadto
źródłowy kod szablonów nadal musi być semantycznie dostępny
odpowiedniemu mechanizmowi instancjacji, choć nie musi być już
tekstowo wklejany przez \code{\#include}.

\subsection{CMake i~\texttt{FILE\_SET CXX\_MODULES}}

Od CMake~3.28 dostępne są mechanizmy skanowania modułów i~specjalny
typ zbioru plików \code{CXX_MODULES}. Przykład projektu:

\begin{lstlisting}[language={},caption={CMakeLists.txt projektu korzystajacego z modu?ow C++20.}]
cmake_minimum_required(VERSION 3.28)
project(GeometryModules LANGUAGES CXX)

add_library(geometry)
target_compile_features(geometry PUBLIC cxx_std_20)
target_sources(geometry PUBLIC
    FILE_SET CXX_MODULES FILES geometry.cppm
)

add_executable(app main.cpp)
target_link_libraries(app PRIVATE geometry)
\end{lstlisting}

\begin{lstlisting}[language=bash,caption={Konfiguracja modu?ow przy uzyciu generatora Ninja.}]
cmake -S . -B build -G Ninja
cmake --build build
./build/app
\end{lstlisting}

Ten przykład pokazuje \emph{składnię} CMake, ale jego zbudowanie
wymaga zgodnego zestawu wersji CMake, kompilatora, generatora i~skanera
zależności. W~szczególności kompilator akceptujący składnię modułów
niekoniecznie współpracuje z~automatycznym skanowaniem w~danej
konfiguracji CMake. Na współczesnych systemach używa się m.in.
Ninja (co najmniej 1.11 w~typowych konfiguracjach) i~odpowiednich
wersji Clanga lub MSVC. Obsługa modułów standardowej biblioteki oraz
\emph{header units} ma odrębne ograniczenia.

\subsection{\texttt{import std;} nie jest tym samym co \texttt{\#include}}

Sam standard C++20 wprowadza językowy mechanizm modułów, ale
\emph{nie zobowiązuje} każdej biblioteki standardowej do dostarczenia
nazwanego modułu \code{std} w~tym standardzie. \code{import std;}
jest standaryzowane w~C++23; obsługa zależy od kompilatora, dostawcy
biblioteki standardowej i~wersji systemu budowania.

\begin{lstlisting}[caption={Nazwany modu? standardowej biblioteki (C++23, zalezny od wsparcia).}]
import std;

int main() {
    std::vector<int> xs{1, 2, 3};
    std::println("size = {}", xs.size());
}
\end{lstlisting}

Powyższy przykład wykorzystuje także \code{std::println} z~C++23.
Nie należy automatycznie kompilować go z~\code{-std=c++20} ani
zakładać, że \code{import std;} działa w~każdym środowisku C++23.

\section{Optymalizacja procesu budowania}
\label{sec:cpp17-build-perf}

\subsection{Incremental build i~zależności nagłówkowe}

Zmiana publicznego nagłówka może wymusić ponowną kompilację wielu TU,
podczas gdy zmiana prywatnego \code{.cpp} zwykle dotyka mniejszej
części projektu. Dlatego ograniczanie niepotrzebnych dołączeń
nagłówków, stosowanie deklaracji wyprzedzających tam, gdzie są
prawidłowe, i~PImpl mogą poprawiać czas budowania.

Moduły zmieniają strukturę zależności, ale nie usuwają konieczności
ich śledzenia; importer musi otrzymać aktualną reprezentację
interfejsu modułu.

\subsection{Precompiled headers, unity builds, ccache}

Popularne techniki przyspieszania kompilacji to:
\begin{description}[style=nextline]
\item[Precompiled headers (PCH)]
Przechowywanie wyniku parsowania stabilnych, obszernych nagłówków.
Pomagają w~projektach z~dużą liczbą TU, ale wymagają spójności
opcji kompilacji i~zarządzania unieważnianiem pamięci podręcznej.
\item[Unity builds]
Łączenie kilku plików źródłowych we wspólną TU. Może ograniczać
narzut kompilacji, ale też ujawniać kolizje anonimowych przestrzeni
nazw, makr i~definicji o~łączności wewnętrznej.
\item[Compilation cache]
Narzędzia takie jak \code{ccache} unikają ponownego wykonywania
kompilacji, jeżeli odpowiednio zdefiniowane wejścia pozostają takie same.
\item[Parallel build]
Kompilowanie niezależnych TU równolegle (np. \code{cmake --build build -j})
wykorzystuje wiele rdzeni, ale zwiększa chwilowe użycie pamięci.
\end{description}

\subsection{Powtarzalność buildów i~zależności transytywne}

Dla bibliotek publicznych warto utrzymywać przejrzystą strukturę
zależności. Przykładowo użytkownik nie powinien polegać na tym,
że \code{\#include <vector>} ,,przypadkiem'' pojawiło się wskutek
dołączenia innego nagłówka. Każdy nagłówek powinien zawierać
potrzebne deklaracje w~sposób niezależny od kolejności dołączeń.

Warto również zapisywać wersje kompilatora i~bibliotek, flagi
wpływające na ABI, wariant budowania oraz sposób instalowania
artefaktów. \emph{Reproducible builds} wymagają dodatkowo kontroli
źródeł niedeterministyczności, np. znaczników czasu i~ścieżek
zapisywanych w~plikach wynikowych.

\section{Praktyczny przewodnik diagnostyczny}
\label{sec:cpp17-troubleshooting}

\begin{center}
\small
\begin{tabularx}{\textwidth}{@{}p{0.34\textwidth}X@{}}
\toprule
Objaw & Najbardziej prawdopodobne źródło problemu\\
\midrule
\code{undefined reference} & Brakujący plik obiektowy, biblioteka,
definicja albo instancjacja szablonu.\\
\code{multiple definition} & Niepoprawna definicja zwykłej
funkcji lub zmiennej w~nagłówku albo podwójnie linkowany obiekt.\\
\code{undefined reference to vtable} & Niezdefiniowana wymagana
funkcja wirtualna (często destruktor) lub brak pliku implementacji.\\
\code{cannot find -lfoo} & Brak biblioteki lub nieprawidłowe
katalogi przeszukiwania linkera.\\
\code{cannot open shared object file} & Problem ścieżki lub
zależności ładowania, niekoniecznie problem kompilacji.\\
Kompilacja klienta działa, program kończy się błędem & Możliwa
niezgodność ABI, naruszenie ODR, zły czas życia obiektu lub błąd programu.\\
Importer nie znajduje modułu & Nie zbudowano BMI/CMI, nie wskazano
ścieżki albo konfiguracja skanera nie obsługuje importu.\\
Import działa, ale linker zgłasza brak funkcji & Brak pliku
obiektowego modułu/implementacji lub potrzebnej biblioteki.\\
Zmiana nagłówka przebudowuje prawie wszystko & Silne zależności
nagłówkowe; rozważ redukcję dołączeń, PImpl, moduły lub PCH.\\
\bottomrule
\end{tabularx}
\end{center}

\subsection{Kiedy wybrać nagłówki, bibliotekę czy moduł?}

\begin{itemize}
\item \textbf{Header-only} sprawdza się dla małych narzędzi, szablonów
      i~funkcji \code{inline}; ceną mogą być zależności kompilacyjne.
\item \textbf{Biblioteka statyczna} upraszcza dystrybucję kodu obiektowego
      wewnątrz aplikacji, ale nie usuwa wymagań ABI w~czasie linkowania.
\item \textbf{Biblioteka współdzielona} pozwala współdzielić kod
      i~wymieniać implementację, lecz wymaga jasnej polityki ABI,
      widoczności symboli i~wersjonowania.
\item \textbf{Moduł C++20} służy do semantycznego podziału interfejsu
      źródłowego i~kontroli importów; nie jest alternatywą wobec
      procesu linkowania ani automatycznym formatem dystrybucji ABI.
\item \textbf{PImpl i~C API} rozważaj na granicach bibliotek o~długim
      cyklu życia, szczególnie gdy stabilność ABI ma duże znaczenie.
\end{itemize}

\section*{Podsumowanie rozdziału}
\addcontentsline{toc}{section}{Podsumowanie rozdziału}

Kompletny program C++ jest wynikiem wspólnej pracy preprocessora,
kompilatora, linkera i~loadera. Organizując projekt, trzeba osobno
rozumieć \emph{deklaracje i~definicje}, \emph{ODR i~łączność},
\emph{symbole i~biblioteki} oraz \emph{ABI}. Nagłówki i~moduły
pomagają udostępniać interfejsy źródłowe; biblioteki statyczne
i~współdzielone dostarczają kod do połączenia lub załadowania.

Dobra organizacja projektu polega na jasnym rozdziale API i~implementacji,
ograniczaniu zależności, korzystaniu z~targetów CMake oraz świadomym
wersjonowaniu interfejsu binarnego. Nowoczesne moduły upraszczają część
problemów nagłówków, ale wprowadzają wymagania dotyczące narzędzi,
BMI/CMI i~kolejności kompilacji.

\section*{Literatura i~dokumentacja}
\addcontentsline{toc}{section}{Literatura i dokumentacja}
\begin{itemize}
\item ISO/IEC 14882:2020, zagadnienia \emph{translation units},
      \emph{linkage}, \emph{One Definition Rule} i~\emph{modules}.
\item \url{https://en.cppreference.com/w/cpp/language/definition}
      --- ODR, deklaracje i~definicje.
\item \url{https://en.cppreference.com/w/cpp/language/modules}
      --- składnia i~semantyka modułów.
\item \url{https://gcc.gnu.org/onlinedocs/gcc/C_002b_002b-Modules.html}
      --- implementacja modułów w~GCC.
\item \url{https://clang.llvm.org/docs/StandardCPlusPlusModules.html}
      --- moduły w~Clangu, BMI i~komendy budowania.
\item \url{https://cmake.org/cmake/help/latest/manual/cmake-cxxmodules.7.html}
      --- CMake i~zależności modułów.
\item \url{https://cmake.org/cmake/help/latest/command/target_sources.html}
      --- \code{FILE_SET CXX_MODULES}.
\item \url{https://gcc.gnu.org/onlinedocs/gcc/Code-Gen-Options.html}
      --- widoczność symboli w~ELF.
\item John Lakos, \emph{Large-Scale C++ Software Design}, Addison-Wesley, 1996.
\end{itemize}

% Komendy powloki sa przeznaczone glownie dla GCC/Clang na Linux/ELF;
% roznice Windows/MSVC i macOS zaznaczono w tekscie.

\chapter{Kompilacja, linkowanie, ABI i moduły C++20}
\label{chap:cpp17-build-abi-modules}

\section{Dlaczego trzeba rozumieć proces budowania?}

W~poprzednich rozdziałach projektowaliśmy klasy, szablony, mechanizmy
zarządzania pamięcią i~współbieżność. Dopóki cały program mieści się
w~jednym pliku, wystarcza zwykle wywołanie \code{g++ main.cpp -o app}.
W~realnym projekcie kod dzieli się jednak na dziesiątki lub tysiące
\emph{jednostek translacji}, biblioteki zewnętrzne i~moduły. To wtedy
pojawiają się komunikaty \code{undefined reference}, \code{multiple definition},
konflikty wersji bibliotek, długi czas budowania oraz błędy ujawniające
się dopiero po wymianie biblioteki współdzielonej.

Musimy rozdzielić trzy różne zagadnienia:
\begin{description}[style=nextline]
\item[Poprawność programu w~języku C++]
Czy deklaracje, definicje, typy i~wywołania spełniają reguły C++,
w~szczególności \emph{One Definition Rule} (ODR)?
\item[Model budowania i~linkowania]
Jakie pliki są kompilowane osobno, gdzie kompilator zapisuje symbole,
w~jakiej kolejności linker przeszukuje biblioteki oraz jak ładowane są
biblioteki dynamiczne?
\item[Zgodność binarna (ABI)]
Czy dwa osobno zbudowane komponenty zgadzają się co do reprezentacji
typów, sposobu wywoływania funkcji, nazw symboli i~pozostałych reguł
interfejsu binarnego?
\end{description}

\nota{Norma ISO C++ opisuje semantykę programu, ale nie narzuca jednego
formatu pliku obiektowego, konwencji \emph{name mangling}, organizacji
bibliotek \code{.so}/\code{.dll} ani poleceń linkera. Poniżej
odróżniamy reguły języka od typowych rozwiązań narzędziowych.}

\section{Od pliku źródłowego do programu}

\begin{center}\small
\begin{tabularx}{\textwidth}{@{}p{0.31\textwidth}X@{}}
\toprule
Etap & Za co odpowiada \\
\midrule
Preprocesor & Rozwija dyrektywy \code{#include}, makra i warunkową kompilację. \\
Kompilator & Sprawdza składnię i typy oraz wytwarza kod dla jednostki translacji. \\
Asembler & Zamienia kod asemblera na plik obiektowy z symbolami i relokacjami. \\
Linker & Łączy pliki obiektowe i biblioteki, rozwiązując odwołania do symboli. \\
Loader systemowy & Ładuje program i wymagane biblioteki współdzielone przed uruchomieniem. \\
\bottomrule
\end{tabularx}
\end{center}
\label{sec:cpp17-pipeline}

\subsection{Preprocesor, kompilator, asembler, linker i~loader}

W~uproszczeniu klasyczny projekt z~nagłówkami przechodzi przez etapy:
\[
\begin{aligned}
\text{źródła i nagłówki}&\ \longrightarrow\ \text{preprocessing}
\ \longrightarrow\ \text{kompilacja}\\
&\phantom{\ \longrightarrow\ }\longrightarrow\ \text{plik obiektowy}
\ \longrightarrow\ \text{linkowanie}
\ \longrightarrow\ \text{program}.
\end{aligned}
\]

Na niektórych platformach kompilator wytwarza tekst asemblerowy
(\code{.s}) jako etap pośredni; w~innych przejście do kodu maszynowego
jest zintegrowane. Linkowanie może odbywać się przy budowaniu programu,
a~dodatkowe rozwiązywanie symboli może następować w~czasie ładowania
lub później.

Prześledźmy etapy na samodzielnym programie.

\begin{lstlisting}[caption={Plik hello.cpp: prosty program do obserwowania kompilacji.}]
#include <iostream>

int twice(int x) { return 2 * x; }

int main() {
    std::cout << twice(21) << '\n';
}
\end{lstlisting}

\begin{lstlisting}[language=bash,caption={Wybrane etapy budowania przy uzyciu GCC/Clang.}]
g++ -std=c++20 -E hello.cpp -o hello.ii  # Preprocesor
# Produkt poprzedniego etapu moze byc duzy.
g++ -std=c++20 -S hello.cpp -o hello.s   # Asembler

g++ -std=c++20 -c hello.cpp -o hello.o   # Kod obiektowy
g++ hello.o -o hello                      # Linkowanie
./hello                                   # 42
\end{lstlisting}

Opcja \code{-c} nie tworzy wykonywalnego programu; produkuje plik
obiektowy, który może zawierać nierozwiązane odwołania do symboli
pochodzących z~innych jednostek translacji. Polecenie \code{-S}
kończy budowanie na asemblerze, a~\code{-E} na wstępnym przetwarzaniu.
Warto odróżnić program sterujący kompilacją (np.~\code{g++}) od samego
frontendu C++, asemblera i~linkera.

\subsection{Pliki obiektowe, symbole i~relokacje}

Plik obiektowy zawiera zwykle kod maszynowy, dane, tablice symboli
oraz \emph{relokacje}, czyli informacje, które adresy muszą zostać
skorygowane po ustaleniu ostatecznego układu programu. W~środowisku
Linux powszechny jest format ELF; Windows używa PE/COFF, a~macOS
Mach-O. Szczegóły tych formatów nie są częścią języka C++.

\begin{lstlisting}[language=bash,caption={Ogladanie symboli i relokacji (Linux/ELF).}]
nm -C hello.o                 # Symbole; -C demangluje nazwy C++
readelf -Ws hello.o          # Tablica symboli ELF
readelf -r hello.o           # Relokacje
objdump -drC hello.o         # Asembler i relokacje
\end{lstlisting}

Przykładowo litera \code{T} w~typowym wyjściu \code{nm} informuje
o~symbolu definiowanym w~sekcji kodu, a~\code{U} o~symbolu
niezdefiniowanym w~danym obiekcie, który trzeba rozwiązać. Nie
należy traktować oznaczeń \code{nm} jako przenośnej własności C++.

\subsection{Jednostka translacji i~kompilacja oddzielna}

W~tradycyjnym modelu \emph{jednostka translacji} (\emph{translation unit},
TU) powstaje z~pliku źródłowego po uwzględnieniu dyrektyw preprocesora,
w~szczególności zawartości dołączanych nagłówków. Każda TU jest
analizowana niezależnie, a~linker łączy otrzymane obiekty.

Rozważmy projekt:
\begin{lstlisting}[caption={Trzy pliki klasycznego projektu C++: API, implementacja i klient.}]
// Plik: math.hpp
#ifndef CPP17_MATH_HPP
#define CPP17_MATH_HPP
int square(int x);
#endif

// Plik: math.cpp
#include "math.hpp"
int square(int x) { return x * x; }

// Plik: main.cpp
#include "math.hpp"
#include <iostream>
int main() { std::cout << square(7) << '\n'; }
\end{lstlisting}

\begin{lstlisting}[language=bash,caption={Niezalezne budowanie TU i koncowe linkowanie.}]
g++ -std=c++20 -Wall -Wextra -c math.cpp -o math.o
g++ -std=c++20 -Wall -Wextra -c main.cpp -o main.o
g++ main.o math.o -o app
./app                        # 49
\end{lstlisting}

\code{math.hpp} dostarcza deklarację klientowi i~implementacji,
ale nie jest zwykle oddzielnym obiektem linkowanym. Zmiana samego
\code{math.cpp} nie musi wymuszać ponownej kompilacji \code{main.cpp},
o~ile nie zmienił się kontrakt w~nagłówku lub pozostałe zależności.

\subsection{Przykładowe komendy w~MSVC}

W~\emph{Developer Command Prompt for Visual Studio} można zbudować
analogiczny przykład za pomocą \code{cl.exe} i~\code{link.exe}:

\begin{lstlisting}[language=bash,caption={Kompilacja dwóch plików w środowisku MSVC/Windows.}]
cl /std:c++20 /EHsc /c math.cpp
cl /std:c++20 /EHsc /c main.cpp
link main.obj math.obj /OUT:app.exe
app.exe
\end{lstlisting}

Przy budowaniu programów C++ w~Visual Studio warto również
kontrolować zgodność wersji MSVC, ustawienia biblioteki uruchomieniowej
(np.~\code{/MD} wobec \code{/MT}) i~konfigurację Debug/Release.
Takie flagi nie mają bezpośrednich odpowiedników narzuconych przez
standard języka; należą do konkretnego środowiska.

\section{Nagłówki, makra i~zasada jednej definicji}
\label{sec:cpp17-headers-odr}

\subsection{Include guards i~\texttt{\#pragma once}}

Zawartość \code{\#include} jest w~uproszczeniu wstawiana do źródła
przez preprocesor. Wielokrotne dołączenie tego samego nagłówka
w~jednej TU mogłoby więc prowadzić do powtórzenia definicji klasy.
Chronimy się przez \emph{include guard}:

\begin{lstlisting}[caption={Klasyczna os?ona nag?owka.}]
#ifndef LIB_WIDGET_HPP
#define LIB_WIDGET_HPP

struct Widget {
    int value{};
};

#endif // LIB_WIDGET_HPP
\end{lstlisting}

W~praktyce popularna jest też dyrektywa \code{\#pragma once},
obsługiwana przez najważniejsze kompilatory, lecz niestandaryzowana
przez ISO C++. Oba rozwiązania chronią \emph{jedną TU} przed
powtórnym dołączeniem pliku. Nie rozwiązują automatycznie problemu
\emph{wielu definicji tego samego nie-inline symbolu w~różnych TU}.

\subsection{ODR: deklaracja kontra definicja}

Deklaracja wprowadza nazwę i~informację o~typie, natomiast definicja
zwykle dostarcza byt danej funkcji lub obiektu. Deklaracji pewnej
funkcji może być wiele, ale \emph{nie-inline} funkcja z~zewnętrzną
łącznością, używana w~programie w~sposób wymagający definicji,
powinna mieć jedną definicję w~programie. Reguły ODR zawierają
wyjątki pozwalające na identyczne definicje w~różnych TU, np. dla
odpowiednich funkcji \code{inline}, klas i~szablonów.

\begin{lstlisting}[caption={Typowa pu?apka: definicja zwyk?ej funkcji w nag?owku.}]
// BAD: wrong.hpp, dolaczony do dwoch plikow .cpp
int twice(int x) { return 2 * x; }
// Dwie zewnetrzne definicje -> naruszenie ODR.
\end{lstlisting}

\begin{lstlisting}[caption={Definicja inline przeznaczona do nag?owka.}]
// Plik: twice.hpp
#pragma once
inline int twice(int x) { return 2 * x; }
\end{lstlisting}

\code{inline} nie nakazuje kompilatorowi wstawienia instrukcji funkcji
bezpośrednio w~miejscu wywołania. Jego znaczenie dla organizacji
programu polega przede wszystkim na dopuszczeniu definicji w~wielu TU
zgodnie z~ODR (z~odpowiednimi warunkami). Wszystkie takie definicje
muszą być zgodne i~spełniać reguły ODR; samo dopisanie \code{inline}
nie naprawia różniących się implementacji w~dwóch plikach.

\subsection{Zmienne \texttt{extern}, \texttt{inline}, \texttt{const} i~\texttt{static}}

\begin{lstlisting}[caption={Jedna definicja zmiennej globalnej, wiele deklaracji.}]
// Plik: settings.hpp
#pragma once
extern int request_limit;       // Tylko deklaracja

// Plik: settings.cpp
#include "settings.hpp"
int request_limit = 100;        // Definicja

// Plik: main.cpp
#include "settings.hpp"
int main() { return request_limit == 100 ? 0 : 1; }
\end{lstlisting}

Od C++17 w~nagłówkach można również definiować \emph{zmienne inline}:
\begin{lstlisting}[caption={Zmienna inline w nag?owku (C++17).}]
#pragma once
inline constexpr int default_timeout_ms = 250;
inline int global_counter = 0;
\end{lstlisting}

W~tym przypadku odpowiednie definicje we wszystkich TU mogą
wskazywać na ten sam obiekt o~zewnętrznej łączności. W~kontraście
zmienna \code{const} zdefiniowana w~przestrzeni nazw bez \code{extern}
ma w~typowym przypadku łączność wewnętrzną; oddzielne TU mogą mieć
wówczas odrębne egzemplarze. \code{static} na poziomie przestrzeni
nazw także zapewnia łączność wewnętrzną, lecz na współczesnym C++
zwykle preferujemy anonimową przestrzeń nazw.

\begin{lstlisting}[caption={Nazwy pomocnicze ograniczone do jednej TU.}]
// Plik: helper.cpp
namespace {
    int internal_counter = 0;
    int twice_local(int x) { return x * 2; }
}

int perform(int x) {
    ++internal_counter;
    return twice_local(x);
}
\end{lstlisting}

\subsection{Szablony: dlaczego definicje zwykle znajdują się w~nagłówkach?}

Kompilator instancjuje szablon dla określonych argumentów typowych
w~miejscu, w~którym wymagana jest jego definicja. Samo udostępnienie
deklaracji z~pliku \code{.hpp} i~ukrycie ogólnej definicji szablonu
w~oddzielnym \code{.cpp} może zatem prowadzić do błędu linkera.

\begin{lstlisting}[caption={Najprostszy, tradycyjny sposob udostepniania szablonu.}]
// Plik: multiply.hpp
#pragma once

template<class T>
T multiply(T a, T b) {
    return a * b;
}
\end{lstlisting}

Można także jawnie instancjować \emph{konkretne} typy w~pliku
implementacji. Oddzielne jednostki translacji nie muszą wtedy
instancjować tego samego kodu niezależnie:

\begin{lstlisting}[caption={Jawna instancjacja okreslonego wariantu szablonu.}]
// Plik: operation.hpp
#pragma once
template<class T>
T add(T a, T b);
extern template int add<int>(int, int);

// Plik: operation.cpp
#include "operation.hpp"
template<class T>
T add(T a, T b) { return a + b; }
template int add<int>(int, int);

// Plik: main.cpp
#include "operation.hpp"
int main() { return add(10, 32) == 42 ? 0 : 1; }
\end{lstlisting}

W~tym przykładzie klient może korzystać z~\code{add<int>}, ale
instancjacja innych typów nie zostaje mu udostępniona automatycznie.
Definicje szablonów umieszczane w~interfejsach modułów podlegają
nieco innemu modelowi widoczności, który omówimy dalej.

\subsection{Naruszenie ODR bez łatwego komunikatu}

Nie wszystkie naruszenia ODR skutkują natychmiastowym błędem
linkera. Szczególnie groźne jest warunkowe definiowanie tej samej
klasy w~dwóch TU w~inny sposób, np. przy użyciu sprzecznych
makr kompilacji:

\begin{lstlisting}[caption={Przyk?ad niezgodnych definicji typu (celowo b?edny).}]
// Plik: config.hpp
struct Config {
#ifdef EXTRA_FIELD
    int extra;
#endif
    int value;
};

// Jedna TU z -DEXTRA_FIELD, inna bez tej flagi:
// rozne definicje klasy -> naruszenie ODR.
\end{lstlisting}

Taki program może się nawet zlinkować, lecz pozostaje niepoprawny.
Dlatego opcje zmieniające publiczne definicje klas i~struktur muszą
być spójne we wszystkich TU danego programu.

\section{Symbole, łączność i~typowe błędy linkera}
\label{sec:cpp17-symbols}

\subsection{Łączność a~widoczność symbolu}

\emph{Linkage} w~języku C++ mówi, czy deklaracje w~różnych zakresach
i~TU mogą nazywać tę samą encję. W~praktyce spotykamy między innymi
łączność zewnętrzną, wewnętrzną, a~w~jednostkach modułów także
łączność modułową. Nie należy jednak utożsamiać tej własności
języka z~\emph{visibility} na poziomie ELF ani z~eksportem symbolu
z~biblioteki dynamicznej.

\begin{center}
\begin{tabularx}{\textwidth}{@{}lX@{}}
\toprule
Mechanizm & Co określa?\\
\midrule
\code{static} w~przestrzeni nazw & Łączność wewnętrzna nazwy w~TU.\\
\code{extern} & M.in. deklarowanie obiektu definiowanego w~innym miejscu.\\
Anonimowa przestrzeń nazw & Wewnętrzna łączność elementów pomocniczych.\\
\code{inline} & Warunki wielokrotnych definicji zgodnych z~ODR.\\
\code{-fvisibility=hidden} & Domyślna widoczność symboli w~typowych bibliotekach ELF.\\
\code{__declspec(dllexport)} & Udostępnianie symboli biblioteki DLL w~MSVC.\\
\bottomrule
\end{tabularx}
\end{center}

\subsection{Name mangling i~\texttt{extern~\textquotedbl{}C\textquotedbl{}}}

C++ dopuszcza przeciążanie funkcji, przestrzenie nazw oraz rozbudowane
typy szablonowe. Aby móc odróżniać symbole po linkowaniu,
implementacje zwykle kodują część tych informacji w~nazwach
binarnych (\emph{name mangling}). Dokładny sposób kodowania nie jest
ustalony przez ISO C++.

\begin{lstlisting}[caption={Przeciazanie i punkt wejscia C ABI.}]
namespace calc {
    int f(int x) { return x + 1; }
    double f(double x) { return x + 1.0; }
}

extern "C" int api_add(int a, int b) {
    return a + b;
}
\end{lstlisting}

Na typowej platformie ELF/Itanium ABI funkcje przeciążone otrzymują
różne symbole zakodowane przez kompilator, podczas gdy \code{api_add}
ma nazwę zgodną z~łącznością języka C. Nie jest to gwarancja pełnej
zgodności ABI z~dowolną biblioteką C: znaczenie mają także deklaracje
typów i~konwencja wywołania ustalona dla danej platformy.

\begin{lstlisting}[language=bash,caption={Demanglowanie i inspekcja symboli (przyk?ady narzedziowe).}]
g++ -std=c++20 -c overload.cpp -o overload.o
nm overload.o
nm -C overload.o
c++filt _Z3fooi      # Ilustracja popularnego schematu manglingu
\end{lstlisting}

\subsection{\texttt{undefined reference}, \texttt{multiple definition}}

Błąd \code{undefined reference} oznacza zwykle, że linker znalazł
odwołanie do symbolu, lecz nie znalazł odpowiedniej definicji.
Przykładowe przyczyny to pominięcie \code{math.o}, pominięcie
biblioteki, niezgodny podpis funkcji albo brak jawnej instancjacji
szablonu. Z~kolei \code{multiple definition} sugeruje m.in. zwykłą
funkcję lub zmienną zdefiniowaną w~nagłówku dołączanym do wielu TU.

Przydatna lista kontrolna:
\begin{enumerate}
\item Sprawdź, czy wszystkie wymagane \code{.cpp} są kompilowane i~linkowane.
\item Porównaj podpis funkcji: przestrzeń nazw, typy, kwalifikatory i~ABI.
\item Upewnij się, że deklaracja i~definicja pochodzą z~tego samego interfejsu.
\item Sprawdź kolejność bibliotek statycznych i~ich zależności.
\item Dla klas polimorficznych sprawdź brakujące definicje funkcji
      wirtualnych, w~tym destruktora wirtualnego, który może wymagać
      definicji nawet jako funkcja \code{= 0}.
\item Dla szablonów sprawdź dostępność definicji w~punkcie instancjacji
      lub odpowiednią instancjację jawną.
\end{enumerate}

\section{Biblioteki statyczne}
\label{sec:cpp17-static-libraries}

\emph{Biblioteka statyczna} jest typowo archiwum plików obiektowych.
W~Linux ma zwykle rozszerzenie \code{.a}, w~MSVC \code{.lib}.
Biblioteka sama nie jest osobnym procesem ani automatycznie tworzonym
wątkiem. Linker wybiera z~niej obiekty potrzebne do rozwiązania
zależności i~tworzy wynikowy program lub inną bibliotekę.

\begin{lstlisting}[language=bash,caption={Budowanie i uzycie biblioteki statycznej (Linux/GCC).}]
g++ -std=c++20 -O2 -c math.cpp -o math.o
ar rcs libmathstuff.a math.o

g++ -std=c++20 -O2 -c main.cpp -o main.o
g++ main.o -L. -lmathstuff -o app
./app
\end{lstlisting}

W~typowym poleceniu GNU linkera biblioteki archiwalne warto umieszczać
\emph{po} obiektach, które odwołują się do ich symboli. Przykładowo
\code{g++ main.o -lmathstuff} jest zazwyczaj właściwszym porządkiem
niż \code{g++ -lmathstuff main.o}. Linkowanie zależności cyklicznych
bibliotek statycznych może wymagać ponownego wskazania biblioteki
lub użycia grupy linkera (np.~\code{--start-group}/\code{--end-group}
w~GNU ld); lepiej jednak unikać takich zależności architektonicznie.

\nota{Statyczna biblioteka nie oznacza automatycznie w~pełni
statycznego programu. Aplikacja może dalej korzystać z~bibliotek
współdzielonych systemu i~standardowej biblioteki C++.}

\section{Biblioteki współdzielone i~dynamiczne ładowanie}
\label{sec:cpp17-shared-libraries}

\subsection{Budowanie biblioteki współdzielonej}

Na systemach ELF biblioteki współdzielone są zazwyczaj plikami
\code{.so}. Obiekty wykorzystywane w~takich bibliotekach zwykle
buduje się jako kod niezależny od położenia w~pamięci (\emph{PIC}).

\begin{lstlisting}[language=bash,caption={Biblioteka .so oraz program klienta (Linux/GCC).}]
g++ -std=c++20 -O2 -fPIC -c math.cpp -o math_pic.o
g++ -shared math_pic.o -o libmathstuff.so

g++ -std=c++20 -O2 main.cpp -L. -lmathstuff \
    '-Wl,-rpath,$ORIGIN' -o app
./app
\end{lstlisting}

Wyrażenie \code{\$ORIGIN} odnosi się w~typowym loaderze ELF do
katalogu, w~którym znajduje się ładowany obiekt. Apostrofy w~poleceniu
powłoki zapobiegają wcześniejszej ekspansji \code{\$ORIGIN} przez
powłokę. W~projekcie produkcyjnym należy przemyśleć ścieżki
instalacyjne, pakowanie bibliotek i~bezpieczeństwo ładowania.

Biblioteki współdzielone pozwalają wielu procesom mapować ich kod
z~tego samego pliku oraz wymieniać implementację bez konieczności
pełnego linkowania programu w~czasie jego budowania. Nie gwarantuje
to jednak, że \emph{dowolna} nowa wersja zachowa zgodność ABI.

\subsection{Wyszukiwanie bibliotek: RPATH/RUNPATH, SONAME}

Na Linux/ELF linker i~loader korzystają z~mechanizmów takich jak
\emph{SONAME}, \emph{RPATH} i~\emph{RUNPATH}. Reguły wyszukiwania
zależą od systemu i~konfiguracji; środowisko wykonania może również
uwzględniać ścieżki systemowe lub zmienne środowiskowe.

\begin{lstlisting}[language=bash,caption={Diagnostyka zaleznosci dynamicznych (Linux).}]
readelf -d ./app | grep -E 'NEEDED|RPATH|RUNPATH'
ldd ./app                  # Tylko dla zaufanych plikow
readelf -Ws libmathstuff.so
\end{lstlisting}

Programów nieznanego pochodzenia nie należy lekkomyślnie przekazywać
do narzędzi, które mogą uruchamiać kod lub loader; do statycznej
inspekcji lepiej nadaje się \code{readelf}.

\subsection{Eksportowanie symboli i~kontrola API}

Domyślne eksportowanie zbyt wielu symboli może powodować kolizje
nazw i~powiększać powierzchnię ABI. W~ELF można budować bibliotekę
z~\code{-fvisibility=hidden}, a~publiczne symbole jawnie oznaczać:

\begin{lstlisting}[caption={Przenosne w granicach GCC/Clang/MSVC makro eksportu.}]
// Plik: api.hpp
#pragma once

#if defined(_WIN32)
  #if defined(MYLIB_BUILD_SHARED)
    #define MYLIB_API __declspec(dllexport)
  #elif defined(MYLIB_USE_SHARED)
    #define MYLIB_API __declspec(dllimport)
  #else
    #define MYLIB_API
  #endif
#elif defined(__GNUC__) || defined(__clang__)
  #define MYLIB_API __attribute__((visibility("default")))
#else
  #define MYLIB_API
#endif

extern "C" MYLIB_API int mylib_add(int, int);
\end{lstlisting}

W~MSVC zwykle stosuje się projekt DLL oraz odpowiednią bibliotekę
importową; na macOS biblioteki mają najczęściej rozszerzenie
\code{.dylib}. Flagi \code{-fPIC}, \code{-shared} i~\code{-Wl,...}
w~powyższej postaci dotyczą konkretnego toolchainu, nie ISO C++.

\subsection{\texttt{dlopen}, \texttt{dlsym} i~granica interfejsu wtyczki}

Bibliotekę można także załadować podczas działania procesu, bez
linkowania jej wprost z~aplikacją. Na POSIX często używa się
\code{dlopen}/\code{dlsym}, natomiast na Windows odpowiedników
\code{LoadLibrary}/\code{GetProcAddress}.

\begin{lstlisting}[caption={Prosta wtyczka C ABI: biblioteka i klient POSIX (dwa pliki).}]
// Plik: plugin.cpp
extern "C" int plugin_square(int x) {
    return x * x;
}

// Plik: loader.cpp
#include <dlfcn.h>
#include <iostream>
#include <stdexcept>

int main() {
    void* lib = dlopen("./libplugin.so", RTLD_NOW);
    if (!lib) throw std::runtime_error(dlerror());

    using Fn = int (*)(int);
    void* address = dlsym(lib, "plugin_square");
    if (!address) {
        dlclose(lib);
        throw std::runtime_error("Missing plugin_square");
    }

    // POSIX dopuszcza konwersje dlsym -> wskaznik do funkcji.
    Fn square = reinterpret_cast<Fn>(address);
    std::cout << square(7) << '\n'; // 49
    dlclose(lib);
}
\end{lstlisting}

\begin{lstlisting}[language=bash,caption={Budowanie przyk?adu wtyczki (Linux/ELF).}]
g++ -std=c++20 -fPIC -shared plugin.cpp -o libplugin.so
g++ -std=c++20 loader.cpp -ldl -o loader
./loader
\end{lstlisting}

\nota{To przykład dla POSIX, a~nie w~pełni przenośna konstrukcja
ISO C++. Również rzutowanie wyniku \code{dlsym} na wskaźnik funkcji
opiera się na wymaganiach POSIX, a~nie na ogólnej gwarancji C++.
Dodatkowo nie wolno wyładowywać biblioteki, dopóki jakakolwiek funkcja
lub obiekt z~tej biblioteki nadal może być użyty. Przy bardziej
złożonych wtyczkach potrzebny jest kontrakt wersji ABI i~zarządzania
czasem życia.}

\section{ABI --- Application Binary Interface}
\label{sec:cpp17-abi}

\subsection{API a~ABI: dwa różne kontrakty}

\emph{API} opisuje przede wszystkim dostępne nazwy, typy, funkcje
i~semantykę na poziomie źródła. \emph{ABI} określa, jak te elementy
reprezentowane są w~kodzie maszynowym. ABI obejmuje między innymi:
\begin{itemize}
\item konwencje wywołania: przekazywanie parametrów, wartość zwracana,
      rejestry i~ułożenie stosu;
\item reprezentację i~wyrównanie obiektów, w~tym pól klas;
\item nazwy symboli C++, RTTI, mechanizmy wyjątków i~wirtualnych wywołań;
\item reguły udostępniania i~ładowania symboli bibliotek;
\item ABI standardowej biblioteki oraz ustawienia kompilatora.
\end{itemize}

Dwie biblioteki mogą być zgodne źródłowo i~jednocześnie binarnie
niezgodne. W~szczególności identyczny nagłówek nie oznacza, że plik
\code{.so} zbudowany innym toolchainem lub z~innym zestawem opcji
na pewno współpracuje z~aplikacją.

\subsection{Zmiana układu obiektu łamie zgodność binarną}

Rozważmy klasę publiczną:
\begin{lstlisting}[caption={Zmiana prywatnych pol nadal moze zmienic ABI.}]
// Wersja biblioteki 1
class Record {
    int id_{};
public:
    int id() const;
};

// Wersja biblioteki 2
class Record {
    int id_{};
    double score_{}; // Nawet prywatne pole zmienia uklad.
public:
    int id() const;
};
\end{lstlisting}

Klient kompiluje np.~\code{sizeof(Record)} według nagłówka,
który widział podczas kompilacji. Gdy zamienimy jedynie bibliotekę
na wersję z~innym układem \code{Record}, może dojść do naruszenia
kontraktu pamięci, nawet jeśli \code{id()} nadal ma taki sam podpis.
Podobne ryzyko dotyczy dodawania wirtualnych funkcji i~zmian
hierarchii dziedziczenia.

\subsection{Wzorzec PImpl}

Aby zmniejszyć zależności kompilacyjne i~ukryć układ implementacji,
można przechowywać wskaźnik do niekompletnego typu. Dzięki temu
zmiana \emph{pól implementacyjnych} często nie zmienia rozmiaru
publicznego obiektu.

\begin{lstlisting}[caption={PImpl: plik widget.hpp.}]
#pragma once
#include <memory>

class Widget {
public:
    Widget();
    ~Widget();

    Widget(Widget&&) noexcept;
    Widget& operator=(Widget&&) noexcept;

    Widget(const Widget&) = delete;
    Widget& operator=(const Widget&) = delete;

    int value() const;
    void set_value(int value);

private:
    struct Impl;
    std::unique_ptr<Impl> impl_;
};
\end{lstlisting}

\begin{lstlisting}[caption={PImpl: plik widget.cpp.}]
#include "widget.hpp"
#include <utility>

struct Widget::Impl {
    int value = 0;
    // Tu moga pojawiac sie dalsze pola.
};

Widget::Widget() : impl_(std::make_unique<Impl>()) {}
Widget::~Widget() = default;
Widget::Widget(Widget&&) noexcept = default;
Widget& Widget::operator=(Widget&&) noexcept = default;

int Widget::value() const { return impl_->value; }
void Widget::set_value(int value) { impl_->value = value; }
\end{lstlisting}

Destruktor jest zdefiniowany w~\code{widget.cpp}, gdzie \code{Impl}
jest kompletny; to istotne przy \code{std::unique_ptr<Impl>}.
PImpl nie zapewnia automatycznie pełnej stabilności ABI: nadal
liczy się zgodność kompilatora, standardowej biblioteki, eksportu
symboli, konwencji wywołania i~pozostałych publicznych deklaracji.
Dodatkowy wskaźnik i~alokacja na stercie mogą zwiększać koszt dostępu.

\subsection{C ABI jako ostrożna granica interoperacyjności}

Dla wtyczek i~bibliotek wykorzystywanych z~różnych języków często
projektuje się wąski interfejs oparty na funkcjach \code{extern "C"}
oraz nieprzezroczystych uchwytach (\emph{opaque handles}). Zamiast
przekazywać \code{std::string} lub klasy z~funkcjami wirtualnymi,
przekazuje się proste typy i~wskaźniki.

\begin{lstlisting}[caption={Przyk?adowy szkic stabilniejszego interfejsu binarnego.}]
#ifdef __cplusplus
extern "C" {
#endif

struct engine_handle;

engine_handle* engine_create(void);
void engine_destroy(engine_handle* p);
int engine_eval(engine_handle* p, double x, double* result);

#ifdef __cplusplus
}
#endif
\end{lstlisting}

Podpis funkcji nadal wymaga uzgodnienia ABI danej platformy,
konwencji wywołania i~kontraktu błędów. Przekraczanie takiej granicy
wyjątkami C++, pamięcią zwalnianą przez niezgodny alokator lub obiektami
standardowej biblioteki jest niebezpieczne. Bezpieczniej zwracać
kody błędu, a~alokację i~zwalnianie łączyć w~pary po tej samej
stronie interfejsu.

\subsection{Wersje kompilatora i~biblioteki standardowej}

Istnieją różne rodziny ABI C++ (np. rodzina Itanium ABI stosowana na
wielu platformach uniksowych oraz ABI MSVC w~Windows). Nawet w~ramach
jednej rodziny znaczenie mają flagi kompilatora, ustawienia standardowej
biblioteki, opcje wyjątków, RTTI i~ABI konkretnych typów.
Przykładowo GCC/libstdc++ używa mechanizmu tzw. \emph{dual ABI}
dla niektórych typów standardowych (m.in. \code{std::string}); mieszanie
niezgodnych ustawień może skutkować błędami linkera lub działania.
Nie należy zakładać, że biblioteki zbudowane dowolnymi wersjami
GCC, Clanga i~MSVC są wymienne binarnie.

\section{System budowania: CMake i~projekt wieloplikowy}
\label{sec:cpp17-cmake}

\subsection{Dlaczego targety są ważniejsze niż globalne flagi?}

Ręczne wywoływanie kompilatora dobrze tłumaczy mechanikę, lecz dla
większych projektów potrzebujemy grafu zależności, wariantów
Debug/Release, testów i~kontroli opcji kompilatora. Współczesny
CMake modeluje budowanie przez \emph{targety}: bibliotekę, program,
interfejsową bibliotekę nagłówkową i~ich właściwości.

Przykładowa struktura:
\begin{lstlisting}[caption={Uk?ad plikow klasycznej biblioteki CMake.}]
project/
  CMakeLists.txt
  include/
    mathstuff/math.hpp
  src/
    math.cpp
    main.cpp
\end{lstlisting}

W~tej konkretnej strukturze pliki źródłowe mogą zaczynać się od
\code{\#include "mathstuff/math.hpp"}, a~\code{main.cpp} również
korzysta z~tej samej ścieżki nagłówka. To dzięki
\code{target_include_directories} kompilator znajduje katalog
\code{include} bez wpisywania ręcznych flag \code{-I}.

\begin{lstlisting}[language={},caption={CMakeLists.txt dla klasycznej biblioteki.}]
cmake_minimum_required(VERSION 3.20)
project(MathStuff LANGUAGES CXX)

add_library(mathstuff STATIC src/math.cpp)
target_compile_features(mathstuff PUBLIC cxx_std_20)
target_include_directories(mathstuff PUBLIC
    $<BUILD_INTERFACE:${CMAKE_CURRENT_SOURCE_DIR}/include>
    $<INSTALL_INTERFACE:include>
)

add_executable(app src/main.cpp)
target_link_libraries(app PRIVATE mathstuff)
\end{lstlisting}

\begin{lstlisting}[language=bash,caption={Konfigurowanie i budowanie projektu z CMake.}]
cmake -S . -B build -DCMAKE_BUILD_TYPE=Debug
cmake --build build -j
./build/app
\end{lstlisting}

\code{PRIVATE}, \code{PUBLIC} i~\code{INTERFACE} określają
\emph{propagację wymagań użycia} (usage requirements) między
targetami:
\begin{center}
\begin{tabularx}{\textwidth}{@{}lX@{}}
\toprule
Słowo & Znaczenie\\
\midrule
\code{PRIVATE} & Wymaganie tylko na potrzeby budowania danego targetu.\\
\code{INTERFACE} & Wymaganie wyłącznie dla jego konsumentów.\\
\code{PUBLIC} & Wymaganie dla targetu i~jego konsumentów.\\
\bottomrule
\end{tabularx}
\end{center}

Gdy publiczny nagłówek biblioteki wymaga pewnego katalogu
nagłówków, ten katalog powinien zwykle propagować się do klientów.
Gdy jednak biblioteka korzysta z~zależności jedynie wewnątrz
\code{.cpp}, często wystarcza \code{PRIVATE}. Właściwy zakres
zależy też od tego, czy target jest biblioteką statyczną lub
współdzieloną i~jakie zależności są konieczne przy ostatecznym
linkowaniu.

\subsection{Biblioteka nagłówkowa i~\texttt{INTERFACE}}

\begin{lstlisting}[language={},caption={Header-only target i wymagania dla klientow.}]
add_library(my_headers INTERFACE)
target_include_directories(my_headers INTERFACE
    ${CMAKE_CURRENT_SOURCE_DIR}/include
)
target_compile_features(my_headers INTERFACE cxx_std_20)

add_executable(tool src/tool.cpp)
target_link_libraries(tool PRIVATE my_headers)
\end{lstlisting}

Taki target nie tworzy archiwum \code{.a}. Reprezentuje zestaw
wymagań i~publiczne nagłówki używane przez konsumentów.

\subsection{Diagnostyka, sanitizatory, LTO i~build types}

\begin{lstlisting}[language=bash,caption={Przyk?adowe flagi diagnostyczne GCC/Clang.}]
g++ -std=c++20 -O0 -g -Wall -Wextra -Wpedantic \
    -fsanitize=address,undefined main.cpp -o app_debug

g++ -std=c++20 -O2 -flto main.cpp math.cpp -o app_lto
\end{lstlisting}

\emph{Link-Time Optimization} (LTO) może umożliwiać optymalizacje
pomiędzy TU, np. dodatkowe wstawianie funkcji do miejsc wywołania.
Jest to jednak kosztowna funkcja toolchainu; jej efekty zależą od
programu i~opcji. Profilowanie i~powtarzalne pomiary są ważniejsze
niż sam fakt włączenia \code{-O3} lub \code{-flto}.

\section{Moduły C++20: nowy mechanizm organizacji kodu}
\label{sec:cpp17-modules}

\subsection{Motywacja: problem tekstowego dołączania nagłówków}

W~klasycznym modelu ten sam popularny nagłówek może być parsowany
oddzielnie w~wielu TU. Makra i~kolejność \code{\#include} wpływają
na wynik preprocessingu. Moduły C++20 umożliwiają jawne eksportowanie
deklaracji i~importowanie ich bez tekstowego wklejania całych
interfejsów do każdego klienta.

Zaletami modułów mogą być izolacja makr, bardziej jednoznaczne
zależności i~krótsze czasy kompilacji. Nie ma jednak gwarancji,
że \emph{każdy} program po migracji przyspieszy, a~wdrożenie
wymaga wsparcia kompilatora i~systemu budowania.

\subsection{Pierwszy moduł: \texttt{export module} i~\texttt{import}}

Rozważmy dwa pliki:
\begin{lstlisting}[caption={Interfejs modu?u geometry: plik geometry.cppm.}]
export module geometry;

export int square(int x) {
    return x * x;
}

int hidden_helper(int x) {
    return x + 1;
}
\end{lstlisting}

\begin{lstlisting}[caption={Klient modu?u: plik main.cpp.}]
#include <iostream>
import geometry;

int main() {
    std::cout << square(7) << '\n'; // 49
    // hidden_helper(7); // Nie jest eksportowana.
}
\end{lstlisting}

\code{export module geometry;} deklaruje \emph{główną jednostkę
interfejsu modułu} (\emph{primary module interface unit}).
\code{export int square(...)} udostępnia deklarację klientowi.
Samo \code{import geometry;} nie oznacza linkowania kodu funkcji
z~powietrza: jeśli moduł generuje kod obiektowy, trzeba go również
uwzględnić podczas linkowania.

\subsection{BMI/CMI: moduł nie jest biblioteką binarną}

Kompilator wytwarza dla interfejsu modułu własną reprezentację,
często nazywaną \emph{BMI} (\emph{Binary Module Interface}) lub
\emph{CMI} (\emph{Compiled Module Interface}). W~Clangu spotyka się
pliki \code{.pcm}, w~GCC \code{.gcm}, a~w~MSVC pliki \code{.ifc}.
Nie są to uniwersalne, przenośne między kompilatorami artefakty
podobne do biblioteki \code{.a}. Importer musi znać interfejs
przetworzony przez zgodną konfigurację kompilatora.

Zależności muszą być budowane w~poprawnym porządku:
\[
\text{źródło interfejsu modułu}
\longrightarrow \text{BMI/CMI}
\longrightarrow \text{kompilacja importera}
\longrightarrow \text{linkowanie obiektów}.
\]

\subsection{Rzeczywiste komendy: Clang i~GCC}

Dla prostego modułu \code{geometry.cppm} w~Clangu obsługującym
standardowe moduły można zastosować:
\begin{lstlisting}[language=bash,caption={Budowanie przyk?adowego modu?u w Clangu (tryb BMI).}]
clang++ -std=c++20 geometry.cppm --precompile -o geometry.pcm
clang++ -std=c++20 -c geometry.pcm -o geometry.o
clang++ -std=c++20 -fmodule-file=geometry=geometry.pcm \
    -c main.cpp -o main.o
clang++ geometry.o main.o -o app
./app
\end{lstlisting}

Dla GCC~14 (wsparcie modułów C++20 było w~tej wersji nadal rozwijane)
trzeba dodatkowo włączyć eksperymentalny mechanizm
\code{-fmodules-ts}. W~przykładach z~GCC~14 przydatne bywa jawne
wskazanie języka przy rozszerzeniu \code{.cppm}:
\begin{lstlisting}[language=bash,caption={Budowanie tego samego modu?u z GCC 14.}]
g++ -std=c++20 -fmodules-ts -x c++ -c geometry.cppm \
    -o geometry.o
# GCC tworzy potrzebne dane interfejsu w gcm.cache/.
g++ -std=c++20 -fmodules-ts -c main.cpp -o main.o
g++ geometry.o main.o -o app
./app
\end{lstlisting}

Dokładne opcje i~formaty modułów mogą zmieniać się między wersjami
narzędzi. W~dużym projekcie lepiej powierzyć zależności systemowi
budowania obsługującemu skanowanie importów, zamiast ręcznie
zarządzać plikami BMI.

\subsection{Jednostka interfejsu i~jednostka implementacji}

Moduł może rozdzielać interfejs od implementacji bez tradycyjnego
nagłówka. Jednostka implementacji deklaruje \code{module name;},
a~nie \code{export module name;}:

\begin{lstlisting}[caption={Dwa pliki jednego modu?u: interface.cppm i implementation.cpp.}]
// Plik: engine.cppm
export module engine;
export int run(int value);

// Plik: engine.cpp
module engine;
int run(int value) { return 2 * value; }
\end{lstlisting}

Importer widzi tylko deklaracje eksportowane przez interfejs.
Jednostkę implementacji również trzeba skompilować i~zlinkować.
Konkretny system budowania musi uwzględnić fakt, że jednostka
implementacji importuje interfejs swojego modułu niejawnie.

\subsection{Global module fragment i~nagłówki klasyczne}

Tradycyjny nagłówek może wymagać makr konfigurujących jego treść.
W~takim przypadku można skorzystać z~\emph{global module fragment}:

\begin{lstlisting}[caption={Global module fragment: tradycyjny include przed modu?em.}]
module;
#include <string>
#include <utility>

export module greeting;

export std::string greet(std::string name) {
    return "Hello " + std::move(name);
}
\end{lstlisting}

Dyrektywa \code{module;} rozpoczyna fragment globalny, a~następna
deklaracja \code{export module greeting;} rozpoczyna właściwą
jednostkę modułu. Importowanie nagłówków jednostkowych przez
\code{import <string>;} jest innym mechanizmem, tzw. \emph{header units},
których dostępność zależy od toolchainu. Nie utożsamiajmy go z~nazwanym
modułem standardowej biblioteki \code{import std;}.

\subsection{Partycje i~ponowny eksport importu}

Duży moduł można podzielić na partycje. Nazwa po dwukropku oznacza
partycję wewnątrz jednego modułu, a~nie odrębny nazwany moduł.

\begin{lstlisting}[caption={Dwie jednostki interfejsu modu?u z partycja.}]
// Plik: math_detail.cppm
export module mathlib:detail;
export int triple(int x) { return 3 * x; }

// Plik: mathlib.cppm
export module mathlib;
export import :detail;
export int times_six(int x) { return 2 * triple(x); }

// Plik: main.cpp
import mathlib;
int main() { return times_six(7) == 42 ? 0 : 1; }
\end{lstlisting}

Klient importuje \code{mathlib}, a~nie \code{mathlib:detail}.
W~tym przykładzie \code{export import :detail;} re-eksportuje
eksporty partycji. Wymaga to prawidłowej kolejności budowania
\emph{najpierw partycji, potem głównego interfejsu, a~na końcu klienta}.

\subsection{Łączność modułowa i~ODR}

C++20 wprowadza reguły własności deklaracji (\emph{module ownership}).
Deklaracje przypisane do nazwanego modułu należą do niego, a~pewne
nieeksportowane nazwy uzyskują \emph{module linkage}. Różni się to od
mechanicznego zastąpienia \code{\#include} przez \code{import}.

Przykładowo dodanie \code{export} nie sprawia automatycznie, że
interfejs C++ będzie zgodny ABI między kompilatorami. Ponadto
źródłowy kod szablonów nadal musi być semantycznie dostępny
odpowiedniemu mechanizmowi instancjacji, choć nie musi być już
tekstowo wklejany przez \code{\#include}.

\subsection{CMake i~\texttt{FILE\_SET CXX\_MODULES}}

Od CMake~3.28 dostępne są mechanizmy skanowania modułów i~specjalny
typ zbioru plików \code{CXX_MODULES}. Przykład projektu:

\begin{lstlisting}[language={},caption={CMakeLists.txt projektu korzystajacego z modu?ow C++20.}]
cmake_minimum_required(VERSION 3.28)
project(GeometryModules LANGUAGES CXX)

add_library(geometry)
target_compile_features(geometry PUBLIC cxx_std_20)
target_sources(geometry PUBLIC
    FILE_SET CXX_MODULES FILES geometry.cppm
)

add_executable(app main.cpp)
target_link_libraries(app PRIVATE geometry)
\end{lstlisting}

\begin{lstlisting}[language=bash,caption={Konfiguracja modu?ow przy uzyciu generatora Ninja.}]
cmake -S . -B build -G Ninja
cmake --build build
./build/app
\end{lstlisting}

Ten przykład pokazuje \emph{składnię} CMake, ale jego zbudowanie
wymaga zgodnego zestawu wersji CMake, kompilatora, generatora i~skanera
zależności. W~szczególności kompilator akceptujący składnię modułów
niekoniecznie współpracuje z~automatycznym skanowaniem w~danej
konfiguracji CMake. Na współczesnych systemach używa się m.in.
Ninja (co najmniej 1.11 w~typowych konfiguracjach) i~odpowiednich
wersji Clanga lub MSVC. Obsługa modułów standardowej biblioteki oraz
\emph{header units} ma odrębne ograniczenia.

\subsection{\texttt{import std;} nie jest tym samym co \texttt{\#include}}

Sam standard C++20 wprowadza językowy mechanizm modułów, ale
\emph{nie zobowiązuje} każdej biblioteki standardowej do dostarczenia
nazwanego modułu \code{std} w~tym standardzie. \code{import std;}
jest standaryzowane w~C++23; obsługa zależy od kompilatora, dostawcy
biblioteki standardowej i~wersji systemu budowania.

\begin{lstlisting}[caption={Nazwany modu? standardowej biblioteki (C++23, zalezny od wsparcia).}]
import std;

int main() {
    std::vector<int> xs{1, 2, 3};
    std::println("size = {}", xs.size());
}
\end{lstlisting}

Powyższy przykład wykorzystuje także \code{std::println} z~C++23.
Nie należy automatycznie kompilować go z~\code{-std=c++20} ani
zakładać, że \code{import std;} działa w~każdym środowisku C++23.

\section{Optymalizacja procesu budowania}
\label{sec:cpp17-build-perf}

\subsection{Incremental build i~zależności nagłówkowe}

Zmiana publicznego nagłówka może wymusić ponowną kompilację wielu TU,
podczas gdy zmiana prywatnego \code{.cpp} zwykle dotyka mniejszej
części projektu. Dlatego ograniczanie niepotrzebnych dołączeń
nagłówków, stosowanie deklaracji wyprzedzających tam, gdzie są
prawidłowe, i~PImpl mogą poprawiać czas budowania.

Moduły zmieniają strukturę zależności, ale nie usuwają konieczności
ich śledzenia; importer musi otrzymać aktualną reprezentację
interfejsu modułu.

\subsection{Precompiled headers, unity builds, ccache}

Popularne techniki przyspieszania kompilacji to:
\begin{description}[style=nextline]
\item[Precompiled headers (PCH)]
Przechowywanie wyniku parsowania stabilnych, obszernych nagłówków.
Pomagają w~projektach z~dużą liczbą TU, ale wymagają spójności
opcji kompilacji i~zarządzania unieważnianiem pamięci podręcznej.
\item[Unity builds]
Łączenie kilku plików źródłowych we wspólną TU. Może ograniczać
narzut kompilacji, ale też ujawniać kolizje anonimowych przestrzeni
nazw, makr i~definicji o~łączności wewnętrznej.
\item[Compilation cache]
Narzędzia takie jak \code{ccache} unikają ponownego wykonywania
kompilacji, jeżeli odpowiednio zdefiniowane wejścia pozostają takie same.
\item[Parallel build]
Kompilowanie niezależnych TU równolegle (np. \code{cmake --build build -j})
wykorzystuje wiele rdzeni, ale zwiększa chwilowe użycie pamięci.
\end{description}

\subsection{Powtarzalność buildów i~zależności transytywne}

Dla bibliotek publicznych warto utrzymywać przejrzystą strukturę
zależności. Przykładowo użytkownik nie powinien polegać na tym,
że \code{\#include <vector>} ,,przypadkiem'' pojawiło się wskutek
dołączenia innego nagłówka. Każdy nagłówek powinien zawierać
potrzebne deklaracje w~sposób niezależny od kolejności dołączeń.

Warto również zapisywać wersje kompilatora i~bibliotek, flagi
wpływające na ABI, wariant budowania oraz sposób instalowania
artefaktów. \emph{Reproducible builds} wymagają dodatkowo kontroli
źródeł niedeterministyczności, np. znaczników czasu i~ścieżek
zapisywanych w~plikach wynikowych.

\section{Praktyczny przewodnik diagnostyczny}
\label{sec:cpp17-troubleshooting}

\begin{center}
\small
\begin{tabularx}{\textwidth}{@{}p{0.34\textwidth}X@{}}
\toprule
Objaw & Najbardziej prawdopodobne źródło problemu\\
\midrule
\code{undefined reference} & Brakujący plik obiektowy, biblioteka,
definicja albo instancjacja szablonu.\\
\code{multiple definition} & Niepoprawna definicja zwykłej
funkcji lub zmiennej w~nagłówku albo podwójnie linkowany obiekt.\\
\code{undefined reference to vtable} & Niezdefiniowana wymagana
funkcja wirtualna (często destruktor) lub brak pliku implementacji.\\
\code{cannot find -lfoo} & Brak biblioteki lub nieprawidłowe
katalogi przeszukiwania linkera.\\
\code{cannot open shared object file} & Problem ścieżki lub
zależności ładowania, niekoniecznie problem kompilacji.\\
Kompilacja klienta działa, program kończy się błędem & Możliwa
niezgodność ABI, naruszenie ODR, zły czas życia obiektu lub błąd programu.\\
Importer nie znajduje modułu & Nie zbudowano BMI/CMI, nie wskazano
ścieżki albo konfiguracja skanera nie obsługuje importu.\\
Import działa, ale linker zgłasza brak funkcji & Brak pliku
obiektowego modułu/implementacji lub potrzebnej biblioteki.\\
Zmiana nagłówka przebudowuje prawie wszystko & Silne zależności
nagłówkowe; rozważ redukcję dołączeń, PImpl, moduły lub PCH.\\
\bottomrule
\end{tabularx}
\end{center}

\subsection{Kiedy wybrać nagłówki, bibliotekę czy moduł?}

\begin{itemize}
\item \textbf{Header-only} sprawdza się dla małych narzędzi, szablonów
      i~funkcji \code{inline}; ceną mogą być zależności kompilacyjne.
\item \textbf{Biblioteka statyczna} upraszcza dystrybucję kodu obiektowego
      wewnątrz aplikacji, ale nie usuwa wymagań ABI w~czasie linkowania.
\item \textbf{Biblioteka współdzielona} pozwala współdzielić kod
      i~wymieniać implementację, lecz wymaga jasnej polityki ABI,
      widoczności symboli i~wersjonowania.
\item \textbf{Moduł C++20} służy do semantycznego podziału interfejsu
      źródłowego i~kontroli importów; nie jest alternatywą wobec
      procesu linkowania ani automatycznym formatem dystrybucji ABI.
\item \textbf{PImpl i~C API} rozważaj na granicach bibliotek o~długim
      cyklu życia, szczególnie gdy stabilność ABI ma duże znaczenie.
\end{itemize}

\section*{Podsumowanie rozdziału}
\addcontentsline{toc}{section}{Podsumowanie rozdziału}

Kompletny program C++ jest wynikiem wspólnej pracy preprocessora,
kompilatora, linkera i~loadera. Organizując projekt, trzeba osobno
rozumieć \emph{deklaracje i~definicje}, \emph{ODR i~łączność},
\emph{symbole i~biblioteki} oraz \emph{ABI}. Nagłówki i~moduły
pomagają udostępniać interfejsy źródłowe; biblioteki statyczne
i~współdzielone dostarczają kod do połączenia lub załadowania.

Dobra organizacja projektu polega na jasnym rozdziale API i~implementacji,
ograniczaniu zależności, korzystaniu z~targetów CMake oraz świadomym
wersjonowaniu interfejsu binarnego. Nowoczesne moduły upraszczają część
problemów nagłówków, ale wprowadzają wymagania dotyczące narzędzi,
BMI/CMI i~kolejności kompilacji.

\section*{Literatura i~dokumentacja}
\addcontentsline{toc}{section}{Literatura i dokumentacja}
\begin{itemize}
\item ISO/IEC 14882:2020, zagadnienia \emph{translation units},
      \emph{linkage}, \emph{One Definition Rule} i~\emph{modules}.
\item \url{https://en.cppreference.com/w/cpp/language/definition}
      --- ODR, deklaracje i~definicje.
\item \url{https://en.cppreference.com/w/cpp/language/modules}
      --- składnia i~semantyka modułów.
\item \url{https://gcc.gnu.org/onlinedocs/gcc/C_002b_002b-Modules.html}
      --- implementacja modułów w~GCC.
\item \url{https://clang.llvm.org/docs/StandardCPlusPlusModules.html}
      --- moduły w~Clangu, BMI i~komendy budowania.
\item \url{https://cmake.org/cmake/help/latest/manual/cmake-cxxmodules.7.html}
      --- CMake i~zależności modułów.
\item \url{https://cmake.org/cmake/help/latest/command/target_sources.html}
      --- \code{FILE_SET CXX_MODULES}.
\item \url{https://gcc.gnu.org/onlinedocs/gcc/Code-Gen-Options.html}
      --- widoczność symboli w~ELF.
\item John Lakos, \emph{Large-Scale C++ Software Design}, Addison-Wesley, 1996.
\end{itemize}

% Komendy powloki sa przeznaczone glownie dla GCC/Clang na Linux/ELF;
% roznice Windows/MSVC i macOS zaznaczono w tekscie.

\chapter{Kompilacja, linkowanie, ABI i moduły C++20}
\label{chap:cpp17-build-abi-modules}

\section{Dlaczego trzeba rozumieć proces budowania?}

W~poprzednich rozdziałach projektowaliśmy klasy, szablony, mechanizmy
zarządzania pamięcią i~współbieżność. Dopóki cały program mieści się
w~jednym pliku, wystarcza zwykle wywołanie \code{g++ main.cpp -o app}.
W~realnym projekcie kod dzieli się jednak na dziesiątki lub tysiące
\emph{jednostek translacji}, biblioteki zewnętrzne i~moduły. To wtedy
pojawiają się komunikaty \code{undefined reference}, \code{multiple definition},
konflikty wersji bibliotek, długi czas budowania oraz błędy ujawniające
się dopiero po wymianie biblioteki współdzielonej.

Musimy rozdzielić trzy różne zagadnienia:
\begin{description}[style=nextline]
\item[Poprawność programu w~języku C++]
Czy deklaracje, definicje, typy i~wywołania spełniają reguły C++,
w~szczególności \emph{One Definition Rule} (ODR)?
\item[Model budowania i~linkowania]
Jakie pliki są kompilowane osobno, gdzie kompilator zapisuje symbole,
w~jakiej kolejności linker przeszukuje biblioteki oraz jak ładowane są
biblioteki dynamiczne?
\item[Zgodność binarna (ABI)]
Czy dwa osobno zbudowane komponenty zgadzają się co do reprezentacji
typów, sposobu wywoływania funkcji, nazw symboli i~pozostałych reguł
interfejsu binarnego?
\end{description}

\nota{Norma ISO C++ opisuje semantykę programu, ale nie narzuca jednego
formatu pliku obiektowego, konwencji \emph{name mangling}, organizacji
bibliotek \code{.so}/\code{.dll} ani poleceń linkera. Poniżej
odróżniamy reguły języka od typowych rozwiązań narzędziowych.}

\section{Od pliku źródłowego do programu}

\begin{center}\small
\begin{tabularx}{\textwidth}{@{}p{0.31\textwidth}X@{}}
\toprule
Etap & Za co odpowiada \\
\midrule
Preprocesor & Rozwija dyrektywy \code{#include}, makra i warunkową kompilację. \\
Kompilator & Sprawdza składnię i typy oraz wytwarza kod dla jednostki translacji. \\
Asembler & Zamienia kod asemblera na plik obiektowy z symbolami i relokacjami. \\
Linker & Łączy pliki obiektowe i biblioteki, rozwiązując odwołania do symboli. \\
Loader systemowy & Ładuje program i wymagane biblioteki współdzielone przed uruchomieniem. \\
\bottomrule
\end{tabularx}
\end{center}
\label{sec:cpp17-pipeline}

\subsection{Preprocesor, kompilator, asembler, linker i~loader}

W~uproszczeniu klasyczny projekt z~nagłówkami przechodzi przez etapy:
\[
\begin{aligned}
\text{źródła i nagłówki}&\ \longrightarrow\ \text{preprocessing}
\ \longrightarrow\ \text{kompilacja}\\
&\phantom{\ \longrightarrow\ }\longrightarrow\ \text{plik obiektowy}
\ \longrightarrow\ \text{linkowanie}
\ \longrightarrow\ \text{program}.
\end{aligned}
\]

Na niektórych platformach kompilator wytwarza tekst asemblerowy
(\code{.s}) jako etap pośredni; w~innych przejście do kodu maszynowego
jest zintegrowane. Linkowanie może odbywać się przy budowaniu programu,
a~dodatkowe rozwiązywanie symboli może następować w~czasie ładowania
lub później.

Prześledźmy etapy na samodzielnym programie.

\begin{lstlisting}[caption={Plik hello.cpp: prosty program do obserwowania kompilacji.}]
#include <iostream>

int twice(int x) { return 2 * x; }

int main() {
    std::cout << twice(21) << '\n';
}
\end{lstlisting}

\begin{lstlisting}[language=bash,caption={Wybrane etapy budowania przy uzyciu GCC/Clang.}]
g++ -std=c++20 -E hello.cpp -o hello.ii  # Preprocesor
# Produkt poprzedniego etapu moze byc duzy.
g++ -std=c++20 -S hello.cpp -o hello.s   # Asembler

g++ -std=c++20 -c hello.cpp -o hello.o   # Kod obiektowy
g++ hello.o -o hello                      # Linkowanie
./hello                                   # 42
\end{lstlisting}

Opcja \code{-c} nie tworzy wykonywalnego programu; produkuje plik
obiektowy, który może zawierać nierozwiązane odwołania do symboli
pochodzących z~innych jednostek translacji. Polecenie \code{-S}
kończy budowanie na asemblerze, a~\code{-E} na wstępnym przetwarzaniu.
Warto odróżnić program sterujący kompilacją (np.~\code{g++}) od samego
frontendu C++, asemblera i~linkera.

\subsection{Pliki obiektowe, symbole i~relokacje}

Plik obiektowy zawiera zwykle kod maszynowy, dane, tablice symboli
oraz \emph{relokacje}, czyli informacje, które adresy muszą zostać
skorygowane po ustaleniu ostatecznego układu programu. W~środowisku
Linux powszechny jest format ELF; Windows używa PE/COFF, a~macOS
Mach-O. Szczegóły tych formatów nie są częścią języka C++.

\begin{lstlisting}[language=bash,caption={Ogladanie symboli i relokacji (Linux/ELF).}]
nm -C hello.o                 # Symbole; -C demangluje nazwy C++
readelf -Ws hello.o          # Tablica symboli ELF
readelf -r hello.o           # Relokacje
objdump -drC hello.o         # Asembler i relokacje
\end{lstlisting}

Przykładowo litera \code{T} w~typowym wyjściu \code{nm} informuje
o~symbolu definiowanym w~sekcji kodu, a~\code{U} o~symbolu
niezdefiniowanym w~danym obiekcie, który trzeba rozwiązać. Nie
należy traktować oznaczeń \code{nm} jako przenośnej własności C++.

\subsection{Jednostka translacji i~kompilacja oddzielna}

W~tradycyjnym modelu \emph{jednostka translacji} (\emph{translation unit},
TU) powstaje z~pliku źródłowego po uwzględnieniu dyrektyw preprocesora,
w~szczególności zawartości dołączanych nagłówków. Każda TU jest
analizowana niezależnie, a~linker łączy otrzymane obiekty.

Rozważmy projekt:
\begin{lstlisting}[caption={Trzy pliki klasycznego projektu C++: API, implementacja i klient.}]
// Plik: math.hpp
#ifndef CPP17_MATH_HPP
#define CPP17_MATH_HPP
int square(int x);
#endif

// Plik: math.cpp
#include "math.hpp"
int square(int x) { return x * x; }

// Plik: main.cpp
#include "math.hpp"
#include <iostream>
int main() { std::cout << square(7) << '\n'; }
\end{lstlisting}

\begin{lstlisting}[language=bash,caption={Niezalezne budowanie TU i koncowe linkowanie.}]
g++ -std=c++20 -Wall -Wextra -c math.cpp -o math.o
g++ -std=c++20 -Wall -Wextra -c main.cpp -o main.o
g++ main.o math.o -o app
./app                        # 49
\end{lstlisting}

\code{math.hpp} dostarcza deklarację klientowi i~implementacji,
ale nie jest zwykle oddzielnym obiektem linkowanym. Zmiana samego
\code{math.cpp} nie musi wymuszać ponownej kompilacji \code{main.cpp},
o~ile nie zmienił się kontrakt w~nagłówku lub pozostałe zależności.

\subsection{Przykładowe komendy w~MSVC}

W~\emph{Developer Command Prompt for Visual Studio} można zbudować
analogiczny przykład za pomocą \code{cl.exe} i~\code{link.exe}:

\begin{lstlisting}[language=bash,caption={Kompilacja dwóch plików w środowisku MSVC/Windows.}]
cl /std:c++20 /EHsc /c math.cpp
cl /std:c++20 /EHsc /c main.cpp
link main.obj math.obj /OUT:app.exe
app.exe
\end{lstlisting}

Przy budowaniu programów C++ w~Visual Studio warto również
kontrolować zgodność wersji MSVC, ustawienia biblioteki uruchomieniowej
(np.~\code{/MD} wobec \code{/MT}) i~konfigurację Debug/Release.
Takie flagi nie mają bezpośrednich odpowiedników narzuconych przez
standard języka; należą do konkretnego środowiska.

\section{Nagłówki, makra i~zasada jednej definicji}
\label{sec:cpp17-headers-odr}

\subsection{Include guards i~\texttt{\#pragma once}}

Zawartość \code{\#include} jest w~uproszczeniu wstawiana do źródła
przez preprocesor. Wielokrotne dołączenie tego samego nagłówka
w~jednej TU mogłoby więc prowadzić do powtórzenia definicji klasy.
Chronimy się przez \emph{include guard}:

\begin{lstlisting}[caption={Klasyczna os?ona nag?owka.}]
#ifndef LIB_WIDGET_HPP
#define LIB_WIDGET_HPP

struct Widget {
    int value{};
};

#endif // LIB_WIDGET_HPP
\end{lstlisting}

W~praktyce popularna jest też dyrektywa \code{\#pragma once},
obsługiwana przez najważniejsze kompilatory, lecz niestandaryzowana
przez ISO C++. Oba rozwiązania chronią \emph{jedną TU} przed
powtórnym dołączeniem pliku. Nie rozwiązują automatycznie problemu
\emph{wielu definicji tego samego nie-inline symbolu w~różnych TU}.

\subsection{ODR: deklaracja kontra definicja}

Deklaracja wprowadza nazwę i~informację o~typie, natomiast definicja
zwykle dostarcza byt danej funkcji lub obiektu. Deklaracji pewnej
funkcji może być wiele, ale \emph{nie-inline} funkcja z~zewnętrzną
łącznością, używana w~programie w~sposób wymagający definicji,
powinna mieć jedną definicję w~programie. Reguły ODR zawierają
wyjątki pozwalające na identyczne definicje w~różnych TU, np. dla
odpowiednich funkcji \code{inline}, klas i~szablonów.

\begin{lstlisting}[caption={Typowa pu?apka: definicja zwyk?ej funkcji w nag?owku.}]
// BAD: wrong.hpp, dolaczony do dwoch plikow .cpp
int twice(int x) { return 2 * x; }
// Dwie zewnetrzne definicje -> naruszenie ODR.
\end{lstlisting}

\begin{lstlisting}[caption={Definicja inline przeznaczona do nag?owka.}]
// Plik: twice.hpp
#pragma once
inline int twice(int x) { return 2 * x; }
\end{lstlisting}

\code{inline} nie nakazuje kompilatorowi wstawienia instrukcji funkcji
bezpośrednio w~miejscu wywołania. Jego znaczenie dla organizacji
programu polega przede wszystkim na dopuszczeniu definicji w~wielu TU
zgodnie z~ODR (z~odpowiednimi warunkami). Wszystkie takie definicje
muszą być zgodne i~spełniać reguły ODR; samo dopisanie \code{inline}
nie naprawia różniących się implementacji w~dwóch plikach.

\subsection{Zmienne \texttt{extern}, \texttt{inline}, \texttt{const} i~\texttt{static}}

\begin{lstlisting}[caption={Jedna definicja zmiennej globalnej, wiele deklaracji.}]
// Plik: settings.hpp
#pragma once
extern int request_limit;       // Tylko deklaracja

// Plik: settings.cpp
#include "settings.hpp"
int request_limit = 100;        // Definicja

// Plik: main.cpp
#include "settings.hpp"
int main() { return request_limit == 100 ? 0 : 1; }
\end{lstlisting}

Od C++17 w~nagłówkach można również definiować \emph{zmienne inline}:
\begin{lstlisting}[caption={Zmienna inline w nag?owku (C++17).}]
#pragma once
inline constexpr int default_timeout_ms = 250;
inline int global_counter = 0;
\end{lstlisting}

W~tym przypadku odpowiednie definicje we wszystkich TU mogą
wskazywać na ten sam obiekt o~zewnętrznej łączności. W~kontraście
zmienna \code{const} zdefiniowana w~przestrzeni nazw bez \code{extern}
ma w~typowym przypadku łączność wewnętrzną; oddzielne TU mogą mieć
wówczas odrębne egzemplarze. \code{static} na poziomie przestrzeni
nazw także zapewnia łączność wewnętrzną, lecz na współczesnym C++
zwykle preferujemy anonimową przestrzeń nazw.

\begin{lstlisting}[caption={Nazwy pomocnicze ograniczone do jednej TU.}]
// Plik: helper.cpp
namespace {
    int internal_counter = 0;
    int twice_local(int x) { return x * 2; }
}

int perform(int x) {
    ++internal_counter;
    return twice_local(x);
}
\end{lstlisting}

\subsection{Szablony: dlaczego definicje zwykle znajdują się w~nagłówkach?}

Kompilator instancjuje szablon dla określonych argumentów typowych
w~miejscu, w~którym wymagana jest jego definicja. Samo udostępnienie
deklaracji z~pliku \code{.hpp} i~ukrycie ogólnej definicji szablonu
w~oddzielnym \code{.cpp} może zatem prowadzić do błędu linkera.

\begin{lstlisting}[caption={Najprostszy, tradycyjny sposob udostepniania szablonu.}]
// Plik: multiply.hpp
#pragma once

template<class T>
T multiply(T a, T b) {
    return a * b;
}
\end{lstlisting}

Można także jawnie instancjować \emph{konkretne} typy w~pliku
implementacji. Oddzielne jednostki translacji nie muszą wtedy
instancjować tego samego kodu niezależnie:

\begin{lstlisting}[caption={Jawna instancjacja okreslonego wariantu szablonu.}]
// Plik: operation.hpp
#pragma once
template<class T>
T add(T a, T b);
extern template int add<int>(int, int);

// Plik: operation.cpp
#include "operation.hpp"
template<class T>
T add(T a, T b) { return a + b; }
template int add<int>(int, int);

// Plik: main.cpp
#include "operation.hpp"
int main() { return add(10, 32) == 42 ? 0 : 1; }
\end{lstlisting}

W~tym przykładzie klient może korzystać z~\code{add<int>}, ale
instancjacja innych typów nie zostaje mu udostępniona automatycznie.
Definicje szablonów umieszczane w~interfejsach modułów podlegają
nieco innemu modelowi widoczności, który omówimy dalej.

\subsection{Naruszenie ODR bez łatwego komunikatu}

Nie wszystkie naruszenia ODR skutkują natychmiastowym błędem
linkera. Szczególnie groźne jest warunkowe definiowanie tej samej
klasy w~dwóch TU w~inny sposób, np. przy użyciu sprzecznych
makr kompilacji:

\begin{lstlisting}[caption={Przyk?ad niezgodnych definicji typu (celowo b?edny).}]
// Plik: config.hpp
struct Config {
#ifdef EXTRA_FIELD
    int extra;
#endif
    int value;
};

// Jedna TU z -DEXTRA_FIELD, inna bez tej flagi:
// rozne definicje klasy -> naruszenie ODR.
\end{lstlisting}

Taki program może się nawet zlinkować, lecz pozostaje niepoprawny.
Dlatego opcje zmieniające publiczne definicje klas i~struktur muszą
być spójne we wszystkich TU danego programu.

\section{Symbole, łączność i~typowe błędy linkera}
\label{sec:cpp17-symbols}

\subsection{Łączność a~widoczność symbolu}

\emph{Linkage} w~języku C++ mówi, czy deklaracje w~różnych zakresach
i~TU mogą nazywać tę samą encję. W~praktyce spotykamy między innymi
łączność zewnętrzną, wewnętrzną, a~w~jednostkach modułów także
łączność modułową. Nie należy jednak utożsamiać tej własności
języka z~\emph{visibility} na poziomie ELF ani z~eksportem symbolu
z~biblioteki dynamicznej.

\begin{center}
\begin{tabularx}{\textwidth}{@{}lX@{}}
\toprule
Mechanizm & Co określa?\\
\midrule
\code{static} w~przestrzeni nazw & Łączność wewnętrzna nazwy w~TU.\\
\code{extern} & M.in. deklarowanie obiektu definiowanego w~innym miejscu.\\
Anonimowa przestrzeń nazw & Wewnętrzna łączność elementów pomocniczych.\\
\code{inline} & Warunki wielokrotnych definicji zgodnych z~ODR.\\
\code{-fvisibility=hidden} & Domyślna widoczność symboli w~typowych bibliotekach ELF.\\
\code{__declspec(dllexport)} & Udostępnianie symboli biblioteki DLL w~MSVC.\\
\bottomrule
\end{tabularx}
\end{center}

\subsection{Name mangling i~\texttt{extern~\textquotedbl{}C\textquotedbl{}}}

C++ dopuszcza przeciążanie funkcji, przestrzenie nazw oraz rozbudowane
typy szablonowe. Aby móc odróżniać symbole po linkowaniu,
implementacje zwykle kodują część tych informacji w~nazwach
binarnych (\emph{name mangling}). Dokładny sposób kodowania nie jest
ustalony przez ISO C++.

\begin{lstlisting}[caption={Przeciazanie i punkt wejscia C ABI.}]
namespace calc {
    int f(int x) { return x + 1; }
    double f(double x) { return x + 1.0; }
}

extern "C" int api_add(int a, int b) {
    return a + b;
}
\end{lstlisting}

Na typowej platformie ELF/Itanium ABI funkcje przeciążone otrzymują
różne symbole zakodowane przez kompilator, podczas gdy \code{api_add}
ma nazwę zgodną z~łącznością języka C. Nie jest to gwarancja pełnej
zgodności ABI z~dowolną biblioteką C: znaczenie mają także deklaracje
typów i~konwencja wywołania ustalona dla danej platformy.

\begin{lstlisting}[language=bash,caption={Demanglowanie i inspekcja symboli (przyk?ady narzedziowe).}]
g++ -std=c++20 -c overload.cpp -o overload.o
nm overload.o
nm -C overload.o
c++filt _Z3fooi      # Ilustracja popularnego schematu manglingu
\end{lstlisting}

\subsection{\texttt{undefined reference}, \texttt{multiple definition}}

Błąd \code{undefined reference} oznacza zwykle, że linker znalazł
odwołanie do symbolu, lecz nie znalazł odpowiedniej definicji.
Przykładowe przyczyny to pominięcie \code{math.o}, pominięcie
biblioteki, niezgodny podpis funkcji albo brak jawnej instancjacji
szablonu. Z~kolei \code{multiple definition} sugeruje m.in. zwykłą
funkcję lub zmienną zdefiniowaną w~nagłówku dołączanym do wielu TU.

Przydatna lista kontrolna:
\begin{enumerate}
\item Sprawdź, czy wszystkie wymagane \code{.cpp} są kompilowane i~linkowane.
\item Porównaj podpis funkcji: przestrzeń nazw, typy, kwalifikatory i~ABI.
\item Upewnij się, że deklaracja i~definicja pochodzą z~tego samego interfejsu.
\item Sprawdź kolejność bibliotek statycznych i~ich zależności.
\item Dla klas polimorficznych sprawdź brakujące definicje funkcji
      wirtualnych, w~tym destruktora wirtualnego, który może wymagać
      definicji nawet jako funkcja \code{= 0}.
\item Dla szablonów sprawdź dostępność definicji w~punkcie instancjacji
      lub odpowiednią instancjację jawną.
\end{enumerate}

\section{Biblioteki statyczne}
\label{sec:cpp17-static-libraries}

\emph{Biblioteka statyczna} jest typowo archiwum plików obiektowych.
W~Linux ma zwykle rozszerzenie \code{.a}, w~MSVC \code{.lib}.
Biblioteka sama nie jest osobnym procesem ani automatycznie tworzonym
wątkiem. Linker wybiera z~niej obiekty potrzebne do rozwiązania
zależności i~tworzy wynikowy program lub inną bibliotekę.

\begin{lstlisting}[language=bash,caption={Budowanie i uzycie biblioteki statycznej (Linux/GCC).}]
g++ -std=c++20 -O2 -c math.cpp -o math.o
ar rcs libmathstuff.a math.o

g++ -std=c++20 -O2 -c main.cpp -o main.o
g++ main.o -L. -lmathstuff -o app
./app
\end{lstlisting}

W~typowym poleceniu GNU linkera biblioteki archiwalne warto umieszczać
\emph{po} obiektach, które odwołują się do ich symboli. Przykładowo
\code{g++ main.o -lmathstuff} jest zazwyczaj właściwszym porządkiem
niż \code{g++ -lmathstuff main.o}. Linkowanie zależności cyklicznych
bibliotek statycznych może wymagać ponownego wskazania biblioteki
lub użycia grupy linkera (np.~\code{--start-group}/\code{--end-group}
w~GNU ld); lepiej jednak unikać takich zależności architektonicznie.

\nota{Statyczna biblioteka nie oznacza automatycznie w~pełni
statycznego programu. Aplikacja może dalej korzystać z~bibliotek
współdzielonych systemu i~standardowej biblioteki C++.}

\section{Biblioteki współdzielone i~dynamiczne ładowanie}
\label{sec:cpp17-shared-libraries}

\subsection{Budowanie biblioteki współdzielonej}

Na systemach ELF biblioteki współdzielone są zazwyczaj plikami
\code{.so}. Obiekty wykorzystywane w~takich bibliotekach zwykle
buduje się jako kod niezależny od położenia w~pamięci (\emph{PIC}).

\begin{lstlisting}[language=bash,caption={Biblioteka .so oraz program klienta (Linux/GCC).}]
g++ -std=c++20 -O2 -fPIC -c math.cpp -o math_pic.o
g++ -shared math_pic.o -o libmathstuff.so

g++ -std=c++20 -O2 main.cpp -L. -lmathstuff \
    '-Wl,-rpath,$ORIGIN' -o app
./app
\end{lstlisting}

Wyrażenie \code{\$ORIGIN} odnosi się w~typowym loaderze ELF do
katalogu, w~którym znajduje się ładowany obiekt. Apostrofy w~poleceniu
powłoki zapobiegają wcześniejszej ekspansji \code{\$ORIGIN} przez
powłokę. W~projekcie produkcyjnym należy przemyśleć ścieżki
instalacyjne, pakowanie bibliotek i~bezpieczeństwo ładowania.

Biblioteki współdzielone pozwalają wielu procesom mapować ich kod
z~tego samego pliku oraz wymieniać implementację bez konieczności
pełnego linkowania programu w~czasie jego budowania. Nie gwarantuje
to jednak, że \emph{dowolna} nowa wersja zachowa zgodność ABI.

\subsection{Wyszukiwanie bibliotek: RPATH/RUNPATH, SONAME}

Na Linux/ELF linker i~loader korzystają z~mechanizmów takich jak
\emph{SONAME}, \emph{RPATH} i~\emph{RUNPATH}. Reguły wyszukiwania
zależą od systemu i~konfiguracji; środowisko wykonania może również
uwzględniać ścieżki systemowe lub zmienne środowiskowe.

\begin{lstlisting}[language=bash,caption={Diagnostyka zaleznosci dynamicznych (Linux).}]
readelf -d ./app | grep -E 'NEEDED|RPATH|RUNPATH'
ldd ./app                  # Tylko dla zaufanych plikow
readelf -Ws libmathstuff.so
\end{lstlisting}

Programów nieznanego pochodzenia nie należy lekkomyślnie przekazywać
do narzędzi, które mogą uruchamiać kod lub loader; do statycznej
inspekcji lepiej nadaje się \code{readelf}.

\subsection{Eksportowanie symboli i~kontrola API}

Domyślne eksportowanie zbyt wielu symboli może powodować kolizje
nazw i~powiększać powierzchnię ABI. W~ELF można budować bibliotekę
z~\code{-fvisibility=hidden}, a~publiczne symbole jawnie oznaczać:

\begin{lstlisting}[caption={Przenosne w granicach GCC/Clang/MSVC makro eksportu.}]
// Plik: api.hpp
#pragma once

#if defined(_WIN32)
  #if defined(MYLIB_BUILD_SHARED)
    #define MYLIB_API __declspec(dllexport)
  #elif defined(MYLIB_USE_SHARED)
    #define MYLIB_API __declspec(dllimport)
  #else
    #define MYLIB_API
  #endif
#elif defined(__GNUC__) || defined(__clang__)
  #define MYLIB_API __attribute__((visibility("default")))
#else
  #define MYLIB_API
#endif

extern "C" MYLIB_API int mylib_add(int, int);
\end{lstlisting}

W~MSVC zwykle stosuje się projekt DLL oraz odpowiednią bibliotekę
importową; na macOS biblioteki mają najczęściej rozszerzenie
\code{.dylib}. Flagi \code{-fPIC}, \code{-shared} i~\code{-Wl,...}
w~powyższej postaci dotyczą konkretnego toolchainu, nie ISO C++.

\subsection{\texttt{dlopen}, \texttt{dlsym} i~granica interfejsu wtyczki}

Bibliotekę można także załadować podczas działania procesu, bez
linkowania jej wprost z~aplikacją. Na POSIX często używa się
\code{dlopen}/\code{dlsym}, natomiast na Windows odpowiedników
\code{LoadLibrary}/\code{GetProcAddress}.

\begin{lstlisting}[caption={Prosta wtyczka C ABI: biblioteka i klient POSIX (dwa pliki).}]
// Plik: plugin.cpp
extern "C" int plugin_square(int x) {
    return x * x;
}

// Plik: loader.cpp
#include <dlfcn.h>
#include <iostream>
#include <stdexcept>

int main() {
    void* lib = dlopen("./libplugin.so", RTLD_NOW);
    if (!lib) throw std::runtime_error(dlerror());

    using Fn = int (*)(int);
    void* address = dlsym(lib, "plugin_square");
    if (!address) {
        dlclose(lib);
        throw std::runtime_error("Missing plugin_square");
    }

    // POSIX dopuszcza konwersje dlsym -> wskaznik do funkcji.
    Fn square = reinterpret_cast<Fn>(address);
    std::cout << square(7) << '\n'; // 49
    dlclose(lib);
}
\end{lstlisting}

\begin{lstlisting}[language=bash,caption={Budowanie przyk?adu wtyczki (Linux/ELF).}]
g++ -std=c++20 -fPIC -shared plugin.cpp -o libplugin.so
g++ -std=c++20 loader.cpp -ldl -o loader
./loader
\end{lstlisting}

\nota{To przykład dla POSIX, a~nie w~pełni przenośna konstrukcja
ISO C++. Również rzutowanie wyniku \code{dlsym} na wskaźnik funkcji
opiera się na wymaganiach POSIX, a~nie na ogólnej gwarancji C++.
Dodatkowo nie wolno wyładowywać biblioteki, dopóki jakakolwiek funkcja
lub obiekt z~tej biblioteki nadal może być użyty. Przy bardziej
złożonych wtyczkach potrzebny jest kontrakt wersji ABI i~zarządzania
czasem życia.}

\section{ABI --- Application Binary Interface}
\label{sec:cpp17-abi}

\subsection{API a~ABI: dwa różne kontrakty}

\emph{API} opisuje przede wszystkim dostępne nazwy, typy, funkcje
i~semantykę na poziomie źródła. \emph{ABI} określa, jak te elementy
reprezentowane są w~kodzie maszynowym. ABI obejmuje między innymi:
\begin{itemize}
\item konwencje wywołania: przekazywanie parametrów, wartość zwracana,
      rejestry i~ułożenie stosu;
\item reprezentację i~wyrównanie obiektów, w~tym pól klas;
\item nazwy symboli C++, RTTI, mechanizmy wyjątków i~wirtualnych wywołań;
\item reguły udostępniania i~ładowania symboli bibliotek;
\item ABI standardowej biblioteki oraz ustawienia kompilatora.
\end{itemize}

Dwie biblioteki mogą być zgodne źródłowo i~jednocześnie binarnie
niezgodne. W~szczególności identyczny nagłówek nie oznacza, że plik
\code{.so} zbudowany innym toolchainem lub z~innym zestawem opcji
na pewno współpracuje z~aplikacją.

\subsection{Zmiana układu obiektu łamie zgodność binarną}

Rozważmy klasę publiczną:
\begin{lstlisting}[caption={Zmiana prywatnych pol nadal moze zmienic ABI.}]
// Wersja biblioteki 1
class Record {
    int id_{};
public:
    int id() const;
};

// Wersja biblioteki 2
class Record {
    int id_{};
    double score_{}; // Nawet prywatne pole zmienia uklad.
public:
    int id() const;
};
\end{lstlisting}

Klient kompiluje np.~\code{sizeof(Record)} według nagłówka,
który widział podczas kompilacji. Gdy zamienimy jedynie bibliotekę
na wersję z~innym układem \code{Record}, może dojść do naruszenia
kontraktu pamięci, nawet jeśli \code{id()} nadal ma taki sam podpis.
Podobne ryzyko dotyczy dodawania wirtualnych funkcji i~zmian
hierarchii dziedziczenia.

\subsection{Wzorzec PImpl}

Aby zmniejszyć zależności kompilacyjne i~ukryć układ implementacji,
można przechowywać wskaźnik do niekompletnego typu. Dzięki temu
zmiana \emph{pól implementacyjnych} często nie zmienia rozmiaru
publicznego obiektu.

\begin{lstlisting}[caption={PImpl: plik widget.hpp.}]
#pragma once
#include <memory>

class Widget {
public:
    Widget();
    ~Widget();

    Widget(Widget&&) noexcept;
    Widget& operator=(Widget&&) noexcept;

    Widget(const Widget&) = delete;
    Widget& operator=(const Widget&) = delete;

    int value() const;
    void set_value(int value);

private:
    struct Impl;
    std::unique_ptr<Impl> impl_;
};
\end{lstlisting}

\begin{lstlisting}[caption={PImpl: plik widget.cpp.}]
#include "widget.hpp"
#include <utility>

struct Widget::Impl {
    int value = 0;
    // Tu moga pojawiac sie dalsze pola.
};

Widget::Widget() : impl_(std::make_unique<Impl>()) {}
Widget::~Widget() = default;
Widget::Widget(Widget&&) noexcept = default;
Widget& Widget::operator=(Widget&&) noexcept = default;

int Widget::value() const { return impl_->value; }
void Widget::set_value(int value) { impl_->value = value; }
\end{lstlisting}

Destruktor jest zdefiniowany w~\code{widget.cpp}, gdzie \code{Impl}
jest kompletny; to istotne przy \code{std::unique_ptr<Impl>}.
PImpl nie zapewnia automatycznie pełnej stabilności ABI: nadal
liczy się zgodność kompilatora, standardowej biblioteki, eksportu
symboli, konwencji wywołania i~pozostałych publicznych deklaracji.
Dodatkowy wskaźnik i~alokacja na stercie mogą zwiększać koszt dostępu.

\subsection{C ABI jako ostrożna granica interoperacyjności}

Dla wtyczek i~bibliotek wykorzystywanych z~różnych języków często
projektuje się wąski interfejs oparty na funkcjach \code{extern "C"}
oraz nieprzezroczystych uchwytach (\emph{opaque handles}). Zamiast
przekazywać \code{std::string} lub klasy z~funkcjami wirtualnymi,
przekazuje się proste typy i~wskaźniki.

\begin{lstlisting}[caption={Przyk?adowy szkic stabilniejszego interfejsu binarnego.}]
#ifdef __cplusplus
extern "C" {
#endif

struct engine_handle;

engine_handle* engine_create(void);
void engine_destroy(engine_handle* p);
int engine_eval(engine_handle* p, double x, double* result);

#ifdef __cplusplus
}
#endif
\end{lstlisting}

Podpis funkcji nadal wymaga uzgodnienia ABI danej platformy,
konwencji wywołania i~kontraktu błędów. Przekraczanie takiej granicy
wyjątkami C++, pamięcią zwalnianą przez niezgodny alokator lub obiektami
standardowej biblioteki jest niebezpieczne. Bezpieczniej zwracać
kody błędu, a~alokację i~zwalnianie łączyć w~pary po tej samej
stronie interfejsu.

\subsection{Wersje kompilatora i~biblioteki standardowej}

Istnieją różne rodziny ABI C++ (np. rodzina Itanium ABI stosowana na
wielu platformach uniksowych oraz ABI MSVC w~Windows). Nawet w~ramach
jednej rodziny znaczenie mają flagi kompilatora, ustawienia standardowej
biblioteki, opcje wyjątków, RTTI i~ABI konkretnych typów.
Przykładowo GCC/libstdc++ używa mechanizmu tzw. \emph{dual ABI}
dla niektórych typów standardowych (m.in. \code{std::string}); mieszanie
niezgodnych ustawień może skutkować błędami linkera lub działania.
Nie należy zakładać, że biblioteki zbudowane dowolnymi wersjami
GCC, Clanga i~MSVC są wymienne binarnie.

\section{System budowania: CMake i~projekt wieloplikowy}
\label{sec:cpp17-cmake}

\subsection{Dlaczego targety są ważniejsze niż globalne flagi?}

Ręczne wywoływanie kompilatora dobrze tłumaczy mechanikę, lecz dla
większych projektów potrzebujemy grafu zależności, wariantów
Debug/Release, testów i~kontroli opcji kompilatora. Współczesny
CMake modeluje budowanie przez \emph{targety}: bibliotekę, program,
interfejsową bibliotekę nagłówkową i~ich właściwości.

Przykładowa struktura:
\begin{lstlisting}[caption={Uk?ad plikow klasycznej biblioteki CMake.}]
project/
  CMakeLists.txt
  include/
    mathstuff/math.hpp
  src/
    math.cpp
    main.cpp
\end{lstlisting}

W~tej konkretnej strukturze pliki źródłowe mogą zaczynać się od
\code{\#include "mathstuff/math.hpp"}, a~\code{main.cpp} również
korzysta z~tej samej ścieżki nagłówka. To dzięki
\code{target_include_directories} kompilator znajduje katalog
\code{include} bez wpisywania ręcznych flag \code{-I}.

\begin{lstlisting}[language={},caption={CMakeLists.txt dla klasycznej biblioteki.}]
cmake_minimum_required(VERSION 3.20)
project(MathStuff LANGUAGES CXX)

add_library(mathstuff STATIC src/math.cpp)
target_compile_features(mathstuff PUBLIC cxx_std_20)
target_include_directories(mathstuff PUBLIC
    $<BUILD_INTERFACE:${CMAKE_CURRENT_SOURCE_DIR}/include>
    $<INSTALL_INTERFACE:include>
)

add_executable(app src/main.cpp)
target_link_libraries(app PRIVATE mathstuff)
\end{lstlisting}

\begin{lstlisting}[language=bash,caption={Konfigurowanie i budowanie projektu z CMake.}]
cmake -S . -B build -DCMAKE_BUILD_TYPE=Debug
cmake --build build -j
./build/app
\end{lstlisting}

\code{PRIVATE}, \code{PUBLIC} i~\code{INTERFACE} określają
\emph{propagację wymagań użycia} (usage requirements) między
targetami:
\begin{center}
\begin{tabularx}{\textwidth}{@{}lX@{}}
\toprule
Słowo & Znaczenie\\
\midrule
\code{PRIVATE} & Wymaganie tylko na potrzeby budowania danego targetu.\\
\code{INTERFACE} & Wymaganie wyłącznie dla jego konsumentów.\\
\code{PUBLIC} & Wymaganie dla targetu i~jego konsumentów.\\
\bottomrule
\end{tabularx}
\end{center}

Gdy publiczny nagłówek biblioteki wymaga pewnego katalogu
nagłówków, ten katalog powinien zwykle propagować się do klientów.
Gdy jednak biblioteka korzysta z~zależności jedynie wewnątrz
\code{.cpp}, często wystarcza \code{PRIVATE}. Właściwy zakres
zależy też od tego, czy target jest biblioteką statyczną lub
współdzieloną i~jakie zależności są konieczne przy ostatecznym
linkowaniu.

\subsection{Biblioteka nagłówkowa i~\texttt{INTERFACE}}

\begin{lstlisting}[language={},caption={Header-only target i wymagania dla klientow.}]
add_library(my_headers INTERFACE)
target_include_directories(my_headers INTERFACE
    ${CMAKE_CURRENT_SOURCE_DIR}/include
)
target_compile_features(my_headers INTERFACE cxx_std_20)

add_executable(tool src/tool.cpp)
target_link_libraries(tool PRIVATE my_headers)
\end{lstlisting}

Taki target nie tworzy archiwum \code{.a}. Reprezentuje zestaw
wymagań i~publiczne nagłówki używane przez konsumentów.

\subsection{Diagnostyka, sanitizatory, LTO i~build types}

\begin{lstlisting}[language=bash,caption={Przyk?adowe flagi diagnostyczne GCC/Clang.}]
g++ -std=c++20 -O0 -g -Wall -Wextra -Wpedantic \
    -fsanitize=address,undefined main.cpp -o app_debug

g++ -std=c++20 -O2 -flto main.cpp math.cpp -o app_lto
\end{lstlisting}

\emph{Link-Time Optimization} (LTO) może umożliwiać optymalizacje
pomiędzy TU, np. dodatkowe wstawianie funkcji do miejsc wywołania.
Jest to jednak kosztowna funkcja toolchainu; jej efekty zależą od
programu i~opcji. Profilowanie i~powtarzalne pomiary są ważniejsze
niż sam fakt włączenia \code{-O3} lub \code{-flto}.

\section{Moduły C++20: nowy mechanizm organizacji kodu}
\label{sec:cpp17-modules}

\subsection{Motywacja: problem tekstowego dołączania nagłówków}

W~klasycznym modelu ten sam popularny nagłówek może być parsowany
oddzielnie w~wielu TU. Makra i~kolejność \code{\#include} wpływają
na wynik preprocessingu. Moduły C++20 umożliwiają jawne eksportowanie
deklaracji i~importowanie ich bez tekstowego wklejania całych
interfejsów do każdego klienta.

Zaletami modułów mogą być izolacja makr, bardziej jednoznaczne
zależności i~krótsze czasy kompilacji. Nie ma jednak gwarancji,
że \emph{każdy} program po migracji przyspieszy, a~wdrożenie
wymaga wsparcia kompilatora i~systemu budowania.

\subsection{Pierwszy moduł: \texttt{export module} i~\texttt{import}}

Rozważmy dwa pliki:
\begin{lstlisting}[caption={Interfejs modu?u geometry: plik geometry.cppm.}]
export module geometry;

export int square(int x) {
    return x * x;
}

int hidden_helper(int x) {
    return x + 1;
}
\end{lstlisting}

\begin{lstlisting}[caption={Klient modu?u: plik main.cpp.}]
#include <iostream>
import geometry;

int main() {
    std::cout << square(7) << '\n'; // 49
    // hidden_helper(7); // Nie jest eksportowana.
}
\end{lstlisting}

\code{export module geometry;} deklaruje \emph{główną jednostkę
interfejsu modułu} (\emph{primary module interface unit}).
\code{export int square(...)} udostępnia deklarację klientowi.
Samo \code{import geometry;} nie oznacza linkowania kodu funkcji
z~powietrza: jeśli moduł generuje kod obiektowy, trzeba go również
uwzględnić podczas linkowania.

\subsection{BMI/CMI: moduł nie jest biblioteką binarną}

Kompilator wytwarza dla interfejsu modułu własną reprezentację,
często nazywaną \emph{BMI} (\emph{Binary Module Interface}) lub
\emph{CMI} (\emph{Compiled Module Interface}). W~Clangu spotyka się
pliki \code{.pcm}, w~GCC \code{.gcm}, a~w~MSVC pliki \code{.ifc}.
Nie są to uniwersalne, przenośne między kompilatorami artefakty
podobne do biblioteki \code{.a}. Importer musi znać interfejs
przetworzony przez zgodną konfigurację kompilatora.

Zależności muszą być budowane w~poprawnym porządku:
\[
\text{źródło interfejsu modułu}
\longrightarrow \text{BMI/CMI}
\longrightarrow \text{kompilacja importera}
\longrightarrow \text{linkowanie obiektów}.
\]

\subsection{Rzeczywiste komendy: Clang i~GCC}

Dla prostego modułu \code{geometry.cppm} w~Clangu obsługującym
standardowe moduły można zastosować:
\begin{lstlisting}[language=bash,caption={Budowanie przyk?adowego modu?u w Clangu (tryb BMI).}]
clang++ -std=c++20 geometry.cppm --precompile -o geometry.pcm
clang++ -std=c++20 -c geometry.pcm -o geometry.o
clang++ -std=c++20 -fmodule-file=geometry=geometry.pcm \
    -c main.cpp -o main.o
clang++ geometry.o main.o -o app
./app
\end{lstlisting}

Dla GCC~14 (wsparcie modułów C++20 było w~tej wersji nadal rozwijane)
trzeba dodatkowo włączyć eksperymentalny mechanizm
\code{-fmodules-ts}. W~przykładach z~GCC~14 przydatne bywa jawne
wskazanie języka przy rozszerzeniu \code{.cppm}:
\begin{lstlisting}[language=bash,caption={Budowanie tego samego modu?u z GCC 14.}]
g++ -std=c++20 -fmodules-ts -x c++ -c geometry.cppm \
    -o geometry.o
# GCC tworzy potrzebne dane interfejsu w gcm.cache/.
g++ -std=c++20 -fmodules-ts -c main.cpp -o main.o
g++ geometry.o main.o -o app
./app
\end{lstlisting}

Dokładne opcje i~formaty modułów mogą zmieniać się między wersjami
narzędzi. W~dużym projekcie lepiej powierzyć zależności systemowi
budowania obsługującemu skanowanie importów, zamiast ręcznie
zarządzać plikami BMI.

\subsection{Jednostka interfejsu i~jednostka implementacji}

Moduł może rozdzielać interfejs od implementacji bez tradycyjnego
nagłówka. Jednostka implementacji deklaruje \code{module name;},
a~nie \code{export module name;}:

\begin{lstlisting}[caption={Dwa pliki jednego modu?u: interface.cppm i implementation.cpp.}]
// Plik: engine.cppm
export module engine;
export int run(int value);

// Plik: engine.cpp
module engine;
int run(int value) { return 2 * value; }
\end{lstlisting}

Importer widzi tylko deklaracje eksportowane przez interfejs.
Jednostkę implementacji również trzeba skompilować i~zlinkować.
Konkretny system budowania musi uwzględnić fakt, że jednostka
implementacji importuje interfejs swojego modułu niejawnie.

\subsection{Global module fragment i~nagłówki klasyczne}

Tradycyjny nagłówek może wymagać makr konfigurujących jego treść.
W~takim przypadku można skorzystać z~\emph{global module fragment}:

\begin{lstlisting}[caption={Global module fragment: tradycyjny include przed modu?em.}]
module;
#include <string>
#include <utility>

export module greeting;

export std::string greet(std::string name) {
    return "Hello " + std::move(name);
}
\end{lstlisting}

Dyrektywa \code{module;} rozpoczyna fragment globalny, a~następna
deklaracja \code{export module greeting;} rozpoczyna właściwą
jednostkę modułu. Importowanie nagłówków jednostkowych przez
\code{import <string>;} jest innym mechanizmem, tzw. \emph{header units},
których dostępność zależy od toolchainu. Nie utożsamiajmy go z~nazwanym
modułem standardowej biblioteki \code{import std;}.

\subsection{Partycje i~ponowny eksport importu}

Duży moduł można podzielić na partycje. Nazwa po dwukropku oznacza
partycję wewnątrz jednego modułu, a~nie odrębny nazwany moduł.

\begin{lstlisting}[caption={Dwie jednostki interfejsu modu?u z partycja.}]
// Plik: math_detail.cppm
export module mathlib:detail;
export int triple(int x) { return 3 * x; }

// Plik: mathlib.cppm
export module mathlib;
export import :detail;
export int times_six(int x) { return 2 * triple(x); }

// Plik: main.cpp
import mathlib;
int main() { return times_six(7) == 42 ? 0 : 1; }
\end{lstlisting}

Klient importuje \code{mathlib}, a~nie \code{mathlib:detail}.
W~tym przykładzie \code{export import :detail;} re-eksportuje
eksporty partycji. Wymaga to prawidłowej kolejności budowania
\emph{najpierw partycji, potem głównego interfejsu, a~na końcu klienta}.

\subsection{Łączność modułowa i~ODR}

C++20 wprowadza reguły własności deklaracji (\emph{module ownership}).
Deklaracje przypisane do nazwanego modułu należą do niego, a~pewne
nieeksportowane nazwy uzyskują \emph{module linkage}. Różni się to od
mechanicznego zastąpienia \code{\#include} przez \code{import}.

Przykładowo dodanie \code{export} nie sprawia automatycznie, że
interfejs C++ będzie zgodny ABI między kompilatorami. Ponadto
źródłowy kod szablonów nadal musi być semantycznie dostępny
odpowiedniemu mechanizmowi instancjacji, choć nie musi być już
tekstowo wklejany przez \code{\#include}.

\subsection{CMake i~\texttt{FILE\_SET CXX\_MODULES}}

Od CMake~3.28 dostępne są mechanizmy skanowania modułów i~specjalny
typ zbioru plików \code{CXX_MODULES}. Przykład projektu:

\begin{lstlisting}[language={},caption={CMakeLists.txt projektu korzystajacego z modu?ow C++20.}]
cmake_minimum_required(VERSION 3.28)
project(GeometryModules LANGUAGES CXX)

add_library(geometry)
target_compile_features(geometry PUBLIC cxx_std_20)
target_sources(geometry PUBLIC
    FILE_SET CXX_MODULES FILES geometry.cppm
)

add_executable(app main.cpp)
target_link_libraries(app PRIVATE geometry)
\end{lstlisting}

\begin{lstlisting}[language=bash,caption={Konfiguracja modu?ow przy uzyciu generatora Ninja.}]
cmake -S . -B build -G Ninja
cmake --build build
./build/app
\end{lstlisting}

Ten przykład pokazuje \emph{składnię} CMake, ale jego zbudowanie
wymaga zgodnego zestawu wersji CMake, kompilatora, generatora i~skanera
zależności. W~szczególności kompilator akceptujący składnię modułów
niekoniecznie współpracuje z~automatycznym skanowaniem w~danej
konfiguracji CMake. Na współczesnych systemach używa się m.in.
Ninja (co najmniej 1.11 w~typowych konfiguracjach) i~odpowiednich
wersji Clanga lub MSVC. Obsługa modułów standardowej biblioteki oraz
\emph{header units} ma odrębne ograniczenia.

\subsection{\texttt{import std;} nie jest tym samym co \texttt{\#include}}

Sam standard C++20 wprowadza językowy mechanizm modułów, ale
\emph{nie zobowiązuje} każdej biblioteki standardowej do dostarczenia
nazwanego modułu \code{std} w~tym standardzie. \code{import std;}
jest standaryzowane w~C++23; obsługa zależy od kompilatora, dostawcy
biblioteki standardowej i~wersji systemu budowania.

\begin{lstlisting}[caption={Nazwany modu? standardowej biblioteki (C++23, zalezny od wsparcia).}]
import std;

int main() {
    std::vector<int> xs{1, 2, 3};
    std::println("size = {}", xs.size());
}
\end{lstlisting}

Powyższy przykład wykorzystuje także \code{std::println} z~C++23.
Nie należy automatycznie kompilować go z~\code{-std=c++20} ani
zakładać, że \code{import std;} działa w~każdym środowisku C++23.

\section{Optymalizacja procesu budowania}
\label{sec:cpp17-build-perf}

\subsection{Incremental build i~zależności nagłówkowe}

Zmiana publicznego nagłówka może wymusić ponowną kompilację wielu TU,
podczas gdy zmiana prywatnego \code{.cpp} zwykle dotyka mniejszej
części projektu. Dlatego ograniczanie niepotrzebnych dołączeń
nagłówków, stosowanie deklaracji wyprzedzających tam, gdzie są
prawidłowe, i~PImpl mogą poprawiać czas budowania.

Moduły zmieniają strukturę zależności, ale nie usuwają konieczności
ich śledzenia; importer musi otrzymać aktualną reprezentację
interfejsu modułu.

\subsection{Precompiled headers, unity builds, ccache}

Popularne techniki przyspieszania kompilacji to:
\begin{description}[style=nextline]
\item[Precompiled headers (PCH)]
Przechowywanie wyniku parsowania stabilnych, obszernych nagłówków.
Pomagają w~projektach z~dużą liczbą TU, ale wymagają spójności
opcji kompilacji i~zarządzania unieważnianiem pamięci podręcznej.
\item[Unity builds]
Łączenie kilku plików źródłowych we wspólną TU. Może ograniczać
narzut kompilacji, ale też ujawniać kolizje anonimowych przestrzeni
nazw, makr i~definicji o~łączności wewnętrznej.
\item[Compilation cache]
Narzędzia takie jak \code{ccache} unikają ponownego wykonywania
kompilacji, jeżeli odpowiednio zdefiniowane wejścia pozostają takie same.
\item[Parallel build]
Kompilowanie niezależnych TU równolegle (np. \code{cmake --build build -j})
wykorzystuje wiele rdzeni, ale zwiększa chwilowe użycie pamięci.
\end{description}

\subsection{Powtarzalność buildów i~zależności transytywne}

Dla bibliotek publicznych warto utrzymywać przejrzystą strukturę
zależności. Przykładowo użytkownik nie powinien polegać na tym,
że \code{\#include <vector>} ,,przypadkiem'' pojawiło się wskutek
dołączenia innego nagłówka. Każdy nagłówek powinien zawierać
potrzebne deklaracje w~sposób niezależny od kolejności dołączeń.

Warto również zapisywać wersje kompilatora i~bibliotek, flagi
wpływające na ABI, wariant budowania oraz sposób instalowania
artefaktów. \emph{Reproducible builds} wymagają dodatkowo kontroli
źródeł niedeterministyczności, np. znaczników czasu i~ścieżek
zapisywanych w~plikach wynikowych.

\section{Praktyczny przewodnik diagnostyczny}
\label{sec:cpp17-troubleshooting}

\begin{center}
\small
\begin{tabularx}{\textwidth}{@{}p{0.34\textwidth}X@{}}
\toprule
Objaw & Najbardziej prawdopodobne źródło problemu\\
\midrule
\code{undefined reference} & Brakujący plik obiektowy, biblioteka,
definicja albo instancjacja szablonu.\\
\code{multiple definition} & Niepoprawna definicja zwykłej
funkcji lub zmiennej w~nagłówku albo podwójnie linkowany obiekt.\\
\code{undefined reference to vtable} & Niezdefiniowana wymagana
funkcja wirtualna (często destruktor) lub brak pliku implementacji.\\
\code{cannot find -lfoo} & Brak biblioteki lub nieprawidłowe
katalogi przeszukiwania linkera.\\
\code{cannot open shared object file} & Problem ścieżki lub
zależności ładowania, niekoniecznie problem kompilacji.\\
Kompilacja klienta działa, program kończy się błędem & Możliwa
niezgodność ABI, naruszenie ODR, zły czas życia obiektu lub błąd programu.\\
Importer nie znajduje modułu & Nie zbudowano BMI/CMI, nie wskazano
ścieżki albo konfiguracja skanera nie obsługuje importu.\\
Import działa, ale linker zgłasza brak funkcji & Brak pliku
obiektowego modułu/implementacji lub potrzebnej biblioteki.\\
Zmiana nagłówka przebudowuje prawie wszystko & Silne zależności
nagłówkowe; rozważ redukcję dołączeń, PImpl, moduły lub PCH.\\
\bottomrule
\end{tabularx}
\end{center}

\subsection{Kiedy wybrać nagłówki, bibliotekę czy moduł?}

\begin{itemize}
\item \textbf{Header-only} sprawdza się dla małych narzędzi, szablonów
      i~funkcji \code{inline}; ceną mogą być zależności kompilacyjne.
\item \textbf{Biblioteka statyczna} upraszcza dystrybucję kodu obiektowego
      wewnątrz aplikacji, ale nie usuwa wymagań ABI w~czasie linkowania.
\item \textbf{Biblioteka współdzielona} pozwala współdzielić kod
      i~wymieniać implementację, lecz wymaga jasnej polityki ABI,
      widoczności symboli i~wersjonowania.
\item \textbf{Moduł C++20} służy do semantycznego podziału interfejsu
      źródłowego i~kontroli importów; nie jest alternatywą wobec
      procesu linkowania ani automatycznym formatem dystrybucji ABI.
\item \textbf{PImpl i~C API} rozważaj na granicach bibliotek o~długim
      cyklu życia, szczególnie gdy stabilność ABI ma duże znaczenie.
\end{itemize}

\section*{Podsumowanie rozdziału}
\addcontentsline{toc}{section}{Podsumowanie rozdziału}

Kompletny program C++ jest wynikiem wspólnej pracy preprocessora,
kompilatora, linkera i~loadera. Organizując projekt, trzeba osobno
rozumieć \emph{deklaracje i~definicje}, \emph{ODR i~łączność},
\emph{symbole i~biblioteki} oraz \emph{ABI}. Nagłówki i~moduły
pomagają udostępniać interfejsy źródłowe; biblioteki statyczne
i~współdzielone dostarczają kod do połączenia lub załadowania.

Dobra organizacja projektu polega na jasnym rozdziale API i~implementacji,
ograniczaniu zależności, korzystaniu z~targetów CMake oraz świadomym
wersjonowaniu interfejsu binarnego. Nowoczesne moduły upraszczają część
problemów nagłówków, ale wprowadzają wymagania dotyczące narzędzi,
BMI/CMI i~kolejności kompilacji.

\section*{Literatura i~dokumentacja}
\addcontentsline{toc}{section}{Literatura i dokumentacja}
\begin{itemize}
\item ISO/IEC 14882:2020, zagadnienia \emph{translation units},
      \emph{linkage}, \emph{One Definition Rule} i~\emph{modules}.
\item \url{https://en.cppreference.com/w/cpp/language/definition}
      --- ODR, deklaracje i~definicje.
\item \url{https://en.cppreference.com/w/cpp/language/modules}
      --- składnia i~semantyka modułów.
\item \url{https://gcc.gnu.org/onlinedocs/gcc/C_002b_002b-Modules.html}
      --- implementacja modułów w~GCC.
\item \url{https://clang.llvm.org/docs/StandardCPlusPlusModules.html}
      --- moduły w~Clangu, BMI i~komendy budowania.
\item \url{https://cmake.org/cmake/help/latest/manual/cmake-cxxmodules.7.html}
      --- CMake i~zależności modułów.
\item \url{https://cmake.org/cmake/help/latest/command/target_sources.html}
      --- \code{FILE_SET CXX_MODULES}.
\item \url{https://gcc.gnu.org/onlinedocs/gcc/Code-Gen-Options.html}
      --- widoczność symboli w~ELF.
\item John Lakos, \emph{Large-Scale C++ Software Design}, Addison-Wesley, 1996.
\end{itemize}

% Komendy powloki sa przeznaczone glownie dla GCC/Clang na Linux/ELF;
% roznice Windows/MSVC i macOS zaznaczono w tekscie.

\chapter{Kompilacja, linkowanie, ABI i moduły C++20}
\label{chap:cpp17-build-abi-modules}

\section{Dlaczego trzeba rozumieć proces budowania?}

W~poprzednich rozdziałach projektowaliśmy klasy, szablony, mechanizmy
zarządzania pamięcią i~współbieżność. Dopóki cały program mieści się
w~jednym pliku, wystarcza zwykle wywołanie \code{g++ main.cpp -o app}.
W~realnym projekcie kod dzieli się jednak na dziesiątki lub tysiące
\emph{jednostek translacji}, biblioteki zewnętrzne i~moduły. To wtedy
pojawiają się komunikaty \code{undefined reference}, \code{multiple definition},
konflikty wersji bibliotek, długi czas budowania oraz błędy ujawniające
się dopiero po wymianie biblioteki współdzielonej.

Musimy rozdzielić trzy różne zagadnienia:
\begin{description}[style=nextline]
\item[Poprawność programu w~języku C++]
Czy deklaracje, definicje, typy i~wywołania spełniają reguły C++,
w~szczególności \emph{One Definition Rule} (ODR)?
\item[Model budowania i~linkowania]
Jakie pliki są kompilowane osobno, gdzie kompilator zapisuje symbole,
w~jakiej kolejności linker przeszukuje biblioteki oraz jak ładowane są
biblioteki dynamiczne?
\item[Zgodność binarna (ABI)]
Czy dwa osobno zbudowane komponenty zgadzają się co do reprezentacji
typów, sposobu wywoływania funkcji, nazw symboli i~pozostałych reguł
interfejsu binarnego?
\end{description}

\nota{Norma ISO C++ opisuje semantykę programu, ale nie narzuca jednego
formatu pliku obiektowego, konwencji \emph{name mangling}, organizacji
bibliotek \code{.so}/\code{.dll} ani poleceń linkera. Poniżej
odróżniamy reguły języka od typowych rozwiązań narzędziowych.}

\section{Od pliku źródłowego do programu}

\begin{center}\small
\begin{tabularx}{\textwidth}{@{}p{0.31\textwidth}X@{}}
\toprule
Etap & Za co odpowiada \\
\midrule
Preprocesor & Rozwija dyrektywy \code{#include}, makra i warunkową kompilację. \\
Kompilator & Sprawdza składnię i typy oraz wytwarza kod dla jednostki translacji. \\
Asembler & Zamienia kod asemblera na plik obiektowy z symbolami i relokacjami. \\
Linker & Łączy pliki obiektowe i biblioteki, rozwiązując odwołania do symboli. \\
Loader systemowy & Ładuje program i wymagane biblioteki współdzielone przed uruchomieniem. \\
\bottomrule
\end{tabularx}
\end{center}
\label{sec:cpp17-pipeline}

\subsection{Preprocesor, kompilator, asembler, linker i~loader}

W~uproszczeniu klasyczny projekt z~nagłówkami przechodzi przez etapy:
\[
\begin{aligned}
\text{źródła i nagłówki}&\ \longrightarrow\ \text{preprocessing}
\ \longrightarrow\ \text{kompilacja}\\
&\phantom{\ \longrightarrow\ }\longrightarrow\ \text{plik obiektowy}
\ \longrightarrow\ \text{linkowanie}
\ \longrightarrow\ \text{program}.
\end{aligned}
\]

Na niektórych platformach kompilator wytwarza tekst asemblerowy
(\code{.s}) jako etap pośredni; w~innych przejście do kodu maszynowego
jest zintegrowane. Linkowanie może odbywać się przy budowaniu programu,
a~dodatkowe rozwiązywanie symboli może następować w~czasie ładowania
lub później.

Prześledźmy etapy na samodzielnym programie.

\begin{lstlisting}[caption={Plik hello.cpp: prosty program do obserwowania kompilacji.}]
#include <iostream>

int twice(int x) { return 2 * x; }

int main() {
    std::cout << twice(21) << '\n';
}
\end{lstlisting}

\begin{lstlisting}[language=bash,caption={Wybrane etapy budowania przy uzyciu GCC/Clang.}]
g++ -std=c++20 -E hello.cpp -o hello.ii  # Preprocesor
# Produkt poprzedniego etapu moze byc duzy.
g++ -std=c++20 -S hello.cpp -o hello.s   # Asembler

g++ -std=c++20 -c hello.cpp -o hello.o   # Kod obiektowy
g++ hello.o -o hello                      # Linkowanie
./hello                                   # 42
\end{lstlisting}

Opcja \code{-c} nie tworzy wykonywalnego programu; produkuje plik
obiektowy, który może zawierać nierozwiązane odwołania do symboli
pochodzących z~innych jednostek translacji. Polecenie \code{-S}
kończy budowanie na asemblerze, a~\code{-E} na wstępnym przetwarzaniu.
Warto odróżnić program sterujący kompilacją (np.~\code{g++}) od samego
frontendu C++, asemblera i~linkera.

\subsection{Pliki obiektowe, symbole i~relokacje}

Plik obiektowy zawiera zwykle kod maszynowy, dane, tablice symboli
oraz \emph{relokacje}, czyli informacje, które adresy muszą zostać
skorygowane po ustaleniu ostatecznego układu programu. W~środowisku
Linux powszechny jest format ELF; Windows używa PE/COFF, a~macOS
Mach-O. Szczegóły tych formatów nie są częścią języka C++.

\begin{lstlisting}[language=bash,caption={Ogladanie symboli i relokacji (Linux/ELF).}]
nm -C hello.o                 # Symbole; -C demangluje nazwy C++
readelf -Ws hello.o          # Tablica symboli ELF
readelf -r hello.o           # Relokacje
objdump -drC hello.o         # Asembler i relokacje
\end{lstlisting}

Przykładowo litera \code{T} w~typowym wyjściu \code{nm} informuje
o~symbolu definiowanym w~sekcji kodu, a~\code{U} o~symbolu
niezdefiniowanym w~danym obiekcie, który trzeba rozwiązać. Nie
należy traktować oznaczeń \code{nm} jako przenośnej własności C++.

\subsection{Jednostka translacji i~kompilacja oddzielna}

W~tradycyjnym modelu \emph{jednostka translacji} (\emph{translation unit},
TU) powstaje z~pliku źródłowego po uwzględnieniu dyrektyw preprocesora,
w~szczególności zawartości dołączanych nagłówków. Każda TU jest
analizowana niezależnie, a~linker łączy otrzymane obiekty.

Rozważmy projekt:
\begin{lstlisting}[caption={Trzy pliki klasycznego projektu C++: API, implementacja i klient.}]
// Plik: math.hpp
#ifndef CPP17_MATH_HPP
#define CPP17_MATH_HPP
int square(int x);
#endif

// Plik: math.cpp
#include "math.hpp"
int square(int x) { return x * x; }

// Plik: main.cpp
#include "math.hpp"
#include <iostream>
int main() { std::cout << square(7) << '\n'; }
\end{lstlisting}

\begin{lstlisting}[language=bash,caption={Niezalezne budowanie TU i koncowe linkowanie.}]
g++ -std=c++20 -Wall -Wextra -c math.cpp -o math.o
g++ -std=c++20 -Wall -Wextra -c main.cpp -o main.o
g++ main.o math.o -o app
./app                        # 49
\end{lstlisting}

\code{math.hpp} dostarcza deklarację klientowi i~implementacji,
ale nie jest zwykle oddzielnym obiektem linkowanym. Zmiana samego
\code{math.cpp} nie musi wymuszać ponownej kompilacji \code{main.cpp},
o~ile nie zmienił się kontrakt w~nagłówku lub pozostałe zależności.

\subsection{Przykładowe komendy w~MSVC}

W~\emph{Developer Command Prompt for Visual Studio} można zbudować
analogiczny przykład za pomocą \code{cl.exe} i~\code{link.exe}:

\begin{lstlisting}[language=bash,caption={Kompilacja dwóch plików w środowisku MSVC/Windows.}]
cl /std:c++20 /EHsc /c math.cpp
cl /std:c++20 /EHsc /c main.cpp
link main.obj math.obj /OUT:app.exe
app.exe
\end{lstlisting}

Przy budowaniu programów C++ w~Visual Studio warto również
kontrolować zgodność wersji MSVC, ustawienia biblioteki uruchomieniowej
(np.~\code{/MD} wobec \code{/MT}) i~konfigurację Debug/Release.
Takie flagi nie mają bezpośrednich odpowiedników narzuconych przez
standard języka; należą do konkretnego środowiska.

\section{Nagłówki, makra i~zasada jednej definicji}
\label{sec:cpp17-headers-odr}

\subsection{Include guards i~\texttt{\#pragma once}}

Zawartość \code{\#include} jest w~uproszczeniu wstawiana do źródła
przez preprocesor. Wielokrotne dołączenie tego samego nagłówka
w~jednej TU mogłoby więc prowadzić do powtórzenia definicji klasy.
Chronimy się przez \emph{include guard}:

\begin{lstlisting}[caption={Klasyczna os?ona nag?owka.}]
#ifndef LIB_WIDGET_HPP
#define LIB_WIDGET_HPP

struct Widget {
    int value{};
};

#endif // LIB_WIDGET_HPP
\end{lstlisting}

W~praktyce popularna jest też dyrektywa \code{\#pragma once},
obsługiwana przez najważniejsze kompilatory, lecz niestandaryzowana
przez ISO C++. Oba rozwiązania chronią \emph{jedną TU} przed
powtórnym dołączeniem pliku. Nie rozwiązują automatycznie problemu
\emph{wielu definicji tego samego nie-inline symbolu w~różnych TU}.

\subsection{ODR: deklaracja kontra definicja}

Deklaracja wprowadza nazwę i~informację o~typie, natomiast definicja
zwykle dostarcza byt danej funkcji lub obiektu. Deklaracji pewnej
funkcji może być wiele, ale \emph{nie-inline} funkcja z~zewnętrzną
łącznością, używana w~programie w~sposób wymagający definicji,
powinna mieć jedną definicję w~programie. Reguły ODR zawierają
wyjątki pozwalające na identyczne definicje w~różnych TU, np. dla
odpowiednich funkcji \code{inline}, klas i~szablonów.

\begin{lstlisting}[caption={Typowa pu?apka: definicja zwyk?ej funkcji w nag?owku.}]
// BAD: wrong.hpp, dolaczony do dwoch plikow .cpp
int twice(int x) { return 2 * x; }
// Dwie zewnetrzne definicje -> naruszenie ODR.
\end{lstlisting}

\begin{lstlisting}[caption={Definicja inline przeznaczona do nag?owka.}]
// Plik: twice.hpp
#pragma once
inline int twice(int x) { return 2 * x; }
\end{lstlisting}

\code{inline} nie nakazuje kompilatorowi wstawienia instrukcji funkcji
bezpośrednio w~miejscu wywołania. Jego znaczenie dla organizacji
programu polega przede wszystkim na dopuszczeniu definicji w~wielu TU
zgodnie z~ODR (z~odpowiednimi warunkami). Wszystkie takie definicje
muszą być zgodne i~spełniać reguły ODR; samo dopisanie \code{inline}
nie naprawia różniących się implementacji w~dwóch plikach.

\subsection{Zmienne \texttt{extern}, \texttt{inline}, \texttt{const} i~\texttt{static}}

\begin{lstlisting}[caption={Jedna definicja zmiennej globalnej, wiele deklaracji.}]
// Plik: settings.hpp
#pragma once
extern int request_limit;       // Tylko deklaracja

// Plik: settings.cpp
#include "settings.hpp"
int request_limit = 100;        // Definicja

// Plik: main.cpp
#include "settings.hpp"
int main() { return request_limit == 100 ? 0 : 1; }
\end{lstlisting}

Od C++17 w~nagłówkach można również definiować \emph{zmienne inline}:
\begin{lstlisting}[caption={Zmienna inline w nag?owku (C++17).}]
#pragma once
inline constexpr int default_timeout_ms = 250;
inline int global_counter = 0;
\end{lstlisting}

W~tym przypadku odpowiednie definicje we wszystkich TU mogą
wskazywać na ten sam obiekt o~zewnętrznej łączności. W~kontraście
zmienna \code{const} zdefiniowana w~przestrzeni nazw bez \code{extern}
ma w~typowym przypadku łączność wewnętrzną; oddzielne TU mogą mieć
wówczas odrębne egzemplarze. \code{static} na poziomie przestrzeni
nazw także zapewnia łączność wewnętrzną, lecz na współczesnym C++
zwykle preferujemy anonimową przestrzeń nazw.

\begin{lstlisting}[caption={Nazwy pomocnicze ograniczone do jednej TU.}]
// Plik: helper.cpp
namespace {
    int internal_counter = 0;
    int twice_local(int x) { return x * 2; }
}

int perform(int x) {
    ++internal_counter;
    return twice_local(x);
}
\end{lstlisting}

\subsection{Szablony: dlaczego definicje zwykle znajdują się w~nagłówkach?}

Kompilator instancjuje szablon dla określonych argumentów typowych
w~miejscu, w~którym wymagana jest jego definicja. Samo udostępnienie
deklaracji z~pliku \code{.hpp} i~ukrycie ogólnej definicji szablonu
w~oddzielnym \code{.cpp} może zatem prowadzić do błędu linkera.

\begin{lstlisting}[caption={Najprostszy, tradycyjny sposob udostepniania szablonu.}]
// Plik: multiply.hpp
#pragma once

template<class T>
T multiply(T a, T b) {
    return a * b;
}
\end{lstlisting}

Można także jawnie instancjować \emph{konkretne} typy w~pliku
implementacji. Oddzielne jednostki translacji nie muszą wtedy
instancjować tego samego kodu niezależnie:

\begin{lstlisting}[caption={Jawna instancjacja okreslonego wariantu szablonu.}]
// Plik: operation.hpp
#pragma once
template<class T>
T add(T a, T b);
extern template int add<int>(int, int);

// Plik: operation.cpp
#include "operation.hpp"
template<class T>
T add(T a, T b) { return a + b; }
template int add<int>(int, int);

// Plik: main.cpp
#include "operation.hpp"
int main() { return add(10, 32) == 42 ? 0 : 1; }
\end{lstlisting}

W~tym przykładzie klient może korzystać z~\code{add<int>}, ale
instancjacja innych typów nie zostaje mu udostępniona automatycznie.
Definicje szablonów umieszczane w~interfejsach modułów podlegają
nieco innemu modelowi widoczności, który omówimy dalej.

\subsection{Naruszenie ODR bez łatwego komunikatu}

Nie wszystkie naruszenia ODR skutkują natychmiastowym błędem
linkera. Szczególnie groźne jest warunkowe definiowanie tej samej
klasy w~dwóch TU w~inny sposób, np. przy użyciu sprzecznych
makr kompilacji:

\begin{lstlisting}[caption={Przyk?ad niezgodnych definicji typu (celowo b?edny).}]
// Plik: config.hpp
struct Config {
#ifdef EXTRA_FIELD
    int extra;
#endif
    int value;
};

// Jedna TU z -DEXTRA_FIELD, inna bez tej flagi:
// rozne definicje klasy -> naruszenie ODR.
\end{lstlisting}

Taki program może się nawet zlinkować, lecz pozostaje niepoprawny.
Dlatego opcje zmieniające publiczne definicje klas i~struktur muszą
być spójne we wszystkich TU danego programu.

\section{Symbole, łączność i~typowe błędy linkera}
\label{sec:cpp17-symbols}

\subsection{Łączność a~widoczność symbolu}

\emph{Linkage} w~języku C++ mówi, czy deklaracje w~różnych zakresach
i~TU mogą nazywać tę samą encję. W~praktyce spotykamy między innymi
łączność zewnętrzną, wewnętrzną, a~w~jednostkach modułów także
łączność modułową. Nie należy jednak utożsamiać tej własności
języka z~\emph{visibility} na poziomie ELF ani z~eksportem symbolu
z~biblioteki dynamicznej.

\begin{center}
\begin{tabularx}{\textwidth}{@{}lX@{}}
\toprule
Mechanizm & Co określa?\\
\midrule
\code{static} w~przestrzeni nazw & Łączność wewnętrzna nazwy w~TU.\\
\code{extern} & M.in. deklarowanie obiektu definiowanego w~innym miejscu.\\
Anonimowa przestrzeń nazw & Wewnętrzna łączność elementów pomocniczych.\\
\code{inline} & Warunki wielokrotnych definicji zgodnych z~ODR.\\
\code{-fvisibility=hidden} & Domyślna widoczność symboli w~typowych bibliotekach ELF.\\
\code{__declspec(dllexport)} & Udostępnianie symboli biblioteki DLL w~MSVC.\\
\bottomrule
\end{tabularx}
\end{center}

\subsection{Name mangling i~\texttt{extern~\textquotedbl{}C\textquotedbl{}}}

C++ dopuszcza przeciążanie funkcji, przestrzenie nazw oraz rozbudowane
typy szablonowe. Aby móc odróżniać symbole po linkowaniu,
implementacje zwykle kodują część tych informacji w~nazwach
binarnych (\emph{name mangling}). Dokładny sposób kodowania nie jest
ustalony przez ISO C++.

\begin{lstlisting}[caption={Przeciazanie i punkt wejscia C ABI.}]
namespace calc {
    int f(int x) { return x + 1; }
    double f(double x) { return x + 1.0; }
}

extern "C" int api_add(int a, int b) {
    return a + b;
}
\end{lstlisting}

Na typowej platformie ELF/Itanium ABI funkcje przeciążone otrzymują
różne symbole zakodowane przez kompilator, podczas gdy \code{api_add}
ma nazwę zgodną z~łącznością języka C. Nie jest to gwarancja pełnej
zgodności ABI z~dowolną biblioteką C: znaczenie mają także deklaracje
typów i~konwencja wywołania ustalona dla danej platformy.

\begin{lstlisting}[language=bash,caption={Demanglowanie i inspekcja symboli (przyk?ady narzedziowe).}]
g++ -std=c++20 -c overload.cpp -o overload.o
nm overload.o
nm -C overload.o
c++filt _Z3fooi      # Ilustracja popularnego schematu manglingu
\end{lstlisting}

\subsection{\texttt{undefined reference}, \texttt{multiple definition}}

Błąd \code{undefined reference} oznacza zwykle, że linker znalazł
odwołanie do symbolu, lecz nie znalazł odpowiedniej definicji.
Przykładowe przyczyny to pominięcie \code{math.o}, pominięcie
biblioteki, niezgodny podpis funkcji albo brak jawnej instancjacji
szablonu. Z~kolei \code{multiple definition} sugeruje m.in. zwykłą
funkcję lub zmienną zdefiniowaną w~nagłówku dołączanym do wielu TU.

Przydatna lista kontrolna:
\begin{enumerate}
\item Sprawdź, czy wszystkie wymagane \code{.cpp} są kompilowane i~linkowane.
\item Porównaj podpis funkcji: przestrzeń nazw, typy, kwalifikatory i~ABI.
\item Upewnij się, że deklaracja i~definicja pochodzą z~tego samego interfejsu.
\item Sprawdź kolejność bibliotek statycznych i~ich zależności.
\item Dla klas polimorficznych sprawdź brakujące definicje funkcji
      wirtualnych, w~tym destruktora wirtualnego, który może wymagać
      definicji nawet jako funkcja \code{= 0}.
\item Dla szablonów sprawdź dostępność definicji w~punkcie instancjacji
      lub odpowiednią instancjację jawną.
\end{enumerate}

\section{Biblioteki statyczne}
\label{sec:cpp17-static-libraries}

\emph{Biblioteka statyczna} jest typowo archiwum plików obiektowych.
W~Linux ma zwykle rozszerzenie \code{.a}, w~MSVC \code{.lib}.
Biblioteka sama nie jest osobnym procesem ani automatycznie tworzonym
wątkiem. Linker wybiera z~niej obiekty potrzebne do rozwiązania
zależności i~tworzy wynikowy program lub inną bibliotekę.

\begin{lstlisting}[language=bash,caption={Budowanie i uzycie biblioteki statycznej (Linux/GCC).}]
g++ -std=c++20 -O2 -c math.cpp -o math.o
ar rcs libmathstuff.a math.o

g++ -std=c++20 -O2 -c main.cpp -o main.o
g++ main.o -L. -lmathstuff -o app
./app
\end{lstlisting}

W~typowym poleceniu GNU linkera biblioteki archiwalne warto umieszczać
\emph{po} obiektach, które odwołują się do ich symboli. Przykładowo
\code{g++ main.o -lmathstuff} jest zazwyczaj właściwszym porządkiem
niż \code{g++ -lmathstuff main.o}. Linkowanie zależności cyklicznych
bibliotek statycznych może wymagać ponownego wskazania biblioteki
lub użycia grupy linkera (np.~\code{--start-group}/\code{--end-group}
w~GNU ld); lepiej jednak unikać takich zależności architektonicznie.

\nota{Statyczna biblioteka nie oznacza automatycznie w~pełni
statycznego programu. Aplikacja może dalej korzystać z~bibliotek
współdzielonych systemu i~standardowej biblioteki C++.}

\section{Biblioteki współdzielone i~dynamiczne ładowanie}
\label{sec:cpp17-shared-libraries}

\subsection{Budowanie biblioteki współdzielonej}

Na systemach ELF biblioteki współdzielone są zazwyczaj plikami
\code{.so}. Obiekty wykorzystywane w~takich bibliotekach zwykle
buduje się jako kod niezależny od położenia w~pamięci (\emph{PIC}).

\begin{lstlisting}[language=bash,caption={Biblioteka .so oraz program klienta (Linux/GCC).}]
g++ -std=c++20 -O2 -fPIC -c math.cpp -o math_pic.o
g++ -shared math_pic.o -o libmathstuff.so

g++ -std=c++20 -O2 main.cpp -L. -lmathstuff \
    '-Wl,-rpath,$ORIGIN' -o app
./app
\end{lstlisting}

Wyrażenie \code{\$ORIGIN} odnosi się w~typowym loaderze ELF do
katalogu, w~którym znajduje się ładowany obiekt. Apostrofy w~poleceniu
powłoki zapobiegają wcześniejszej ekspansji \code{\$ORIGIN} przez
powłokę. W~projekcie produkcyjnym należy przemyśleć ścieżki
instalacyjne, pakowanie bibliotek i~bezpieczeństwo ładowania.

Biblioteki współdzielone pozwalają wielu procesom mapować ich kod
z~tego samego pliku oraz wymieniać implementację bez konieczności
pełnego linkowania programu w~czasie jego budowania. Nie gwarantuje
to jednak, że \emph{dowolna} nowa wersja zachowa zgodność ABI.

\subsection{Wyszukiwanie bibliotek: RPATH/RUNPATH, SONAME}

Na Linux/ELF linker i~loader korzystają z~mechanizmów takich jak
\emph{SONAME}, \emph{RPATH} i~\emph{RUNPATH}. Reguły wyszukiwania
zależą od systemu i~konfiguracji; środowisko wykonania może również
uwzględniać ścieżki systemowe lub zmienne środowiskowe.

\begin{lstlisting}[language=bash,caption={Diagnostyka zaleznosci dynamicznych (Linux).}]
readelf -d ./app | grep -E 'NEEDED|RPATH|RUNPATH'
ldd ./app                  # Tylko dla zaufanych plikow
readelf -Ws libmathstuff.so
\end{lstlisting}

Programów nieznanego pochodzenia nie należy lekkomyślnie przekazywać
do narzędzi, które mogą uruchamiać kod lub loader; do statycznej
inspekcji lepiej nadaje się \code{readelf}.

\subsection{Eksportowanie symboli i~kontrola API}

Domyślne eksportowanie zbyt wielu symboli może powodować kolizje
nazw i~powiększać powierzchnię ABI. W~ELF można budować bibliotekę
z~\code{-fvisibility=hidden}, a~publiczne symbole jawnie oznaczać:

\begin{lstlisting}[caption={Przenosne w granicach GCC/Clang/MSVC makro eksportu.}]
// Plik: api.hpp
#pragma once

#if defined(_WIN32)
  #if defined(MYLIB_BUILD_SHARED)
    #define MYLIB_API __declspec(dllexport)
  #elif defined(MYLIB_USE_SHARED)
    #define MYLIB_API __declspec(dllimport)
  #else
    #define MYLIB_API
  #endif
#elif defined(__GNUC__) || defined(__clang__)
  #define MYLIB_API __attribute__((visibility("default")))
#else
  #define MYLIB_API
#endif

extern "C" MYLIB_API int mylib_add(int, int);
\end{lstlisting}

W~MSVC zwykle stosuje się projekt DLL oraz odpowiednią bibliotekę
importową; na macOS biblioteki mają najczęściej rozszerzenie
\code{.dylib}. Flagi \code{-fPIC}, \code{-shared} i~\code{-Wl,...}
w~powyższej postaci dotyczą konkretnego toolchainu, nie ISO C++.

\subsection{\texttt{dlopen}, \texttt{dlsym} i~granica interfejsu wtyczki}

Bibliotekę można także załadować podczas działania procesu, bez
linkowania jej wprost z~aplikacją. Na POSIX często używa się
\code{dlopen}/\code{dlsym}, natomiast na Windows odpowiedników
\code{LoadLibrary}/\code{GetProcAddress}.

\begin{lstlisting}[caption={Prosta wtyczka C ABI: biblioteka i klient POSIX (dwa pliki).}]
// Plik: plugin.cpp
extern "C" int plugin_square(int x) {
    return x * x;
}

// Plik: loader.cpp
#include <dlfcn.h>
#include <iostream>
#include <stdexcept>

int main() {
    void* lib = dlopen("./libplugin.so", RTLD_NOW);
    if (!lib) throw std::runtime_error(dlerror());

    using Fn = int (*)(int);
    void* address = dlsym(lib, "plugin_square");
    if (!address) {
        dlclose(lib);
        throw std::runtime_error("Missing plugin_square");
    }

    // POSIX dopuszcza konwersje dlsym -> wskaznik do funkcji.
    Fn square = reinterpret_cast<Fn>(address);
    std::cout << square(7) << '\n'; // 49
    dlclose(lib);
}
\end{lstlisting}

\begin{lstlisting}[language=bash,caption={Budowanie przyk?adu wtyczki (Linux/ELF).}]
g++ -std=c++20 -fPIC -shared plugin.cpp -o libplugin.so
g++ -std=c++20 loader.cpp -ldl -o loader
./loader
\end{lstlisting}

\nota{To przykład dla POSIX, a~nie w~pełni przenośna konstrukcja
ISO C++. Również rzutowanie wyniku \code{dlsym} na wskaźnik funkcji
opiera się na wymaganiach POSIX, a~nie na ogólnej gwarancji C++.
Dodatkowo nie wolno wyładowywać biblioteki, dopóki jakakolwiek funkcja
lub obiekt z~tej biblioteki nadal może być użyty. Przy bardziej
złożonych wtyczkach potrzebny jest kontrakt wersji ABI i~zarządzania
czasem życia.}

\section{ABI --- Application Binary Interface}
\label{sec:cpp17-abi}

\subsection{API a~ABI: dwa różne kontrakty}

\emph{API} opisuje przede wszystkim dostępne nazwy, typy, funkcje
i~semantykę na poziomie źródła. \emph{ABI} określa, jak te elementy
reprezentowane są w~kodzie maszynowym. ABI obejmuje między innymi:
\begin{itemize}
\item konwencje wywołania: przekazywanie parametrów, wartość zwracana,
      rejestry i~ułożenie stosu;
\item reprezentację i~wyrównanie obiektów, w~tym pól klas;
\item nazwy symboli C++, RTTI, mechanizmy wyjątków i~wirtualnych wywołań;
\item reguły udostępniania i~ładowania symboli bibliotek;
\item ABI standardowej biblioteki oraz ustawienia kompilatora.
\end{itemize}

Dwie biblioteki mogą być zgodne źródłowo i~jednocześnie binarnie
niezgodne. W~szczególności identyczny nagłówek nie oznacza, że plik
\code{.so} zbudowany innym toolchainem lub z~innym zestawem opcji
na pewno współpracuje z~aplikacją.

\subsection{Zmiana układu obiektu łamie zgodność binarną}

Rozważmy klasę publiczną:
\begin{lstlisting}[caption={Zmiana prywatnych pol nadal moze zmienic ABI.}]
// Wersja biblioteki 1
class Record {
    int id_{};
public:
    int id() const;
};

// Wersja biblioteki 2
class Record {
    int id_{};
    double score_{}; // Nawet prywatne pole zmienia uklad.
public:
    int id() const;
};
\end{lstlisting}

Klient kompiluje np.~\code{sizeof(Record)} według nagłówka,
który widział podczas kompilacji. Gdy zamienimy jedynie bibliotekę
na wersję z~innym układem \code{Record}, może dojść do naruszenia
kontraktu pamięci, nawet jeśli \code{id()} nadal ma taki sam podpis.
Podobne ryzyko dotyczy dodawania wirtualnych funkcji i~zmian
hierarchii dziedziczenia.

\subsection{Wzorzec PImpl}

Aby zmniejszyć zależności kompilacyjne i~ukryć układ implementacji,
można przechowywać wskaźnik do niekompletnego typu. Dzięki temu
zmiana \emph{pól implementacyjnych} często nie zmienia rozmiaru
publicznego obiektu.

\begin{lstlisting}[caption={PImpl: plik widget.hpp.}]
#pragma once
#include <memory>

class Widget {
public:
    Widget();
    ~Widget();

    Widget(Widget&&) noexcept;
    Widget& operator=(Widget&&) noexcept;

    Widget(const Widget&) = delete;
    Widget& operator=(const Widget&) = delete;

    int value() const;
    void set_value(int value);

private:
    struct Impl;
    std::unique_ptr<Impl> impl_;
};
\end{lstlisting}

\begin{lstlisting}[caption={PImpl: plik widget.cpp.}]
#include "widget.hpp"
#include <utility>

struct Widget::Impl {
    int value = 0;
    // Tu moga pojawiac sie dalsze pola.
};

Widget::Widget() : impl_(std::make_unique<Impl>()) {}
Widget::~Widget() = default;
Widget::Widget(Widget&&) noexcept = default;
Widget& Widget::operator=(Widget&&) noexcept = default;

int Widget::value() const { return impl_->value; }
void Widget::set_value(int value) { impl_->value = value; }
\end{lstlisting}

Destruktor jest zdefiniowany w~\code{widget.cpp}, gdzie \code{Impl}
jest kompletny; to istotne przy \code{std::unique_ptr<Impl>}.
PImpl nie zapewnia automatycznie pełnej stabilności ABI: nadal
liczy się zgodność kompilatora, standardowej biblioteki, eksportu
symboli, konwencji wywołania i~pozostałych publicznych deklaracji.
Dodatkowy wskaźnik i~alokacja na stercie mogą zwiększać koszt dostępu.

\subsection{C ABI jako ostrożna granica interoperacyjności}

Dla wtyczek i~bibliotek wykorzystywanych z~różnych języków często
projektuje się wąski interfejs oparty na funkcjach \code{extern "C"}
oraz nieprzezroczystych uchwytach (\emph{opaque handles}). Zamiast
przekazywać \code{std::string} lub klasy z~funkcjami wirtualnymi,
przekazuje się proste typy i~wskaźniki.

\begin{lstlisting}[caption={Przyk?adowy szkic stabilniejszego interfejsu binarnego.}]
#ifdef __cplusplus
extern "C" {
#endif

struct engine_handle;

engine_handle* engine_create(void);
void engine_destroy(engine_handle* p);
int engine_eval(engine_handle* p, double x, double* result);

#ifdef __cplusplus
}
#endif
\end{lstlisting}

Podpis funkcji nadal wymaga uzgodnienia ABI danej platformy,
konwencji wywołania i~kontraktu błędów. Przekraczanie takiej granicy
wyjątkami C++, pamięcią zwalnianą przez niezgodny alokator lub obiektami
standardowej biblioteki jest niebezpieczne. Bezpieczniej zwracać
kody błędu, a~alokację i~zwalnianie łączyć w~pary po tej samej
stronie interfejsu.

\subsection{Wersje kompilatora i~biblioteki standardowej}

Istnieją różne rodziny ABI C++ (np. rodzina Itanium ABI stosowana na
wielu platformach uniksowych oraz ABI MSVC w~Windows). Nawet w~ramach
jednej rodziny znaczenie mają flagi kompilatora, ustawienia standardowej
biblioteki, opcje wyjątków, RTTI i~ABI konkretnych typów.
Przykładowo GCC/libstdc++ używa mechanizmu tzw. \emph{dual ABI}
dla niektórych typów standardowych (m.in. \code{std::string}); mieszanie
niezgodnych ustawień może skutkować błędami linkera lub działania.
Nie należy zakładać, że biblioteki zbudowane dowolnymi wersjami
GCC, Clanga i~MSVC są wymienne binarnie.

\section{System budowania: CMake i~projekt wieloplikowy}
\label{sec:cpp17-cmake}

\subsection{Dlaczego targety są ważniejsze niż globalne flagi?}

Ręczne wywoływanie kompilatora dobrze tłumaczy mechanikę, lecz dla
większych projektów potrzebujemy grafu zależności, wariantów
Debug/Release, testów i~kontroli opcji kompilatora. Współczesny
CMake modeluje budowanie przez \emph{targety}: bibliotekę, program,
interfejsową bibliotekę nagłówkową i~ich właściwości.

Przykładowa struktura:
\begin{lstlisting}[caption={Uk?ad plikow klasycznej biblioteki CMake.}]
project/
  CMakeLists.txt
  include/
    mathstuff/math.hpp
  src/
    math.cpp
    main.cpp
\end{lstlisting}

W~tej konkretnej strukturze pliki źródłowe mogą zaczynać się od
\code{\#include "mathstuff/math.hpp"}, a~\code{main.cpp} również
korzysta z~tej samej ścieżki nagłówka. To dzięki
\code{target_include_directories} kompilator znajduje katalog
\code{include} bez wpisywania ręcznych flag \code{-I}.

\begin{lstlisting}[language={},caption={CMakeLists.txt dla klasycznej biblioteki.}]
cmake_minimum_required(VERSION 3.20)
project(MathStuff LANGUAGES CXX)

add_library(mathstuff STATIC src/math.cpp)
target_compile_features(mathstuff PUBLIC cxx_std_20)
target_include_directories(mathstuff PUBLIC
    $<BUILD_INTERFACE:${CMAKE_CURRENT_SOURCE_DIR}/include>
    $<INSTALL_INTERFACE:include>
)

add_executable(app src/main.cpp)
target_link_libraries(app PRIVATE mathstuff)
\end{lstlisting}

\begin{lstlisting}[language=bash,caption={Konfigurowanie i budowanie projektu z CMake.}]
cmake -S . -B build -DCMAKE_BUILD_TYPE=Debug
cmake --build build -j
./build/app
\end{lstlisting}

\code{PRIVATE}, \code{PUBLIC} i~\code{INTERFACE} określają
\emph{propagację wymagań użycia} (usage requirements) między
targetami:
\begin{center}
\begin{tabularx}{\textwidth}{@{}lX@{}}
\toprule
Słowo & Znaczenie\\
\midrule
\code{PRIVATE} & Wymaganie tylko na potrzeby budowania danego targetu.\\
\code{INTERFACE} & Wymaganie wyłącznie dla jego konsumentów.\\
\code{PUBLIC} & Wymaganie dla targetu i~jego konsumentów.\\
\bottomrule
\end{tabularx}
\end{center}

Gdy publiczny nagłówek biblioteki wymaga pewnego katalogu
nagłówków, ten katalog powinien zwykle propagować się do klientów.
Gdy jednak biblioteka korzysta z~zależności jedynie wewnątrz
\code{.cpp}, często wystarcza \code{PRIVATE}. Właściwy zakres
zależy też od tego, czy target jest biblioteką statyczną lub
współdzieloną i~jakie zależności są konieczne przy ostatecznym
linkowaniu.

\subsection{Biblioteka nagłówkowa i~\texttt{INTERFACE}}

\begin{lstlisting}[language={},caption={Header-only target i wymagania dla klientow.}]
add_library(my_headers INTERFACE)
target_include_directories(my_headers INTERFACE
    ${CMAKE_CURRENT_SOURCE_DIR}/include
)
target_compile_features(my_headers INTERFACE cxx_std_20)

add_executable(tool src/tool.cpp)
target_link_libraries(tool PRIVATE my_headers)
\end{lstlisting}

Taki target nie tworzy archiwum \code{.a}. Reprezentuje zestaw
wymagań i~publiczne nagłówki używane przez konsumentów.

\subsection{Diagnostyka, sanitizatory, LTO i~build types}

\begin{lstlisting}[language=bash,caption={Przyk?adowe flagi diagnostyczne GCC/Clang.}]
g++ -std=c++20 -O0 -g -Wall -Wextra -Wpedantic \
    -fsanitize=address,undefined main.cpp -o app_debug

g++ -std=c++20 -O2 -flto main.cpp math.cpp -o app_lto
\end{lstlisting}

\emph{Link-Time Optimization} (LTO) może umożliwiać optymalizacje
pomiędzy TU, np. dodatkowe wstawianie funkcji do miejsc wywołania.
Jest to jednak kosztowna funkcja toolchainu; jej efekty zależą od
programu i~opcji. Profilowanie i~powtarzalne pomiary są ważniejsze
niż sam fakt włączenia \code{-O3} lub \code{-flto}.

\section{Moduły C++20: nowy mechanizm organizacji kodu}
\label{sec:cpp17-modules}

\subsection{Motywacja: problem tekstowego dołączania nagłówków}

W~klasycznym modelu ten sam popularny nagłówek może być parsowany
oddzielnie w~wielu TU. Makra i~kolejność \code{\#include} wpływają
na wynik preprocessingu. Moduły C++20 umożliwiają jawne eksportowanie
deklaracji i~importowanie ich bez tekstowego wklejania całych
interfejsów do każdego klienta.

Zaletami modułów mogą być izolacja makr, bardziej jednoznaczne
zależności i~krótsze czasy kompilacji. Nie ma jednak gwarancji,
że \emph{każdy} program po migracji przyspieszy, a~wdrożenie
wymaga wsparcia kompilatora i~systemu budowania.

\subsection{Pierwszy moduł: \texttt{export module} i~\texttt{import}}

Rozważmy dwa pliki:
\begin{lstlisting}[caption={Interfejs modu?u geometry: plik geometry.cppm.}]
export module geometry;

export int square(int x) {
    return x * x;
}

int hidden_helper(int x) {
    return x + 1;
}
\end{lstlisting}

\begin{lstlisting}[caption={Klient modu?u: plik main.cpp.}]
#include <iostream>
import geometry;

int main() {
    std::cout << square(7) << '\n'; // 49
    // hidden_helper(7); // Nie jest eksportowana.
}
\end{lstlisting}

\code{export module geometry;} deklaruje \emph{główną jednostkę
interfejsu modułu} (\emph{primary module interface unit}).
\code{export int square(...)} udostępnia deklarację klientowi.
Samo \code{import geometry;} nie oznacza linkowania kodu funkcji
z~powietrza: jeśli moduł generuje kod obiektowy, trzeba go również
uwzględnić podczas linkowania.

\subsection{BMI/CMI: moduł nie jest biblioteką binarną}

Kompilator wytwarza dla interfejsu modułu własną reprezentację,
często nazywaną \emph{BMI} (\emph{Binary Module Interface}) lub
\emph{CMI} (\emph{Compiled Module Interface}). W~Clangu spotyka się
pliki \code{.pcm}, w~GCC \code{.gcm}, a~w~MSVC pliki \code{.ifc}.
Nie są to uniwersalne, przenośne między kompilatorami artefakty
podobne do biblioteki \code{.a}. Importer musi znać interfejs
przetworzony przez zgodną konfigurację kompilatora.

Zależności muszą być budowane w~poprawnym porządku:
\[
\text{źródło interfejsu modułu}
\longrightarrow \text{BMI/CMI}
\longrightarrow \text{kompilacja importera}
\longrightarrow \text{linkowanie obiektów}.
\]

\subsection{Rzeczywiste komendy: Clang i~GCC}

Dla prostego modułu \code{geometry.cppm} w~Clangu obsługującym
standardowe moduły można zastosować:
\begin{lstlisting}[language=bash,caption={Budowanie przyk?adowego modu?u w Clangu (tryb BMI).}]
clang++ -std=c++20 geometry.cppm --precompile -o geometry.pcm
clang++ -std=c++20 -c geometry.pcm -o geometry.o
clang++ -std=c++20 -fmodule-file=geometry=geometry.pcm \
    -c main.cpp -o main.o
clang++ geometry.o main.o -o app
./app
\end{lstlisting}

Dla GCC~14 (wsparcie modułów C++20 było w~tej wersji nadal rozwijane)
trzeba dodatkowo włączyć eksperymentalny mechanizm
\code{-fmodules-ts}. W~przykładach z~GCC~14 przydatne bywa jawne
wskazanie języka przy rozszerzeniu \code{.cppm}:
\begin{lstlisting}[language=bash,caption={Budowanie tego samego modu?u z GCC 14.}]
g++ -std=c++20 -fmodules-ts -x c++ -c geometry.cppm \
    -o geometry.o
# GCC tworzy potrzebne dane interfejsu w gcm.cache/.
g++ -std=c++20 -fmodules-ts -c main.cpp -o main.o
g++ geometry.o main.o -o app
./app
\end{lstlisting}

Dokładne opcje i~formaty modułów mogą zmieniać się między wersjami
narzędzi. W~dużym projekcie lepiej powierzyć zależności systemowi
budowania obsługującemu skanowanie importów, zamiast ręcznie
zarządzać plikami BMI.

\subsection{Jednostka interfejsu i~jednostka implementacji}

Moduł może rozdzielać interfejs od implementacji bez tradycyjnego
nagłówka. Jednostka implementacji deklaruje \code{module name;},
a~nie \code{export module name;}:

\begin{lstlisting}[caption={Dwa pliki jednego modu?u: interface.cppm i implementation.cpp.}]
// Plik: engine.cppm
export module engine;
export int run(int value);

// Plik: engine.cpp
module engine;
int run(int value) { return 2 * value; }
\end{lstlisting}

Importer widzi tylko deklaracje eksportowane przez interfejs.
Jednostkę implementacji również trzeba skompilować i~zlinkować.
Konkretny system budowania musi uwzględnić fakt, że jednostka
implementacji importuje interfejs swojego modułu niejawnie.

\subsection{Global module fragment i~nagłówki klasyczne}

Tradycyjny nagłówek może wymagać makr konfigurujących jego treść.
W~takim przypadku można skorzystać z~\emph{global module fragment}:

\begin{lstlisting}[caption={Global module fragment: tradycyjny include przed modu?em.}]
module;
#include <string>
#include <utility>

export module greeting;

export std::string greet(std::string name) {
    return "Hello " + std::move(name);
}
\end{lstlisting}

Dyrektywa \code{module;} rozpoczyna fragment globalny, a~następna
deklaracja \code{export module greeting;} rozpoczyna właściwą
jednostkę modułu. Importowanie nagłówków jednostkowych przez
\code{import <string>;} jest innym mechanizmem, tzw. \emph{header units},
których dostępność zależy od toolchainu. Nie utożsamiajmy go z~nazwanym
modułem standardowej biblioteki \code{import std;}.

\subsection{Partycje i~ponowny eksport importu}

Duży moduł można podzielić na partycje. Nazwa po dwukropku oznacza
partycję wewnątrz jednego modułu, a~nie odrębny nazwany moduł.

\begin{lstlisting}[caption={Dwie jednostki interfejsu modu?u z partycja.}]
// Plik: math_detail.cppm
export module mathlib:detail;
export int triple(int x) { return 3 * x; }

// Plik: mathlib.cppm
export module mathlib;
export import :detail;
export int times_six(int x) { return 2 * triple(x); }

// Plik: main.cpp
import mathlib;
int main() { return times_six(7) == 42 ? 0 : 1; }
\end{lstlisting}

Klient importuje \code{mathlib}, a~nie \code{mathlib:detail}.
W~tym przykładzie \code{export import :detail;} re-eksportuje
eksporty partycji. Wymaga to prawidłowej kolejności budowania
\emph{najpierw partycji, potem głównego interfejsu, a~na końcu klienta}.

\subsection{Łączność modułowa i~ODR}

C++20 wprowadza reguły własności deklaracji (\emph{module ownership}).
Deklaracje przypisane do nazwanego modułu należą do niego, a~pewne
nieeksportowane nazwy uzyskują \emph{module linkage}. Różni się to od
mechanicznego zastąpienia \code{\#include} przez \code{import}.

Przykładowo dodanie \code{export} nie sprawia automatycznie, że
interfejs C++ będzie zgodny ABI między kompilatorami. Ponadto
źródłowy kod szablonów nadal musi być semantycznie dostępny
odpowiedniemu mechanizmowi instancjacji, choć nie musi być już
tekstowo wklejany przez \code{\#include}.

\subsection{CMake i~\texttt{FILE\_SET CXX\_MODULES}}

Od CMake~3.28 dostępne są mechanizmy skanowania modułów i~specjalny
typ zbioru plików \code{CXX_MODULES}. Przykład projektu:

\begin{lstlisting}[language={},caption={CMakeLists.txt projektu korzystajacego z modu?ow C++20.}]
cmake_minimum_required(VERSION 3.28)
project(GeometryModules LANGUAGES CXX)

add_library(geometry)
target_compile_features(geometry PUBLIC cxx_std_20)
target_sources(geometry PUBLIC
    FILE_SET CXX_MODULES FILES geometry.cppm
)

add_executable(app main.cpp)
target_link_libraries(app PRIVATE geometry)
\end{lstlisting}

\begin{lstlisting}[language=bash,caption={Konfiguracja modu?ow przy uzyciu generatora Ninja.}]
cmake -S . -B build -G Ninja
cmake --build build
./build/app
\end{lstlisting}

Ten przykład pokazuje \emph{składnię} CMake, ale jego zbudowanie
wymaga zgodnego zestawu wersji CMake, kompilatora, generatora i~skanera
zależności. W~szczególności kompilator akceptujący składnię modułów
niekoniecznie współpracuje z~automatycznym skanowaniem w~danej
konfiguracji CMake. Na współczesnych systemach używa się m.in.
Ninja (co najmniej 1.11 w~typowych konfiguracjach) i~odpowiednich
wersji Clanga lub MSVC. Obsługa modułów standardowej biblioteki oraz
\emph{header units} ma odrębne ograniczenia.

\subsection{\texttt{import std;} nie jest tym samym co \texttt{\#include}}

Sam standard C++20 wprowadza językowy mechanizm modułów, ale
\emph{nie zobowiązuje} każdej biblioteki standardowej do dostarczenia
nazwanego modułu \code{std} w~tym standardzie. \code{import std;}
jest standaryzowane w~C++23; obsługa zależy od kompilatora, dostawcy
biblioteki standardowej i~wersji systemu budowania.

\begin{lstlisting}[caption={Nazwany modu? standardowej biblioteki (C++23, zalezny od wsparcia).}]
import std;

int main() {
    std::vector<int> xs{1, 2, 3};
    std::println("size = {}", xs.size());
}
\end{lstlisting}

Powyższy przykład wykorzystuje także \code{std::println} z~C++23.
Nie należy automatycznie kompilować go z~\code{-std=c++20} ani
zakładać, że \code{import std;} działa w~każdym środowisku C++23.

\section{Optymalizacja procesu budowania}
\label{sec:cpp17-build-perf}

\subsection{Incremental build i~zależności nagłówkowe}

Zmiana publicznego nagłówka może wymusić ponowną kompilację wielu TU,
podczas gdy zmiana prywatnego \code{.cpp} zwykle dotyka mniejszej
części projektu. Dlatego ograniczanie niepotrzebnych dołączeń
nagłówków, stosowanie deklaracji wyprzedzających tam, gdzie są
prawidłowe, i~PImpl mogą poprawiać czas budowania.

Moduły zmieniają strukturę zależności, ale nie usuwają konieczności
ich śledzenia; importer musi otrzymać aktualną reprezentację
interfejsu modułu.

\subsection{Precompiled headers, unity builds, ccache}

Popularne techniki przyspieszania kompilacji to:
\begin{description}[style=nextline]
\item[Precompiled headers (PCH)]
Przechowywanie wyniku parsowania stabilnych, obszernych nagłówków.
Pomagają w~projektach z~dużą liczbą TU, ale wymagają spójności
opcji kompilacji i~zarządzania unieważnianiem pamięci podręcznej.
\item[Unity builds]
Łączenie kilku plików źródłowych we wspólną TU. Może ograniczać
narzut kompilacji, ale też ujawniać kolizje anonimowych przestrzeni
nazw, makr i~definicji o~łączności wewnętrznej.
\item[Compilation cache]
Narzędzia takie jak \code{ccache} unikają ponownego wykonywania
kompilacji, jeżeli odpowiednio zdefiniowane wejścia pozostają takie same.
\item[Parallel build]
Kompilowanie niezależnych TU równolegle (np. \code{cmake --build build -j})
wykorzystuje wiele rdzeni, ale zwiększa chwilowe użycie pamięci.
\end{description}

\subsection{Powtarzalność buildów i~zależności transytywne}

Dla bibliotek publicznych warto utrzymywać przejrzystą strukturę
zależności. Przykładowo użytkownik nie powinien polegać na tym,
że \code{\#include <vector>} ,,przypadkiem'' pojawiło się wskutek
dołączenia innego nagłówka. Każdy nagłówek powinien zawierać
potrzebne deklaracje w~sposób niezależny od kolejności dołączeń.

Warto również zapisywać wersje kompilatora i~bibliotek, flagi
wpływające na ABI, wariant budowania oraz sposób instalowania
artefaktów. \emph{Reproducible builds} wymagają dodatkowo kontroli
źródeł niedeterministyczności, np. znaczników czasu i~ścieżek
zapisywanych w~plikach wynikowych.

\section{Praktyczny przewodnik diagnostyczny}
\label{sec:cpp17-troubleshooting}

\begin{center}
\small
\begin{tabularx}{\textwidth}{@{}p{0.34\textwidth}X@{}}
\toprule
Objaw & Najbardziej prawdopodobne źródło problemu\\
\midrule
\code{undefined reference} & Brakujący plik obiektowy, biblioteka,
definicja albo instancjacja szablonu.\\
\code{multiple definition} & Niepoprawna definicja zwykłej
funkcji lub zmiennej w~nagłówku albo podwójnie linkowany obiekt.\\
\code{undefined reference to vtable} & Niezdefiniowana wymagana
funkcja wirtualna (często destruktor) lub brak pliku implementacji.\\
\code{cannot find -lfoo} & Brak biblioteki lub nieprawidłowe
katalogi przeszukiwania linkera.\\
\code{cannot open shared object file} & Problem ścieżki lub
zależności ładowania, niekoniecznie problem kompilacji.\\
Kompilacja klienta działa, program kończy się błędem & Możliwa
niezgodność ABI, naruszenie ODR, zły czas życia obiektu lub błąd programu.\\
Importer nie znajduje modułu & Nie zbudowano BMI/CMI, nie wskazano
ścieżki albo konfiguracja skanera nie obsługuje importu.\\
Import działa, ale linker zgłasza brak funkcji & Brak pliku
obiektowego modułu/implementacji lub potrzebnej biblioteki.\\
Zmiana nagłówka przebudowuje prawie wszystko & Silne zależności
nagłówkowe; rozważ redukcję dołączeń, PImpl, moduły lub PCH.\\
\bottomrule
\end{tabularx}
\end{center}

\subsection{Kiedy wybrać nagłówki, bibliotekę czy moduł?}

\begin{itemize}
\item \textbf{Header-only} sprawdza się dla małych narzędzi, szablonów
      i~funkcji \code{inline}; ceną mogą być zależności kompilacyjne.
\item \textbf{Biblioteka statyczna} upraszcza dystrybucję kodu obiektowego
      wewnątrz aplikacji, ale nie usuwa wymagań ABI w~czasie linkowania.
\item \textbf{Biblioteka współdzielona} pozwala współdzielić kod
      i~wymieniać implementację, lecz wymaga jasnej polityki ABI,
      widoczności symboli i~wersjonowania.
\item \textbf{Moduł C++20} służy do semantycznego podziału interfejsu
      źródłowego i~kontroli importów; nie jest alternatywą wobec
      procesu linkowania ani automatycznym formatem dystrybucji ABI.
\item \textbf{PImpl i~C API} rozważaj na granicach bibliotek o~długim
      cyklu życia, szczególnie gdy stabilność ABI ma duże znaczenie.
\end{itemize}

\section*{Podsumowanie rozdziału}
\addcontentsline{toc}{section}{Podsumowanie rozdziału}

Kompletny program C++ jest wynikiem wspólnej pracy preprocessora,
kompilatora, linkera i~loadera. Organizując projekt, trzeba osobno
rozumieć \emph{deklaracje i~definicje}, \emph{ODR i~łączność},
\emph{symbole i~biblioteki} oraz \emph{ABI}. Nagłówki i~moduły
pomagają udostępniać interfejsy źródłowe; biblioteki statyczne
i~współdzielone dostarczają kod do połączenia lub załadowania.

Dobra organizacja projektu polega na jasnym rozdziale API i~implementacji,
ograniczaniu zależności, korzystaniu z~targetów CMake oraz świadomym
wersjonowaniu interfejsu binarnego. Nowoczesne moduły upraszczają część
problemów nagłówków, ale wprowadzają wymagania dotyczące narzędzi,
BMI/CMI i~kolejności kompilacji.

\section*{Literatura i~dokumentacja}
\addcontentsline{toc}{section}{Literatura i dokumentacja}
\begin{itemize}
\item ISO/IEC 14882:2020, zagadnienia \emph{translation units},
      \emph{linkage}, \emph{One Definition Rule} i~\emph{modules}.
\item \url{https://en.cppreference.com/w/cpp/language/definition}
      --- ODR, deklaracje i~definicje.
\item \url{https://en.cppreference.com/w/cpp/language/modules}
      --- składnia i~semantyka modułów.
\item \url{https://gcc.gnu.org/onlinedocs/gcc/C_002b_002b-Modules.html}
      --- implementacja modułów w~GCC.
\item \url{https://clang.llvm.org/docs/StandardCPlusPlusModules.html}
      --- moduły w~Clangu, BMI i~komendy budowania.
\item \url{https://cmake.org/cmake/help/latest/manual/cmake-cxxmodules.7.html}
      --- CMake i~zależności modułów.
\item \url{https://cmake.org/cmake/help/latest/command/target_sources.html}
      --- \code{FILE_SET CXX_MODULES}.
\item \url{https://gcc.gnu.org/onlinedocs/gcc/Code-Gen-Options.html}
      --- widoczność symboli w~ELF.
\item John Lakos, \emph{Large-Scale C++ Software Design}, Addison-Wesley, 1996.
\end{itemize}
